% Options for packages loaded elsewhere
\PassOptionsToPackage{unicode}{hyperref}
\PassOptionsToPackage{hyphens}{url}
\documentclass[
]{article}
\usepackage{xcolor}
\usepackage{amsmath,amssymb}
\setcounter{secnumdepth}{-\maxdimen} % remove section numbering
\usepackage{iftex}
\ifPDFTeX
  \usepackage[T1]{fontenc}
  \usepackage[utf8]{inputenc}
  \usepackage{textcomp} % provide euro and other symbols
\else % if luatex or xetex
  \usepackage{unicode-math} % this also loads fontspec
  \defaultfontfeatures{Scale=MatchLowercase}
  \defaultfontfeatures[\rmfamily]{Ligatures=TeX,Scale=1}
\fi
\usepackage{lmodern}
\ifPDFTeX\else
  % xetex/luatex font selection
\fi
% Use upquote if available, for straight quotes in verbatim environments
\IfFileExists{upquote.sty}{\usepackage{upquote}}{}
\IfFileExists{microtype.sty}{% use microtype if available
  \usepackage[]{microtype}
  \UseMicrotypeSet[protrusion]{basicmath} % disable protrusion for tt fonts
}{}
\makeatletter
\@ifundefined{KOMAClassName}{% if non-KOMA class
  \IfFileExists{parskip.sty}{%
    \usepackage{parskip}
  }{% else
    \setlength{\parindent}{0pt}
    \setlength{\parskip}{6pt plus 2pt minus 1pt}}
}{% if KOMA class
  \KOMAoptions{parskip=half}}
\makeatother
\usepackage{longtable,booktabs,array}
\newcounter{none} % for unnumbered tables
\usepackage{multirow}
\usepackage{calc} % for calculating minipage widths
% Correct order of tables after \paragraph or \subparagraph
\usepackage{etoolbox}
\makeatletter
\patchcmd\longtable{\par}{\if@noskipsec\mbox{}\fi\par}{}{}
\makeatother
% Allow footnotes in longtable head/foot
\IfFileExists{footnotehyper.sty}{\usepackage{footnotehyper}}{\usepackage{footnote}}
\makesavenoteenv{longtable}
\usepackage{graphicx}
\makeatletter
\newsavebox\pandoc@box
\newcommand*\pandocbounded[1]{% scales image to fit in text height/width
  \sbox\pandoc@box{#1}%
  \Gscale@div\@tempa{\textheight}{\dimexpr\ht\pandoc@box+\dp\pandoc@box\relax}%
  \Gscale@div\@tempb{\linewidth}{\wd\pandoc@box}%
  \ifdim\@tempb\p@<\@tempa\p@\let\@tempa\@tempb\fi% select the smaller of both
  \ifdim\@tempa\p@<\p@\scalebox{\@tempa}{\usebox\pandoc@box}%
  \else\usebox{\pandoc@box}%
  \fi%
}
% Set default figure placement to htbp
\def\fps@figure{htbp}
\makeatother
\ifLuaTeX
  \usepackage{luacolor}
  \usepackage[soul]{lua-ul}
\else
  \usepackage{soul}
\fi
\setlength{\emergencystretch}{3em} % prevent overfull lines
\providecommand{\tightlist}{%
  \setlength{\itemsep}{0pt}\setlength{\parskip}{0pt}}
\usepackage{bookmark}
\IfFileExists{xurl.sty}{\usepackage{xurl}}{} % add URL line breaks if available
\urlstyle{same}
\hypersetup{
  hidelinks,
  pdfcreator={LaTeX via pandoc}}

\author{}
\date{}

\begin{document}

\begin{quote}
BỘ GIÁO DỤC VÀ ĐÀO TẠO

\textbf{TRƯỜNG ĐẠI HỌC CÔNG NGHỆ TP. HCM}
\end{quote}


\includegraphics[width=2.31519in,height=2.78208in]{./media/image1.jpeg}

\begin{quote}
\textbf{HOÀNG TRẦN GIA PHÚC}

\textbf{NGHIÊN CỨU VÀ PHÁT TRIỂN MÔ HÌNH HỌC SÂU CHO BÀI TOÁN PHÂN LOẠI
CHẤT LƯỢNG TỎI}

\textbf{LUẬN VĂN THẠC SĨ}

Chuyên ngành: Công nghệ Thông tin Mã số ngành:
\end{quote}

TP. HỒ CHÍ MINH, tháng \hl{09} năm 2026

\begin{quote}
BỘ GIÁO DỤC VÀ ĐÀO TẠO

\textbf{TRƯỜNG ĐẠI HỌC CÔNG NGHỆ TP. HCM}
\end{quote}


\includegraphics[width=2.31519in,height=2.78208in]{./media/image1.jpeg}
\includegraphics[width=1.95704in,height=2.34708in]{./media/image2.jpeg}

\begin{quote}
\textbf{HOÀNG TRẦN GIA PHÚC}

\textbf{NGHIÊN CỨU VÀ PHÁT TRIỂN MÔ HÌNH HỌC SÂU CHO BÀI TOÁN PHÂN LOẠI
CHẤT LƯỢNG TỎI}

\textbf{LUẬN VĂN THẠC SĨ}

Chuyên ngành: Công nghệ Thông tin Mã số ngành:
\end{quote}

\textbf{CÁN BỘ HƯỚNG DẪN KHOA HỌC: TS. LÊ CUNG TƯỞNG}

\begin{quote}
TP. HỒ CHÍ MINH, tháng \hl{09} năm 2026

CÔNG TRÌNH ĐƯỢC HOÀN THÀNH TẠI

\textbf{TRƯỜNG ĐẠI HỌC CÔNG NGHỆ TP. HCM}

Cán bộ hướng dẫn khoa học: TS. Lê Cung Tưởng

\emph{(Ghi rõ họ, tên, học hàm, học vị và chữ ký)}

Đề án Thạc sĩ được bảo vệ tại Trường Đại học Công nghệ TP. HCM,
ngày\ldots{} tháng \ldots{} năm\ldots.

Thành phần Hội đồng đánh giá Đề án Thạc sĩ gồm:

\emph{(Ghi rõ họ, tên, học hàm, học vị của Hội đồng chấm bảo vệ Đề án
Thạc sĩ)}
\end{quote}

{\def\LTcaptype{none} % do not increment counter
\begin{longtable}[]{@{}
  >{\centering\arraybackslash}p{(\linewidth - 4\tabcolsep) * \real{0.0657}}
  >{\raggedright\arraybackslash}p{(\linewidth - 4\tabcolsep) * \real{0.3638}}
  >{\centering\arraybackslash}p{(\linewidth - 4\tabcolsep) * \real{0.3421}}@{}}
\toprule\noalign{}
\begin{minipage}[b]{\linewidth}\centering
\textbf{TT}
\end{minipage} & \begin{minipage}[b]{\linewidth}\centering
\begin{quote}
\textbf{Họ và tên}
\end{quote}
\end{minipage} & \begin{minipage}[b]{\linewidth}\centering
\begin{quote}
\textbf{Chức danh Hội đồng}
\end{quote}
\end{minipage} \\
\midrule\noalign{}
\endhead
\bottomrule\noalign{}
\endlastfoot
1 & & \begin{minipage}[t]{\linewidth}\centering
\begin{quote}
Chủ tịch
\end{quote}
\end{minipage} \\
2 & & \begin{minipage}[t]{\linewidth}\centering
\begin{quote}
Phản biện 1
\end{quote}
\end{minipage} \\
3 & & \begin{minipage}[t]{\linewidth}\centering
\begin{quote}
Phản biện 2
\end{quote}
\end{minipage} \\
4 & & \begin{minipage}[t]{\linewidth}\centering
\begin{quote}
Ủy viên
\end{quote}
\end{minipage} \\
5 & & \begin{minipage}[t]{\linewidth}\centering
\begin{quote}
Ủy viên, Thư ký
\end{quote}
\end{minipage} \\
\end{longtable}
}

\begin{quote}
Xác nhận của Chủ tịch Hội đồng đánh giá Luận sau khi Đề án đã được sửa
chữa (nếu có).

\textbf{Chủ tịch Hội đồng đánh giá Đề án}
\end{quote}

{\def\LTcaptype{none} % do not increment counter
\begin{longtable}[]{@{}
  >{\raggedright\arraybackslash}p{(\linewidth - 2\tabcolsep) * \real{0.4373}}
  >{\raggedright\arraybackslash}p{(\linewidth - 2\tabcolsep) * \real{0.5415}}@{}}
\toprule\noalign{}
\begin{minipage}[b]{\linewidth}\centering
TRƯỜNG ĐH CÔNG NGHỆ TP. HCM

\textbf{VIỆN ĐÀO TẠO SAU ĐẠI HỌC}
\end{minipage} & \begin{minipage}[b]{\linewidth}\centering
\textbf{CỘNG HÒA XÃ HỘI CHỦ NGHĨA VIỆT NAM}

\textbf{Độc lập -- Tự do -- Hạnh phúc}
\end{minipage} \\
\midrule\noalign{}
\endhead
\bottomrule\noalign{}
\endlastfoot
\end{longtable}
}

\emph{TP. HCM, ngày}..\emph{\ldots{} tháng\ldots{}}.. \emph{năm
20}..\emph{\ldots{}}

\textbf{\hl{NHIỆM VỤ LUẬN VĂN THẠC SĨ}}

Họ tên học viên:
..........................................................................Giới
tính:

Ngày, tháng, năm sinh:
................................................................Nơi
sinh:

Chuyên ngành:
.............................................................................MSHV:

\textbf{I- Tên đề tài:}

\textbf{II- Nhiệm vụ và nội dung:}

\textbf{III- Ngày giao nhiệm vụ:} \emph{(Ngày bắt đầu thực hiện LV ghi
trong QĐ giao đề tài)}

\textbf{IV- Ngày hoàn thành nhiệm vụ:}

\textbf{V- Cán bộ hướng dẫn:} \emph{(Ghi rõ học hàm, học vị, họ, tên)}

\textbf{CÁN BỘ HƯỚNG DẪN KHOA QUẢN LÝ CHUYÊN NGÀNH}

(Họ tên và chữ ký) (Họ tên và chữ ký)

\section{\texorpdfstring{\hl{LỜI CAM
ĐOAN}}{LỜI CAM ĐOAN}}\label{lux1eddi-cam-ux111oan}

Tôi xin cam kết rằng đây là công trình nghiên cứu do chính tôi thực
hiện. Toàn bộ số liệu, kết quả và nội dung trình bày trong Đề án là
trung thực, không sao chép từ bất cứ nguồn nào và chưa từng được công bố
trong bất kỳ công trình nghiên cứu nào khác.

Tôi xin cam đoan rằng tất cả các tài liệu tham khảo được sử dụng trong
Đề án đều đã được trích dẫn đầy đủ và ghi rõ nguồn gốc theo đúng quy
định.

\begin{quote}
\textbf{Học viên thực hiện Đề án}

\textbf{Hoàng Trần Gia Phúc}
\end{quote}

\section{\texorpdfstring{\hl{LỜI CÁM
ƠN}}{LỜI CÁM ƠN}}\label{lux1eddi-cuxe1m-ux1a1n}

Trước hết, tôi xin bày tỏ lòng biết ơn sâu sắc đến Tiến sĩ Lê Cung
Tưởng, người đã tận tình hướng dẫn, chỉ bảo và hỗ trợ tôi trong suốt quá
trình thực hiện Đề án này với sự tận tâm và chuyên môn cao.

Tôi xin chân thành cảm ơn Quý Thầy, Cô tại Khoa Công nghệ Thông tin,
Trường Đại học Công nghệ TP. HCM (HUTECH). Những người đã tận tâm giảng
dạy, truyền đạt kiến thức và tạo nền tảng vững chắc giúp tôi có thể hoàn
thành quá trình học tập cũng như nghiên cứu.

Tôi cũng xin gửi lời tri ân gia đình, bạn bè và đồng nghiệp đã luôn đồng
hành, động viên và ủng hộ tôi trong suốt hành trình học tập và thực hiện
Đề án này.

Mặc dù đã nỗ lực hết mình, song tôi nhận thức rằng Đề án này vẫn có thể
còn những thiếu sót. Tôi rất mong nhận được những ý kiến đóng góp quý
giá từ Quý Thầy, Cô để công trình được hoàn thiện hơn.

Một lần nữa, xin trân trọng cảm ơn tất cả những ai đã hỗ trợ và đồng
hành cùng tôi trong quá trình học tập và nghiên cứu.

\begin{quote}
\textbf{Học viên thực hiện Đề án}

\textbf{Hoàng Trần Gia Phúc}
\end{quote}

\section{\texorpdfstring{\hl{TÓM TẮT}}{TÓM TẮT}}\label{tuxf3m-tux1eaft}

\section{\texorpdfstring{\hl{ABSTRACT}}{ABSTRACT}}\label{abstract}

\section{\texorpdfstring{\hl{MỤC
LỤC}}{MỤC LỤC}}\label{mux1ee5c-lux1ee5c}

\section{}\label{section}

\hyperref[lux1eddi-cam-ux111oan]{\hl{LỜI CAM ĐOAN}
\hyperref[lux1eddi-cam-ux111oan]{1}}

\hyperref[lux1eddi-cuxe1m-ux1a1n]{\hl{LỜI CÁM ƠN}
\hyperref[lux1eddi-cuxe1m-ux1a1n]{2}}

\hyperref[tuxf3m-tux1eaft]{\hl{TÓM TẮT} \hyperref[tuxf3m-tux1eaft]{3}}

\hyperref[abstract]{\hl{ABSTRACT} \hyperref[abstract]{3}}

\hyperref[mux1ee5c-lux1ee5c]{\hl{MỤC LỤC}
\hyperref[mux1ee5c-lux1ee5c]{3}}

\hyperref[danh-mux1ee5c-cuxe1c-tux1eeb-viux1ebft-tux1eaft]{\hl{DANH MỤC
CÁC TỪ VIẾT TẮT}
\hyperref[danh-mux1ee5c-cuxe1c-tux1eeb-viux1ebft-tux1eaft]{7}}

\hyperref[danh-mux1ee5c-cuxe1c-bux1ea3ng]{\hl{DANH MỤC CÁC BẢNG}
\hyperref[danh-mux1ee5c-cuxe1c-bux1ea3ng]{8}}

\hyperref[danh-mux1ee5c-cuxe1c-huxecnh]{\hl{DANH MỤC CÁC HÌNH}
\hyperref[danh-mux1ee5c-cuxe1c-huxecnh]{9}}

\hyperref[chux1b0ux1a1ng-1-giux1edbi-thiux1ec7u-tux1ed5ng-quan]{CHƯƠNG
1: GIỚI THIỆU TỔNG QUAN
\hyperref[chux1b0ux1a1ng-1-giux1edbi-thiux1ec7u-tux1ed5ng-quan]{10}}

\hyperref[_bookmark9]{\textbf{1.1} \textbf{Tính cấp thiết của đề tài.}
\hyperref[_bookmark9]{11}}

\hyperref[_Toc232111981]{\textbf{1.2} \textbf{Mục tiêu nghiên cứu}
\hyperref[_Toc232111981]{16}}

\hyperref[_Toc232111982]{\emph{\textbf{1.2.1}} \emph{\textbf{Mục tiêu
tổng quát}} \hyperref[_Toc232111982]{16}}

\hyperref[_Toc232111983]{\emph{\textbf{1.2.2}} \emph{\textbf{Mục tiêu cụ
thể}} \hyperref[_Toc232111983]{16}}

\hyperref[_Toc232111984]{\textbf{1.3} \textbf{Đối tượng, phạm vi nghiên
cứu} \hyperref[_Toc232111984]{19}}

\hyperref[_Toc232111985]{\emph{\textbf{1.3.1}} \emph{\textbf{Đối tượng
nghiên cứu}} \hyperref[_Toc232111985]{19}}

\hyperref[_Toc232111986]{\emph{\textbf{1.3.2}} \emph{\textbf{Phạm vi
nghiên cứu}} \hyperref[_Toc232111986]{21}}

\hyperref[_Toc232111987]{\textbf{1.4} \textbf{Phương pháp nghiên cứu}
\hyperref[_Toc232111987]{23}}

\hyperref[_Toc232111988]{\textbf{1.5} \textbf{Ý nghĩa nghiên cứu}
\hyperref[_Toc232111988]{27}}

\hyperref[_Toc232111989]{\emph{\textbf{1.5.1}} \emph{\textbf{Ý nghĩa
khoa học}} \hyperref[_Toc232111989]{27}}

\hyperref[_Toc232111990]{\emph{\textbf{1.5.2}} \emph{\textbf{Ý nghĩa
thực tiễn}} \hyperref[_Toc232111990]{27}}

\hyperref[_Toc232111991]{\emph{\textbf{1.5.3}} \emph{\textbf{Ý nghĩa
kinh tế và xã hội}} \hyperref[_Toc232111991]{28}}

\hyperref[_Toc232111992]{\emph{\textbf{1.5.4}} \emph{\textbf{Ý nghĩa mở
rộng và định hướng phát triển}} \hyperref[_Toc232111992]{28}}

\hyperref[_Toc232111993]{\textbf{\hl{1.6}} \textbf{\hl{Cấu trúc luận
văn}} \hyperref[_Toc232111993]{29}}

\hyperref[_Toc232111994]{\textbf{\hl{1.7}} \textbf{\hl{Tiến độ thực hiện
luận văn}} \hyperref[_Toc232111994]{30}}

\hyperref[tuxf3m-tux1eaft-chux1b0ux1a1ng-1]{TÓM TẮT CHƯƠNG 1
\hyperref[tuxf3m-tux1eaft-chux1b0ux1a1ng-1]{31}}

\hyperref[chux1b0ux1a1ng-2-cux1a1-sux1edf-luxfd-thuyux1ebft-vuxe0-muxf4-huxecnh-nghiuxean-cux1ee9u]{CHƯƠNG
2: CƠ SỞ LÝ THUYẾT VÀ MÔ HÌNH NGHIÊN CỨU
\hyperref[chux1b0ux1a1ng-2-cux1a1-sux1edf-luxfd-thuyux1ebft-vuxe0-muxf4-huxecnh-nghiuxean-cux1ee9u]{32}}

\hyperref[_Toc232111997]{\textbf{2.1} \textbf{Tổng quan về bài toán phân
loại chất lượng tỏi} \hyperref[_Toc232111997]{32}}

\hyperref[_Toc232111998]{\textbf{2.2} \textbf{Tổng quan về bộ dữ liệu
ảnh tỏi sử dụng trong nghiên cứu} \hyperref[_Toc232111998]{33}}

\hyperref[_Toc232111999]{\textbf{2.3} \textbf{Cơ sở lý thuyết về học
sâu} \hyperref[_Toc232111999]{35}}

\hyperref[_Toc232112000]{\emph{\textbf{2.3.1}} \emph{\textbf{Khái niệm
học sâu}} \hyperref[_Toc232112000]{36}}

\hyperref[_Toc232112001]{\emph{\textbf{2.3.2}} \emph{\textbf{Cấu trúc
tổng quát của mạng nơ-ron}} \hyperref[_Toc232112001]{36}}

\hyperref[_Toc232112002]{\emph{\textbf{2.3.3}} \emph{\textbf{Quá trình
học của mô hình}} \hyperref[_Toc232112002]{37}}

\hyperref[_Toc232112003]{\emph{\textbf{2.3.4}} \emph{\textbf{Vai trò của
học sâu trong phân loại ảnh}} \hyperref[_Toc232112003]{38}}

\hyperref[_Toc232112004]{\textbf{2.4} \textbf{Mạng nơ-ron tích chập
trong phân loại ảnh} \hyperref[_Toc232112004]{38}}

\hyperref[_Toc232112005]{\emph{\textbf{2.4.1}} \emph{\textbf{Khái niệm
mạng nơ-ron tích chập}} \hyperref[_Toc232112005]{39}}

\hyperref[_Toc232112006]{\emph{\textbf{2.4.2}} \emph{\textbf{Lớp tích
chập}} \hyperref[_Toc232112006]{40}}

\hyperref[_Toc232112007]{\emph{\textbf{2.4.3}} \emph{\textbf{Hàm kích
hoạt trong CNN}} \hyperref[_Toc232112007]{41}}

\hyperref[_Toc232112008]{\emph{\textbf{2.4.4}} \emph{\textbf{Lớp
pooling}} \hyperref[_Toc232112008]{41}}

\hyperref[_Toc232112009]{\emph{\textbf{2.4.5}} \emph{\textbf{Lớp phân
loại}} \hyperref[_Toc232112009]{42}}

\hyperref[_Toc232112010]{\emph{\textbf{2.4.6}} \emph{\textbf{Vai trò của
CNN trong phân loại ảnh}} \hyperref[_Toc232112010]{42}}

\hyperref[_Toc232112011]{\textbf{2.5} \textbf{Học chuyển giao trong phân
loại ảnh} \hyperref[_Toc232112011]{42}}

\hyperref[_Toc232112012]{\emph{\textbf{2.5.1}} \emph{\textbf{Khái niệm
học chuyển giao}} \hyperref[_Toc232112012]{43}}

\hyperref[_Toc232112013]{\emph{\textbf{2.5.2}} \emph{\textbf{Pre-trained
model trong phân loại ảnh}} \hyperref[_Toc232112013]{44}}

\hyperref[_Toc232112014]{\emph{\textbf{2.5.3}} \emph{\textbf{Các chiến
lược sử dụng học chuyển giao}} \hyperref[_Toc232112014]{44}}

\hyperref[_Toc232112015]{\emph{\textbf{2.5.4}} \emph{\textbf{Fine-tuning
trong phân loại ảnh}} \hyperref[_Toc232112015]{45}}

\hyperref[_Toc232112016]{\emph{\textbf{2.5.5}} \emph{\textbf{Vai trò của
học chuyển giao trong nghiên cứu}} \hyperref[_Toc232112016]{45}}

\hyperref[_Toc232112017]{\textbf{2.6} \textbf{Các mô hình CNN được sử
dụng trong nghiên cứu} \hyperref[_Toc232112017]{46}}

\hyperref[_Toc232112018]{\emph{\textbf{2.5.6}} \emph{\textbf{Mô hình
NASNetMobile}} \hyperref[_Toc232112018]{46}}

\hyperref[_Toc232112019]{\emph{\textbf{2.5.7}} \emph{\textbf{Mô hình
Inception}} \hyperref[_Toc232112019]{49}}

\hyperref[_Toc232112020]{\emph{\textbf{2.5.8}} \emph{\textbf{Mô hình
EfficientNet}} \hyperref[_Toc232112020]{53}}

\hyperref[_Toc232112021]{\emph{\textbf{2.5.9}} \emph{\textbf{Mô hình
MobileNet}} \hyperref[_Toc232112021]{56}}

\hyperref[_Toc232112022]{\emph{\textbf{2.5.10}} \emph{\textbf{Mô hình
DenseNet}} \hyperref[_Toc232112022]{58}}

\hyperref[tuxf3m-tux1eaft-chux1b0ux1a1ng-2]{TÓM TẮT CHƯƠNG 2
\hyperref[tuxf3m-tux1eaft-chux1b0ux1a1ng-2]{62}}

\hyperref[chux1b0ux1a1ng-3-muxf4-huxecnh-ux111ux1ec1-xuux1ea5t]{CHƯƠNG
3: MÔ HÌNH ĐỀ XUẤT
\hyperref[chux1b0ux1a1ng-3-muxf4-huxecnh-ux111ux1ec1-xuux1ea5t]{63}}

\hyperref[_Toc232112025]{\textbf{3.1 Các mô hình học sâu tiền huấn luyện
cho phân loại} \hyperref[_Toc232112025]{63}}

\hyperref[_Toc232112026]{\emph{\textbf{3.1.1 Mục tiêu}}
\hyperref[_Toc232112026]{63}}

\hyperref[_Toc232112027]{\emph{\textbf{3.1.2 Các mô hình được lựa chọn}}
\hyperref[_Toc232112027]{63}}

\hyperref[_bookmark21]{\emph{\textbf{3.1.3 Bộ dữ liệu và tiền xử lý}}
\hyperref[_bookmark21]{63}}

\hyperref[_Toc232112029]{\emph{\textbf{3.1.4 Xây dựng và huấn luyện mô
hình}} \hyperref[_Toc232112029]{63}}

\hyperref[_Toc232112030]{\textbf{3.2 Mô hình đề xuất CPGC-EfficientNetB4}
\hyperref[_Toc232112030]{65}}

\hyperref[_Toc232112031]{\emph{\textbf{3.2.1 Hạn chế của mô hình
EfficientNetB4}} \hyperref[_Toc232112031]{65}}

\hyperref[_Toc232112032]{\emph{\textbf{3.2.2 Ý tưởng cải tiến}}
\hyperref[_Toc232112032]{65}}

\hyperref[_Toc232112033]{\emph{\textbf{3.2.3 Kiến trúc}}
\hyperref[_Toc232112033]{65}}

\hyperref[chux1b0ux1a1ng-4-kux1ebft-quux1ea3-thux1ef1c-nghiux1ec7m]{CHƯƠNG
4: KẾT QUẢ THỰC NGHIỆM
\hyperref[chux1b0ux1a1ng-4-kux1ebft-quux1ea3-thux1ef1c-nghiux1ec7m]{66}}

\hyperref[_Toc232112035]{\textbf{4.1 Tiêu chí đánh giá và công cụ hỗ
trợ} \hyperref[_Toc232112035]{66}}

\hyperref[_Toc232112036]{\emph{\textbf{4.1.1 Tiêu chí đánh giá}}
\hyperref[_Toc232112036]{66}}

\hyperref[_Toc232112037]{\emph{\textbf{4.1.2 Công cụ hỗ trợ}}
\hyperref[_Toc232112037]{67}}

\hyperref[_Toc232112038]{\textbf{4.2 Kết quả của 5 mô hình Pre-trained}
\hyperref[_Toc232112038]{67}}

\hyperref[_Toc232112039]{\textbf{4.3 Kết quả của mô hình đề xuất}
\hyperref[_Toc232112039]{68}}

\hyperref[chux1b0ux1a1ng-5-kux1ebft-luux1eadn]{CHƯƠNG 5: KẾT LUẬN
\hyperref[chux1b0ux1a1ng-5-kux1ebft-luux1eadn]{69}}

\hyperref[_Toc232112041]{\textbf{5.1 Tổng kết quá trình nghiên cứu}
\hyperref[_Toc232112041]{69}}

\hyperref[_Toc232112042]{\textbf{5.2 Hạn chế và hướng phát triển}
\hyperref[_Toc232112042]{69}}

\hyperref[_Toc232112043]{\emph{\textbf{5.2.1 Hạn chế:}}
\hyperref[_Toc232112043]{69}}

\hyperref[_Toc232112044]{\emph{\textbf{5.2.2 Hướng phát triển trong
tương lai:}} \hyperref[_Toc232112044]{69}}

\hyperref[_Toc232112045]{\textbf{5.3 Khả năng ứng dụng thực tiễn}
\hyperref[_Toc232112045]{69}}

\hyperref[tuxe0i-liux1ec7u-tham-khux1ea3o]{TÀI LIỆU THAM KHẢO
\hyperref[tuxe0i-liux1ec7u-tham-khux1ea3o]{70}}

\section{\texorpdfstring{\hl{DANH MỤC CÁC TỪ VIẾT
TẮT}}{DANH MỤC CÁC TỪ VIẾT TẮT}}\label{danh-mux1ee5c-cuxe1c-tux1eeb-viux1ebft-tux1eaft}

{\def\LTcaptype{none} % do not increment counter
\begin{longtable}[]{@{}
  >{\centering\arraybackslash}p{(\linewidth - 6\tabcolsep) * \real{0.0691}}
  >{\centering\arraybackslash}p{(\linewidth - 6\tabcolsep) * \real{0.1241}}
  >{\raggedright\arraybackslash}p{(\linewidth - 6\tabcolsep) * \real{0.4122}}
  >{\raggedright\arraybackslash}p{(\linewidth - 6\tabcolsep) * \real{0.3482}}@{}}
\toprule\noalign{}
\begin{minipage}[b]{\linewidth}\centering
\textbf{STT}
\end{minipage} & \begin{minipage}[b]{\linewidth}\centering
\textbf{Viết tắt}
\end{minipage} & \begin{minipage}[b]{\linewidth}\centering
\textbf{Tiếng Anh}
\end{minipage} & \begin{minipage}[b]{\linewidth}\centering
\textbf{Tiếng Việt}
\end{minipage} \\
\midrule\noalign{}
\endhead
\bottomrule\noalign{}
\endlastfoot
1 & CNN & Convolutional Neural Network & Mạng nơ-ron tích chập \\
2 & SE & Squeeze-and-Excitation & Cơ chế nén và kích hoạt kênh \\
3 & ECA & Efficient Channel Attention & Cơ chế chú ý kênh hiệu quả \\
4 & GCD & Garlic Clove Dataset & Bộ dữ liệu ảnh tép tỏi \\
5 & IoT & Internet of Things & Internet vạn vật \\
6 & XAI & Explainable Artificial Intelligence & Trí tuệ nhân tạo có khả
năng giải thích \\
7 & BVTV & ~ & Bảo vệ Thực vật \\
8 & TP. HCM & ~ & Thành phố Hồ Chí Minh \\
\multirow{2}{=}{\centering\arraybackslash 9} &
\multirow{2}{=}{\centering\arraybackslash HUTECH} & Ho Chi Minh City
University of & \multirow{2}{=}{Trường Đại học Công nghệ TP.HCM} \\
& & Technology \\
\multirow{2}{=}{\centering\arraybackslash 10} &
\multirow{2}{=}{\centering\arraybackslash CVPR} &
\multirow{2}{=}{Conference on Computer Vision and Pattern Recognition} &
\multirow{2}{=}{Hội nghị về Thị giác Máy tính và Nhận dạng Mẫu} \\
 \\
\multirow{2}{=}{\centering\arraybackslash 11} &
\multirow{2}{=}{\centering\arraybackslash ACM} & Association for
Computing & \multirow{2}{=}{Hiệp hội Máy tính Hoa Kỳ} \\
& & Machinery \\
12 & CPGC & Classification of Peeled Garlic Cloves & Khung phân loại tép
tỏi đã bóc vỏ \\
13 & GIS & Geographic Information System & Hệ thống thông tin địa lý \\
\end{longtable}
}

\section{\texorpdfstring{\hl{DANH MỤC CÁC
BẢNG}}{DANH MỤC CÁC BẢNG}}\label{danh-mux1ee5c-cuxe1c-bux1ea3ng}

Bảng 2.1
\ldots\ldots\ldots\ldots\ldots\ldots\ldots\ldots\ldots\ldots\ldots\ldots\ldots\ldots\ldots\ldots\ldots\ldots\ldots\ldots\ldots\ldots\ldots\ldots\ldots\ldots Trang
33

\section{\texorpdfstring{\hl{DANH MỤC CÁC
HÌNH}}{DANH MỤC CÁC HÌNH}}\label{danh-mux1ee5c-cuxe1c-huxecnh}

Hình 1.1 Hình ảnh
\ldots\ldots\ldots\ldots\ldots\ldots\ldots\ldots\ldots\ldots\ldots\ldots\ldots\ldots\ldots\ldots\ldots\ldots\ldots\ldots\ldots\ldots\ldots\ldots Trang
28

`

\section{\texorpdfstring{CHƯƠNG 1: GIỚI THIỆU TỔNG QUAN
}{CHƯƠNG 1: GIỚI THIỆU TỔNG QUAN }}\label{chux1b0ux1a1ng-1-giux1edbi-thiux1ec7u-tux1ed5ng-quan}

\begin{enumerate}
\def\labelenumi{\arabic{enumi}.}
\item
  \protect\phantomsection\label{_bookmark9}{}\textbf{Tính cấp thiết của
  đề tài.}
\end{enumerate}

Nông nghiệp là một trong những lĩnh vực có vai trò quan trọng trong nền
kinh tế Việt Nam, góp phần bảo đảm an ninh lương thực, tạo việc làm và
cung cấp nguyên liệu cho ngành công nghiệp chế biến thực phẩm. Theo số
liệu năm 2025, khu vực nông, lâm nghiệp và thủy sản chiếm 11,64\% trong
cơ cấu nền kinh tế Việt Nam {[}26{]}.
Trong bối cảnh sản xuất nông nghiệp ngày càng hướng đến tự động hóa và
nâng cao giá trị sau thu hoạch, việc ứng dụng công nghệ vào các khâu sơ
chế, phân loại và kiểm soát chất lượng nông sản trở thành yêu cầu cấp
thiết.

Tỏi là một loại nông sản được sử dụng rộng rãi trong đời sống, chế biến
thực phẩm và các sản phẩm gia vị. Trong dây chuyền chế biến tỏi, nguyên
liệu thường trải qua các công đoạn như tách tép, bóc vỏ, làm sạch, phân
loại và đưa vào chế biến tiếp theo. Mặc dù các công đoạn tách tép và bóc
vỏ đã có thể được hỗ trợ bởi máy móc, nhưng sau quá trình bóc vỏ, tỏi
vẫn thường xuất hiện nhiều trạng thái khác nhau như tỏi đã bóc sạch, tỏi
còn sót vỏ, tỏi bị hư hỏng, dập nát hoặc biến màu. Do đó, công đoạn phân
loại sau bóc vỏ vẫn giữ vai trò quan trọng trong việc bảo đảm chất lượng
nguyên liệu đầu vào.

Hiện nay, việc phân loại tỏi sau bóc vỏ tại nhiều cơ sở sản xuất vẫn chủ
yếu được thực hiện thủ công. Người lao động phải quan sát từng tép tỏi
để phân biệt các nhóm như tỏi sạch, tỏi còn vỏ và tỏi hư. Phương pháp
này có nhiều hạn chế: tốn thời gian, cần nhiều nhân công, phụ thuộc vào
kinh nghiệm người kiểm tra và dễ xảy ra sai sót khi làm việc trong thời
gian dài. Đặc biệt, trong dây chuyền sản xuất thực phẩm, nếu tỏi hư hoặc
tỏi còn sót vỏ bị lẫn vào nhóm tỏi sạch, chất lượng sản phẩm cuối cùng
có thể bị ảnh hưởng, làm giảm độ đồng nhất và giá trị thương mại của sản
phẩm.

Sự phát triển của trí tuệ nhân tạo, thị giác máy tính và học sâu đã mở
ra khả năng tự động hóa công đoạn kiểm tra chất lượng nông sản. Các
nghiên cứu gần đây cho thấy học máy và học sâu đang được ứng dụng mạnh
mẽ trong kiểm soát chất lượng thực phẩm, phát hiện lỗi sản phẩm, kiểm
tra bề mặt và phân loại tự động trong dây chuyền sản xuất
{[}27{]}. Đối với nông sản,
các hệ thống thị giác nhân tạo có thể thay thế hoặc hỗ trợ con người
trong việc phát hiện lỗi, đánh giá chất lượng và phân loại sản phẩm dựa
trên ảnh đầu vào.

Đối với bài toán phân loại chất lượng tỏi, mô hình học sâu có thể học
các đặc trưng trực quan như màu sắc, kết cấu bề mặt, vùng hư hỏng, phần
vỏ còn sót và hình dạng bất thường của tép tỏi. Đây là những đặc điểm mà
phương pháp thủ công thường đánh giá bằng mắt thường, nhưng có thể được
mô hình hóa bằng các mạng nơ-ron tích chập hiện đại. Một số nghiên cứu
liên quan đến tỏi cũng đã chứng minh tiềm năng của học sâu trong xử lý
ảnh hậu thu hoạch, chẳng hạn phân loại tỏi sau cắt rễ bằng CNN hoặc phát
hiện hư hỏng tỏi bằng mô hình YOLO cải tiến {[}28{]}, {[}40{]}.

Bên cạnh nhu cầu thực tiễn, đề tài còn phù hợp với định hướng phát triển
công nghệ trí tuệ nhân tạo tại Việt Nam. Quyết định số 127/QĐ-TTg ngày
26/01/2021 của Thủ tướng Chính phủ đã ban hành Chiến lược quốc gia về
nghiên cứu, phát triển và ứng dụng trí tuệ nhân tạo đến năm 2030, trong
đó AI được định hướng ứng dụng vào nhiều lĩnh vực, bao gồm sản xuất và
nông nghiệp công nghệ cao
{[}2{]}.
Vì vậy, việc nghiên cứu và phát triển mô hình học sâu cho bài toán phân
loại chất lượng tỏi là cần thiết, có ý nghĩa khoa học và có khả năng ứng
dụng thực tiễn trong lĩnh vực chế biến nông sản.

Từ những lý do trên, luận văn lựa chọn đề tài ``\textbf{Nghiên cứu và
phát triển mô hình học sâu cho bài toán phân loại chất lượng tỏi}'' nhằm
xây dựng một mô hình phân loại ảnh có khả năng nhận diện ba trạng thái
chính của tỏi sau bóc vỏ, gồm \textbf{tỏi sạch}, \textbf{tỏi còn vỏ} và
\textbf{tỏi hư}, góp phần hỗ trợ tự động hóa dây chuyền sản xuất và nâng
cao chất lượng sản phẩm.

\begin{itemize}
\item
  \textbf{Tình hình nghiên cứu ngoài nước liên quan đến đề tài}
\end{itemize}

Trên thế giới, trí tuệ nhân tạo (AI) được ứng dụng rộng rãi trong nhiều
khía cạnh của nông nghiệp, từ quản lý cây trồng, chăn nuôi, tối ưu hóa
tài nguyên nước và đất, đến lai tạo giống cây trồng mới. Một số ứng dụng
và nghiên cứu tiêu biểu bao gồm:

\begin{itemize}
\item
  \textbf{Hệ thống thị giác nhân tạo cho kiểm tra và phân loại trái
  cây:} Nghiên cứu của Rojas Santelices và cộng sự năm 2025 tổng hợp 56
  công trình trong giai đoạn 2015--2024 về hệ thống thị giác nhân tạo
  dùng cho kiểm tra và phân loại trái cây. Kết quả cho thấy các hệ thống
  kiểm tra chất lượng nông sản hiện nay thường sử dụng camera RGB, hệ
  thống chiếu sáng LED, kết hợp với các thuật toán xử lý ảnh và mô hình
  học sâu như ResNet, VGG hoặc transfer learning. Công trình cũng chỉ ra
  rằng kiểm tra thủ công có nhiều hạn chế như tốn thời gian, mang tính
  chủ quan và dễ sai sót, trong khi hệ thống thị giác máy tính có thể hỗ
  trợ phát hiện lỗi, phân loại chất lượng và ra quyết định tự động trong
  dây chuyền sản xuất. Nghiên cứu này có ý nghĩa nền tảng đối với đề tài
  vì bài toán phân loại tỏi sau bóc vỏ cũng thuộc nhóm kiểm tra chất
  lượng nông sản bằng hình ảnh {[}38{]}.
\item
  \textbf{Thiết kế và thực nghiệm hệ thống phân loại chất lượng dưa leo
  dựa trên thị giác máy tính:} Liu và cộng sự năm 2024 đề xuất một hệ
  thống phân loại chất lượng dưa leo dựa trên thị giác máy tính và học
  sâu. Hệ thống sử dụng ảnh RGB để dự đoán khối lượng và hỗ trợ phân
  loại sản phẩm theo chất lượng. Kết quả thực nghiệm cho thấy hệ thống
  đạt năng suất tối đa 2,3 tấn/giờ; mô hình MassNet đạt sai số MAPE
  3,9\% và RMSE 6,7 g trong dự đoán khối lượng. Công trình này cho thấy
  việc kết hợp mô hình học sâu với thiết bị phân loại cơ khí có thể hỗ
  trợ triển khai trong dây chuyền sản xuất thực tế {[}39{]}.
\item
  \textbf{Phân loại ảnh tỏi sau cắt rễ bằng mạng nơ-ron tích chập sâu
  theo hướng phân loại đa nhãn và đa lớp:} Nhóm tác giả Pham Thi Quynh
  Anh, Dang Quoc Thuyet và Yuichi Kobayashi năm 2022 sử dụng mạng CNN để
  phân loại tỏi sau khi cắt rễ. Công trình phân loại tỏi theo hai nhóm
  nhãn: nhóm chất lượng cắt rễ gồm good, bad, untrimmed, scratched và
  nhóm trạng thái bề mặt gồm clean, muddy. Mô hình multi-label đạt độ
  chính xác tổng thể 95,2\%; khi bổ sung lớp nền, độ chính xác tăng lên
  98,0\%; thời gian xử lý mỗi ảnh khoảng 0,021 giây. Công trình này
  chứng minh rằng CNN có thể được sử dụng hiệu quả trong xử lý ảnh hậu
  thu hoạch tỏi và có tiềm năng triển khai cho robot phân loại tỏi thời
  gian thực {[}28{]}.
\item
  \textbf{Garlic-YOLO-DD: Thuật toán phát hiện đối tượng nhẹ cho bài
  toán phát hiện hư hỏng tỏi:} Gao và cộng sự năm 2026 đề xuất mô hình
  Garlic-YOLO-DD cho bài toán phát hiện hư hỏng tỏi. Mô hình được xây
  dựng dựa trên YOLOv11n và cải tiến bằng các thành phần như ADown,
  SimAM attention và BiFPN nhằm giảm số tham số, tăng khả năng nhận diện
  vùng hư hỏng nhỏ và cải thiện tốc độ suy luận. Kết quả cho thấy mô
  hình đạt mAP@50 là 92,58\%, chỉ sử dụng 1,497 triệu tham số và đạt tốc
  độ 167 FPS trên bộ dữ liệu hư hỏng tỏi tự xây dựng {[}40{]}.
\item
  \textbf{DH-GarlicNet: Phương pháp nhận diện chính xác hư hỏng tỏi dựa
  trên mạng phần dư cải tiến:} He và cộng sự năm 2026 đề xuất mô hình
  DH-GarlicNet để nhận diện tỏi hư dựa trên kiến trúc ResNet34 cải tiến.
  Mô hình tích hợp depthwise convolution, SE attention và hàm kích hoạt
  SiLU để tăng khả năng trích xuất đặc trưng vùng hư hỏng. Kết quả thực
  nghiệm cho thấy mô hình đạt độ chính xác 98,95\% và loss 0,0793 trong
  nhiệm vụ phát hiện hư hỏng tỏi. Phân tích Grad-CAM cho thấy mô hình có
  khả năng tập trung vào các vùng hư hỏng trên ảnh, chứng minh hiệu quả
  trong học đặc trưng cục bộ {[}41{]}.
\end{itemize}

\begin{itemize}
\item
  Các công trình trên cho thấy xu hướng nghiên cứu quốc tế hiện nay
  không chỉ dừng lại ở phân loại ảnh nông sản thông thường mà đang tập
  trung vào các mô hình có khả năng nhận diện lỗi nhỏ, học đặc trưng cục
  bộ, sử dụng attention và tối ưu mô hình để triển khai trong môi trường
  thực tế. Đây là cơ sở quan trọng cho đề tài phân loại chất lượng tỏi
  sau bóc vỏ, vì bài toán này cũng cần nhận diện các đặc điểm nhỏ như vỏ
  còn sót, vùng hư hỏng, vết dập hoặc thay đổi màu sắc trên bề mặt tép
  tỏi.
\end{itemize}

\begin{itemize}
\item
  \textbf{Tình hình nghiên cứu trong nước liên quan đến đề tài}
\end{itemize}

Tại Việt Nam, việc ứng dụng trí tuệ nhân tạo trong nông nghiệp đang ngày
càng được chú trọng. Các nghiên cứu và ứng dụng hiện nay tập trung nhiều
vào nhận diện sâu bệnh cây trồng, giám sát sản xuất, truy xuất nguồn
gốc, quản lý dinh dưỡng và hỗ trợ ra quyết định trong canh tác. Mặc dù
số lượng nghiên cứu trực tiếp về phân loại chất lượng tỏi sau bóc vỏ còn
hạn chế, nhiều công trình trong nước đã cho thấy tiềm năng lớn của AI và
thị giác máy tính trong nông nghiệp.

\begin{itemize}
\item
  \textbf{Ứng dụng "AI Plant Doctor":} Trong chẩn đoán bệnh trên cây
  trồng, ứng dụng "AI Plant Doctor" được phát triển bởi BOM-Software,
  phối hợp với Trung tâm Thúc đẩy Quốc tế hóa Khoa học và Công nghệ Việt
  Nam (VISTIP), Viện Nghiên cứu Cây ăn quả Miền Nam (SOFRI), và Trung
  tâm Quốc tế Sinh học Nông nghiệp (CABI) trong khuôn khổ chương trình
  Plantwise. Ứng dụng này sử dụng AI và công nghệ Big Data để hỗ trợ
  nhận diện sâu bệnh trên cây trồng, ban đầu tập trung vào cây thanh
  long. Ứng dụng ra mắt vào ngày 18/12/2020 tại tỉnh Tiền Giang, cho
  phép nông dân chụp ảnh cây trồng bằng điện thoại thông minh (tương
  thích cả Android và iOS) để nhận diện bệnh và nhận khuyến nghị xử lý.
  Hiện tại, ứng dụng đang được mở rộng để hỗ trợ thêm nhiều loại cây
  trồng khác giúp nông dân phát hiện sớm sâu bệnh, giảm thiểu thiệt hại,
  và cung cấp giải pháp xử lý kịp thời mà không cần phụ thuộc hoàn toàn
  vào chuyên gia tại chỗ {[}29{]}.
\item
  \textbf{Ứng dụng AI trong quản lý bệnh hại cây trồng:} Bài tổng hợp
  của CESTI năm 2025 trình bày nhiều ứng dụng AI trong quản lý bệnh hại
  cây trồng tại Việt Nam. Trong đó, ứng dụng AI Doctor tại Hưng Yên hỗ
  trợ nhận diện và quản lý bốn loại sâu bệnh phổ biến trên cây nhãn; ứng
  dụng 2Nông được tích hợp AI để chẩn đoán sâu bệnh và thiếu dinh dưỡng
  trên cà phê, tiêu, sầu riêng và lúa; hệ thống T-Pest sử dụng mô hình
  YOLOv5-Ghost để nhận diện sâu bệnh lúa theo thời gian thực. Ngoài ra,
  một hệ thống AI nhận diện bệnh hại dưa lưới trong nhà màng đạt độ
  chính xác trên 90\%. Các công trình này cho thấy AI tại Việt Nam đang
  được triển khai nhiều trong nhận diện hình ảnh nông nghiệp, tạo cơ sở
  thực tiễn cho việc nghiên cứu phân loại chất lượng tỏi {[}30{]}.
\item
  \textbf{Nghiên cứu thuật toán học sâu mạng nơ-ron tích chập và ứng
  dụng trong nhận dạng hình ảnh ong mật mang phấn:} Nghiên cứu của Học
  viện Nông nghiệp Việt Nam năm 2023 ứng dụng mạng nơ-ron tích chập
  trong nhận dạng ảnh ong mật mang phấn. Nhóm nghiên cứu khảo sát các
  kiến trúc CNN, học chuyển giao như VGG16, VGG19, ResNet50, ResNet101
  và xây dựng mô hình CNN cơ bản nhằm giảm chi phí tính toán. Nghiên cứu
  cũng sử dụng các kỹ thuật tăng cường dữ liệu như xoay ảnh, dịch
  chuyển, cắt, phóng to/thu nhỏ và lật ảnh để cải thiện khả năng tổng
  quát hóa. Công trình này liên quan đến đề tài vì cùng sử dụng CNN cho
  bài toán nhận dạng ảnh trong nông nghiệp, đồng thời cho thấy hướng
  thiết kế mô hình nhẹ, hiệu quả là cần thiết khi muốn triển khai thực
  tế {[}31{]}.
\item
  \textbf{Ứng dụng trí tuệ nhân tạo vào phát hiện độc tố vi nấm trong
  thực phẩm:} Bài viết của Khoa Công nghệ thực phẩm, Học viện Nông
  nghiệp Việt Nam năm 2025 trình bày xu hướng ứng dụng AI trong phát
  hiện độc tố vi nấm trong thực phẩm. Bài viết chỉ ra rằng các phương
  pháp truyền thống như LC-MS/MS và ELISA tuy có độ chính xác cao nhưng
  yêu cầu thiết bị hiện đại, quy trình phức tạp và nhân sự chuyên môn;
  trong khi đó, AI có khả năng xử lý dữ liệu lớn, tự động hóa nhận diện
  và hỗ trợ kiểm nghiệm nhanh hơn. Nội dung này liên quan gián tiếp đến
  đề tài vì phân loại tỏi hư, tỏi sạch và tỏi còn vỏ cũng thuộc nhóm bài
  toán kiểm soát chất lượng và an toàn trong chế biến thực phẩm {[}32{]}.
\end{itemize}

\begin{itemize}
\item
  Tuy nhiên, trong phạm vi các nguồn đã khảo sát, các nghiên cứu tại
  Việt Nam hiện nay chủ yếu tập trung vào phát hiện bệnh cây trồng, quản
  lý canh tác hoặc an toàn thực phẩm nói chung. Chưa có nhiều công trình
  tập trung trực tiếp vào bài toán phân loại chất lượng tỏi sau bóc vỏ
  trong dây chuyền chế biến. Điều này cho thấy đề tài có tính mới ở mức
  ứng dụng, đặc biệt trong bối cảnh các cơ sở chế biến tỏi vẫn còn phụ
  thuộc nhiều vào phân loại thủ công.
\end{itemize}

\begin{itemize}
\item
  \textbf{Khoảng trống nghiên cứu và lý do lựa chọn đề tài}
\end{itemize}

Từ các nghiên cứu ngoài nước và trong nước có thể nhận thấy rằng học sâu
đã chứng minh hiệu quả trong nhiều bài toán kiểm tra chất lượng nông
sản, phát hiện hư hỏng và phân loại sản phẩm sau thu hoạch. Tuy nhiên,
phần lớn các công trình hiện nay tập trung vào trái cây, rau củ, bệnh lá
cây hoặc hư hỏng trên củ tỏi nguyên củ. Các nghiên cứu trực tiếp về tỏi
chủ yếu xử lý bài toán cắt rễ, phát hiện hư hỏng củ hoặc nhận diện bệnh,
trong khi bài toán phân loại tỏi sau bóc vỏ thành các nhóm tỏi sạch, tỏi
còn vỏ và tỏi hư vẫn chưa được khai thác nhiều.

Bài toán phân loại tỏi sau bóc vỏ có những đặc thù riêng so với các bài
toán phân loại nông sản thông thường. Thứ nhất, đối tượng cần phân loại
có kích thước nhỏ và hình dạng không đồng nhất. Thứ hai, dấu hiệu phân
biệt giữa các lớp đôi khi rất tinh vi, chẳng hạn phần vỏ còn sót chỉ
chiếm một vùng nhỏ trên bề mặt tép tỏi. Thứ ba, tỏi hư có thể biểu hiện
qua nhiều dạng khác nhau như đổi màu, dập nát, mềm nhũn hoặc xuất hiện
vùng thâm. Những đặc điểm này đòi hỏi mô hình phải có khả năng khai thác
cả đặc trưng tổng thể và đặc trưng cục bộ.

Chính vì vậy, luận văn lựa chọn hướng nghiên cứu mô hình học sâu cho bài
toán phân loại chất lượng tỏi sau bóc vỏ. Việc phát triển mô hình có khả
năng phân loại chính xác ba nhóm \textbf{tỏi sạch}, \textbf{tỏi còn vỏ}
và \textbf{tỏi hư} không chỉ có ý nghĩa về mặt học thuật mà còn có khả
năng ứng dụng thực tiễn trong dây chuyền chế biến thực phẩm. Hệ thống
này có thể hỗ trợ giảm phụ thuộc vào lao động thủ công, nâng cao tính ổn
định của chất lượng sản phẩm và tạo nền tảng cho việc xây dựng dây
chuyền phân loại tỏi tự động trong tương lai.

\begin{enumerate}
\def\labelenumi{\arabic{enumi}.}
\setcounter{enumi}{1}
\item
  \protect\phantomsection\label{_Toc232111981}{}\textbf{Mục tiêu nghiên
  cứu}

  \begin{enumerate}
  \def\labelenumii{\arabic{enumii}.}
  \item
    \protect\phantomsection\label{_Toc232111982}{}\emph{\textbf{Mục tiêu
    tổng quát}}
  \end{enumerate}
\end{enumerate}

Mục tiêu tổng quát của luận văn là nghiên cứu và phát triển mô hình học
sâu nhằm phân loại chất lượng tỏi sau quá trình bóc vỏ dựa trên dữ liệu
hình ảnh. Mô hình được xây dựng với khả năng nhận diện và phân loại tỏi
thành ba nhóm chính: \textbf{tỏi đã bóc sạch}, \textbf{tỏi còn sót vỏ}
và \textbf{tỏi hư hỏng}. Thông qua đó, luận văn hướng đến việc hỗ trợ tự
động hóa công đoạn kiểm tra chất lượng tỏi trong dây chuyền chế biến
thực phẩm, góp phần giảm phụ thuộc vào lao động thủ công, nâng cao năng
suất và bảo đảm tính ổn định của sản phẩm đầu ra.

Bên cạnh việc xây dựng mô hình phân loại ảnh thông thường, luận văn còn
hướng đến đề xuất một mô hình học sâu cải tiến dựa trên kiến trúc CNN
phù hợp nhất sau quá trình thực nghiệm. Mô hình đề xuất không chỉ học
trực tiếp bài toán phân loại ba lớp, mà còn được bổ sung các nhiệm vụ
phụ có ý nghĩa theo logic đánh giá chất lượng tỏi, từ đó giúp mô hình
học đặc trưng tốt hơn và cải thiện hiệu quả phân loại.

\begin{enumerate}
\def\labelenumi{\arabic{enumi}.}
\setcounter{enumi}{1}
\item
  \protect\phantomsection\label{_Toc232111983}{}\emph{\textbf{Mục tiêu
  cụ thể}}
\end{enumerate}

Luận văn tập trung vào các mục tiêu cụ thể như sau:

\begin{enumerate}
\def\labelenumi{\alph{enumi}.}
\item
  \textbf{Thu thập, xây dựng và tiền xử lý bộ dữ liệu ảnh tỏi}
\end{enumerate}

Luận văn tiến hành thu thập và xây dựng bộ dữ liệu ảnh tỏi phục vụ cho
bài toán phân loại chất lượng sau bóc vỏ. Dữ liệu được tổ chức thành ba
nhóm chính gồm \textbf{Fully\_Peeled\_Garlic} tương ứng với tỏi đã bóc
sạch, \textbf{Partially\_Peeled\_Garlic} tương ứng với tỏi còn sót vỏ và
\textbf{Spoiled\_Garlic} tương ứng với tỏi hư hỏng.

Sau khi thu thập, dữ liệu được kiểm tra, làm sạch và tiền xử lý nhằm bảo
đảm chất lượng đầu vào cho quá trình huấn luyện mô hình. Các bước tiền
xử lý có thể bao gồm chuẩn hóa kích thước ảnh, kiểm tra nhãn, loại bỏ
ảnh lỗi, chia dữ liệu thành các tập huấn luyện, kiểm định và kiểm tra.
Ngoài ra, luận văn cũng xem xét áp dụng các kỹ thuật tăng cường dữ liệu
nhằm giúp mô hình học được nhiều biến thể khác nhau của ảnh tỏi trong
điều kiện thực tế.

\begin{enumerate}
\def\labelenumi{\alph{enumi}.}
\setcounter{enumi}{1}
\item
  \textbf{Thực nghiệm và so sánh các mô hình CNN gốc}
\end{enumerate}

Luận văn tiến hành thực nghiệm trên nhiều mô hình học sâu thuộc nhóm
mạng nơ-ron tích chập nhằm đánh giá khả năng phân loại ảnh tỏi. Các mô
hình được lựa chọn bao gồm \textbf{NASNetMobile}, \textbf{InceptionV3},
\textbf{EfficientNetB4}, \textbf{MobileNetV2} và \textbf{DenseNet121}.
Đây là các kiến trúc CNN đã được sử dụng phổ biến trong bài toán phân
loại ảnh và có khả năng học đặc trưng hình ảnh ở nhiều mức độ khác nhau.

Thông qua quá trình thực nghiệm, luận văn so sánh hiệu quả của các mô
hình trên cùng bộ dữ liệu và cùng điều kiện đánh giá. Mục tiêu của bước
này là xác định mô hình nền phù hợp nhất với bài toán phân loại chất
lượng tỏi, làm cơ sở để phát triển mô hình đề xuất ở các bước tiếp theo.

\begin{enumerate}
\def\labelenumi{\alph{enumi}.}
\setcounter{enumi}{2}
\item
  \textbf{Lựa chọn mô hình phù hợp nhất làm backbone cho mô hình đề
  xuất}
\end{enumerate}

Dựa trên kết quả thực nghiệm giữa các mô hình CNN nền, luận văn lựa chọn
mô hình có hiệu quả tốt nhất để làm backbone chính. Trong định hướng
hiện tại, \textbf{EfficientNetB4} được lựa chọn làm backbone do có khả
năng trích xuất đặc trưng ảnh tốt, cân bằng giữa độ chính xác và chi phí
tính toán, đồng thời cho kết quả phù hợp với bài toán phân loại chất
lượng tỏi.

EfficientNetB4 được sử dụng để trích xuất các đặc trưng quan trọng từ
ảnh tỏi, bao gồm màu sắc, kết cấu bề mặt, vùng hư hỏng, phần vỏ còn sót
và hình dạng bất thường. Các đặc trưng này là cơ sở để mô hình đưa ra
quyết định phân loại ảnh vào một trong ba nhóm chất lượng.

\begin{enumerate}
\def\labelenumi{\alph{enumi}.}
\setcounter{enumi}{3}
\item
  \textbf{Đề xuất mô hình học sâu cải tiến cho bài toán phân loại chất
  lượng tỏi}
\end{enumerate}

Trên nền tảng EfficientNetB4, luận văn đề xuất mô hình học sâu cải tiến
theo hướng kết hợp đặc trưng và học đa nhiệm. Mô hình đề xuất sử dụng
EfficientNetB4 làm backbone chính để trích xuất đặc trưng ảnh, sau đó
kết hợp hai cơ chế tổng hợp đặc trưng là \textbf{Global Average Pooling
(GAP)} và \textbf{Global Max Pooling (GMP)}. Trong đó, GAP giúp mô hình
nắm bắt thông tin tổng quát trên toàn ảnh, còn GMP giúp giữ lại các đặc
trưng nổi bật, chẳng hạn vùng hư hỏng hoặc phần vỏ còn sót trên bề mặt
tép tỏi.

Khác với kiến trúc phân loại thông thường chỉ sử dụng một đầu ra Softmax
cho ba lớp, mô hình đề xuất bổ sung thêm hai nhiệm vụ phụ nhằm hỗ trợ mô
hình học theo logic phân loại thực tế của bài toán tỏi. Cụ thể, mô hình
gồm một nhiệm vụ chính và hai nhiệm vụ phụ:

\begin{itemize}
\item
  Nhiệm vụ chính là phân loại ảnh tỏi thành ba lớp:
  \textbf{Fully\_Peeled\_Garlic}, \textbf{Partially\_Peeled\_Garlic} và
  \textbf{Spoiled\_Garlic}.
\item
  Nhiệm vụ phụ thứ nhất là nhận diện ảnh tỏi thuộc nhóm Spoiled hay
  Non-spoiled, tức là xác định mẫu tỏi có bị hư hỏng hay không.
\item
  Nhiệm vụ phụ thứ hai là phân biệt Fully\_Peeled và Partially\_Peeled
  đối với các ảnh thuộc nhóm không hư hỏng, nhằm giúp mô hình học rõ hơn
  sự khác biệt giữa tỏi đã bóc sạch và tỏi còn sót vỏ.
\end{itemize}

Với thiết kế này, mô hình không chỉ học câu hỏi ``\emph{\textbf{ảnh này
thuộc lớp nào}}'', mà còn học thêm các câu hỏi phụ có ý nghĩa về mặt
chất lượng: ``\emph{\textbf{ảnh này có bị hư không?}}'' và
``\emph{\textbf{nếu không hư thì ảnh này là tỏi bóc sạch hay tỏi còn
vỏ?}}''. Cách tiếp cận này được kỳ vọng giúp mô hình học đặc trưng sâu
hơn, giảm nhầm lẫn giữa các lớp có biểu hiện gần giống nhau và cải thiện
hiệu quả phân loại tổng thể.

\begin{enumerate}
\def\labelenumi{\alph{enumi}.}
\setcounter{enumi}{4}
\item
  \textbf{Xử lý vấn đề mất cân bằng dữ liệu bằng hàm mất mát phù hợp}
\end{enumerate}

Trong bài toán phân loại chất lượng tỏi, số lượng mẫu giữa các lớp có
thể không hoàn toàn cân bằng. Thông thường, nhóm tỏi sạch có thể chiếm
tỷ lệ lớn hơn so với nhóm tỏi còn vỏ hoặc tỏi hư. Sự mất cân bằng này có
thể làm mô hình thiên lệch về lớp có nhiều mẫu, dẫn đến giảm khả năng
nhận diện các lớp ít mẫu nhưng quan trọng.

Do đó, luận văn xem xét sử dụng \textbf{Class-Balanced Loss} hoặc biến
thể kết hợp với Focal Loss trong quá trình huấn luyện. Mục tiêu của việc
sử dụng hàm mất mát này là giúp mô hình chú ý hơn đến các lớp có ít mẫu
và các mẫu khó phân loại, từ đó cải thiện khả năng nhận diện tỏi hư và
tỏi còn sót vỏ.

\begin{enumerate}
\def\labelenumi{\alph{enumi}.}
\setcounter{enumi}{5}
\item
  \textbf{Đánh giá và phân tích hiệu quả của mô hình}
\end{enumerate}

Luận văn đánh giá hiệu quả của các mô hình thông qua các chỉ số chuẩn
trong bài toán phân loại ảnh, bao gồm \textbf{Accuracy},
\textbf{Precision}, \textbf{Recall}, \textbf{F1-score}, \textbf{AUC} và
\textbf{Confusion Matrix}. Các chỉ số này giúp đánh giá toàn diện hiệu
quả của mô hình, không chỉ ở độ chính xác tổng thể mà còn ở khả năng
nhận diện từng lớp cụ thể.

Đặc biệt, ma trận nhầm lẫn được sử dụng để phân tích các lỗi phân loại
thường gặp, chẳng hạn mô hình nhầm tỏi còn vỏ thành tỏi sạch, hoặc nhầm
tỏi hư thành tỏi sạch. Đây là những lỗi có ý nghĩa thực tiễn quan trọng
vì có thể ảnh hưởng trực tiếp đến chất lượng sản phẩm trong dây chuyền
chế biến.

\begin{enumerate}
\def\labelenumi{\alph{enumi}.}
\setcounter{enumi}{6}
\item
  \textbf{Đề xuất định hướng ứng dụng thực tiễn}
\end{enumerate}

Từ kết quả nghiên cứu và thực nghiệm, luận văn đề xuất định hướng ứng
dụng mô hình vào hệ thống phân loại chất lượng tỏi tự động. Hệ thống có
thể được phát triển theo hướng kết hợp camera, mô hình học sâu và cơ cấu
phân loại trên dây chuyền sản xuất. Khi được triển khai, mô hình có thể
hỗ trợ phát hiện và phân loại tỏi sạch, tỏi còn vỏ và tỏi hư một cách
nhanh chóng, ổn định hơn so với phương pháp thủ công.

Kết quả của luận văn có thể làm cơ sở cho việc xây dựng hệ thống kiểm
tra chất lượng tỏi trong thực tế, đồng thời có thể mở rộng sang các loại
nông sản khác có đặc điểm tương tự như hành, gừng, khoai tây hoặc các
sản phẩm cần kiểm tra bề mặt sau sơ chế.

\begin{enumerate}
\def\labelenumi{\arabic{enumi}.}
\setcounter{enumi}{2}
\item
  \protect\phantomsection\label{_Toc232111984}{}\textbf{Đối tượng, phạm
  vi nghiên cứu}

  \begin{enumerate}
  \def\labelenumii{\arabic{enumii}.}
  \item
    \protect\phantomsection\label{_Toc232111985}{}\emph{\textbf{Đối
    tượng nghiên cứu}}
  \end{enumerate}
\end{enumerate}

Đối tượng nghiên cứu của luận văn là các mô hình học sâu ứng dụng trong
nhận diện hình ảnh, cụ thể là bài toán phân loại chất lượng tỏi sau quá
trình bóc vỏ. Nghiên cứu tập trung xây dựng, huấn luyện, tối ưu và đánh
giá các mô hình phân loại ảnh nhằm nhận diện ba nhóm tỏi chính gồm:
\textbf{tỏi đã bóc sạch}, \textbf{tỏi còn sót vỏ} và \textbf{tỏi hư
hỏng}.

Cụ thể, đối tượng nghiên cứu của luận văn bao gồm:

\begin{enumerate}
\def\labelenumi{\alph{enumi}.}
\item
  \textbf{Hình ảnh dữ liệu tỏi sau bóc vỏ}
\end{enumerate}

Luận văn sử dụng bộ dữ liệu hình ảnh tỏi được chia thành ba nhóm chính:
\textbf{Fully\_Peeled\_Garlic}, \textbf{Partially\_Peeled\_Garlic} và
\textbf{Spoiled\_Garlic}. Trong đó, Fully\_Peeled\_Garlic là nhóm tỏi đã
được bóc vỏ hoàn chỉnh, bề mặt tương đối sạch và không có dấu hiệu hư
hỏng rõ rệt; Partially\_Peeled\_Garlic là nhóm tỏi vẫn còn sót lại một
phần vỏ sau quá trình bóc; Spoiled\_Garlic là nhóm tỏi có dấu hiệu hư
hỏng, dập nát, biến màu hoặc không đạt yêu cầu chất lượng.

Bộ dữ liệu này được sử dụng để huấn luyện, kiểm định và đánh giá hiệu
quả của các mô hình học sâu. Hình ảnh trong dữ liệu đóng vai trò là
nguồn đầu vào chính để mô hình học các đặc trưng trực quan như màu sắc,
hình dạng, kết cấu bề mặt, vùng hư hỏng và phần vỏ còn sót trên tép tỏi.

\begin{enumerate}
\def\labelenumi{\alph{enumi}.}
\setcounter{enumi}{1}
\item
  \textbf{Các mô hình học sâu dùng trong phân loại ảnh}
\end{enumerate}

Luận văn tập trung nghiên cứu các mô hình mạng nơ-ron tích chập, hay còn
gọi là \textbf{Convolutional Neural Network -- CNN}, trong bài toán phân
loại chất lượng tỏi. Các mô hình CNN nền được lựa chọn để thực nghiệm và
so sánh bao gồm \textbf{NASNetMobile}, \textbf{InceptionV3},
\textbf{EfficientNetB4}, \textbf{MobileNetV2} và \textbf{DenseNet121}.

Các mô hình này được lựa chọn vì đã được sử dụng rộng rãi trong các bài
toán phân loại ảnh, có khả năng trích xuất đặc trưng hình ảnh ở nhiều
mức độ khác nhau và phù hợp để đánh giá hiệu quả trên bộ dữ liệu tỏi.
Thông qua quá trình thực nghiệm, luận văn xác định mô hình có kết quả
phù hợp nhất để làm cơ sở phát triển mô hình đề xuất.

\begin{enumerate}
\def\labelenumi{\alph{enumi}.}
\setcounter{enumi}{2}
\item
  \textbf{Mô hình đề xuất cho bài toán phân loại chất lượng tỏi}
\end{enumerate}

Trên cơ sở kết quả thực nghiệm các mô hình CNN nền, luận văn lựa chọn
\textbf{EfficientNetB4} làm backbone chính để phát triển mô hình đề
xuất. Mô hình đề xuất có cấu trúc dựa trên EfficientNetB4 kết hợp với
\textbf{Global Average Pooling (GAP)}, \textbf{Global Max Pooling
(GMP)}, nhánh phân loại chính ba lớp và hai nhánh phụ gồm
\textbf{Spoilage Aux} và \textbf{Peeling Aux}.

Trong đó, nhánh phân loại chính thực hiện nhiệm vụ phân loại ảnh tỏi
thành ba lớp: Fully\_Peeled\_Garlic, Partially\_Peeled\_Garlic và
Spoiled\_Garlic. Nhánh phụ Spoilage Aux hỗ trợ mô hình học khả năng phân
biệt tỏi hư và tỏi không hư. Nhánh phụ Peeling Aux hỗ trợ mô hình phân
biệt tỏi đã bóc sạch và tỏi còn sót vỏ đối với các mẫu không thuộc nhóm
tỏi hư.

Cách thiết kế này giúp mô hình không chỉ học trực tiếp nhãn phân loại ba
lớp, mà còn học thêm logic đánh giá chất lượng tỏi theo từng bước: trước
hết xác định tỏi có bị hư hay không, sau đó nếu không hư thì tiếp tục
phân biệt tỏi đã bóc sạch hay còn sót vỏ. Đây là hướng tiếp cận phù hợp
với đặc thù của bài toán, vì các dấu hiệu như vùng hư hỏng hoặc phần vỏ
còn sót thường có kích thước nhỏ và dễ gây nhầm lẫn.

\begin{enumerate}
\def\labelenumi{\alph{enumi}.}
\setcounter{enumi}{3}
\item
  \textbf{Các phương pháp xử lý và tối ưu mô hình}
\end{enumerate}

Luận văn nghiên cứu và áp dụng các phương pháp xử lý dữ liệu, tăng cường
dữ liệu và tối ưu mô hình nhằm nâng cao hiệu quả phân loại. Các bước xử
lý có thể bao gồm chuẩn hóa kích thước ảnh, kiểm tra nhãn, chia dữ liệu
thành các tập huấn luyện, kiểm định và kiểm tra. Bên cạnh đó, các kỹ
thuật \textbf{Data Augmentation} như xoay ảnh, lật ảnh, thay đổi độ
sáng, phóng to hoặc dịch chuyển ảnh được sử dụng nhằm giúp mô hình tăng
khả năng khái quát hóa và hạn chế hiện tượng quá khớp.

Ngoài ra, do dữ liệu giữa các lớp có thể không hoàn toàn cân bằng, luận
văn sử dụng hoặc xem xét sử dụng \textbf{Class-Balanced Loss (CBLoss)}
nhằm giảm ảnh hưởng của mất cân bằng dữ liệu. Kỹ thuật này giúp mô hình
chú ý hơn đến các lớp có số lượng mẫu ít hơn, đặc biệt là nhóm tỏi còn
sót vỏ và tỏi hư hỏng.

\begin{enumerate}
\def\labelenumi{\alph{enumi}.}
\setcounter{enumi}{4}
\item
  \textbf{Phạm vi ứng dụng của mô hình}
\end{enumerate}

Đối tượng ứng dụng của nghiên cứu là công đoạn phân loại chất lượng tỏi
sau bóc vỏ trong dây chuyền chế biến thực phẩm. Mô hình được kỳ vọng có
thể hỗ trợ nhận diện và phân loại tự động các mẫu tỏi dựa trên ảnh, từ
đó giảm phụ thuộc vào lao động thủ công, nâng cao năng suất và bảo đảm
tính ổn định của chất lượng sản phẩm.

Ngoài bài toán phân loại tỏi, hướng nghiên cứu này còn có tiềm năng mở
rộng sang các loại nông sản khác có đặc điểm tương tự như hành, gừng,
khoai tây hoặc các sản phẩm cần kiểm tra bề mặt sau sơ chế.

\begin{enumerate}
\def\labelenumi{\arabic{enumi}.}
\item
  \protect\phantomsection\label{_Toc232111986}{}\emph{\textbf{Phạm vi
  nghiên cứu}}
\end{enumerate}

\begin{enumerate}
\def\labelenumi{\alph{enumi}.}
\item
  \textbf{Phạm vi dữ liệu}
\end{enumerate}

Dữ liệu sử dụng trong luận văn là bộ dữ liệu ảnh tỏi sau bóc vỏ, được
phân chia thành ba lớp chính: \textbf{Fully\_Peeled\_Garlic},
\textbf{Partially\_Peeled\_Garlic} và \textbf{Spoiled\_Garlic}. Bộ dữ
liệu phản ánh các trạng thái chất lượng thường gặp của tỏi sau công đoạn
bóc vỏ, bao gồm tỏi đã bóc sạch, tỏi còn sót vỏ và tỏi bị hư hỏng.

Trong phạm vi nghiên cứu này, dữ liệu được sử dụng dưới dạng ảnh tĩnh.
Luận văn tập trung vào việc xử lý, huấn luyện và đánh giá mô hình trên
các ảnh đã được gán nhãn sẵn. Nghiên cứu không đi sâu vào quá trình thu
hoạch, bảo quản, tách tép hoặc vận hành máy bóc vỏ tỏi. Các yếu tố như
điều kiện ánh sáng, góc chụp, phông nền và chất lượng ảnh được xem xét ở
mức phục vụ cho quá trình tiền xử lý và tăng cường dữ liệu.

\begin{enumerate}
\def\labelenumi{\alph{enumi}.}
\setcounter{enumi}{1}
\item
  \textbf{Phạm vi bài toán}
\end{enumerate}

Luận văn tập trung vào bài toán \textbf{phân loại ảnh tỏi thành ba lớp}.
Đầu vào của mô hình là ảnh tỏi sau bóc vỏ, đầu ra là nhãn phân loại
tương ứng với một trong ba nhóm: tỏi đã bóc sạch, tỏi còn sót vỏ hoặc
tỏi hư hỏng.

Nghiên cứu chưa đi sâu vào bài toán phát hiện vị trí của từng tép tỏi
trên ảnh, phân đoạn vùng hư hỏng, định vị chính xác phần vỏ còn sót hoặc
điều khiển cơ cấu phân loại tự động trên băng chuyền. Các bài toán này
có thể được xem là hướng phát triển tiếp theo sau khi mô hình phân loại
đạt kết quả ổn định.

\begin{enumerate}
\def\labelenumi{\alph{enumi}.}
\setcounter{enumi}{2}
\item
  \textbf{Phạm vi kỹ thuật}
\end{enumerate}

Về mặt kỹ thuật, luận văn tập trung nghiên cứu và ứng dụng các mô hình
học sâu, đặc biệt là mạng nơ-ron tích chập, cho bài toán phân loại ảnh.
Các mô hình CNN nền được đưa vào thực nghiệm gồm \textbf{NASNetMobile},
\textbf{InceptionV3}, \textbf{EfficientNetB4}, \textbf{MobileNetV2} và
\textbf{DenseNet121}. Các mô hình này được huấn luyện và đánh giá trên
cùng bộ dữ liệu để so sánh hiệu quả phân loại.

Sau quá trình thực nghiệm, luận văn lựa chọn mô hình phù hợp nhất để
phát triển mô hình đề xuất. Trong định hướng nghiên cứu hiện tại,
EfficientNetB4 được sử dụng làm backbone chính. Mô hình đề xuất kết hợp
EfficientNetB4 với GAP/GMP, nhánh phân loại chính ba lớp và hai nhánh
phụ gồm Spoilage Aux và Peeling Aux. Ngoài ra, luận văn sử dụng CBLoss
nhằm hỗ trợ xử lý vấn đề mất cân bằng dữ liệu giữa các lớp.

Phạm vi kỹ thuật của luận văn giới hạn trong các phương pháp học sâu cho
phân loại ảnh. Các kỹ thuật xử lý ảnh truyền thống chỉ đóng vai trò hỗ
trợ tiền xử lý dữ liệu, không phải hướng nghiên cứu chính.

\begin{enumerate}
\def\labelenumi{\alph{enumi}.}
\setcounter{enumi}{3}
\item
  \textbf{Phạm vi thực nghiệm}
\end{enumerate}

Các thực nghiệm của luận văn được thực hiện trong môi trường lập trình
Python với các thư viện hỗ trợ học sâu như TensorFlow/Keras hoặc
PyTorch, tùy theo quá trình triển khai mô hình. Việc huấn luyện mô hình
được thực hiện trên môi trường có hỗ trợ GPU nhằm tăng tốc quá trình
tính toán.

Trong quá trình thực nghiệm, dữ liệu được chia thành các tập huấn luyện,
kiểm định và kiểm tra. Mô hình được huấn luyện trên tập huấn luyện, tinh
chỉnh và lựa chọn thông qua tập kiểm định, sau đó đánh giá cuối cùng
trên tập kiểm tra. Các kỹ thuật như lưu mô hình tốt nhất, theo dõi quá
trình huấn luyện thông qua đồ thị Accuracy/Loss và điều chỉnh siêu tham
số được sử dụng để bảo đảm quá trình thực nghiệm có tính hệ thống.

\begin{enumerate}
\def\labelenumi{\alph{enumi}.}
\setcounter{enumi}{4}
\item
  \textbf{Phạm vi đánh giá}
\end{enumerate}

Luận văn đánh giá mô hình thông qua các chỉ số phổ biến trong bài toán
phân loại ảnh như \textbf{Accuracy}, \textbf{Precision},
\textbf{Recall}, \textbf{F1-score}, \textbf{AUC} và \textbf{Confusion
Matrix}. Trong đó, Accuracy phản ánh độ chính xác tổng thể của mô hình;
Precision và Recall giúp đánh giá khả năng dự đoán đúng từng lớp;
F1-score thể hiện sự cân bằng giữa Precision và Recall; AUC hỗ trợ đánh
giá khả năng phân biệt của mô hình; Confusion Matrix giúp phân tích chi
tiết các trường hợp mô hình nhầm lẫn giữa các lớp.

Đối với bài toán phân loại chất lượng tỏi, việc phân tích lỗi có ý nghĩa
quan trọng. Một số lỗi như nhầm tỏi hư thành tỏi sạch hoặc nhầm tỏi còn
vỏ thành tỏi sạch có thể ảnh hưởng trực tiếp đến chất lượng sản phẩm đầu
ra. Do đó, ngoài kết quả tổng thể, luận văn còn chú trọng phân tích hiệu
quả nhận diện trên từng lớp để đánh giá tính phù hợp của mô hình trong
thực tế.

\begin{enumerate}
\def\labelenumi{\alph{enumi}.}
\setcounter{enumi}{5}
\item
  \textbf{Phạm vi ứng dụng thực tiễn}
\end{enumerate}

Luận văn hướng đến việc xây dựng mô hình làm nền tảng cho hệ thống phân
loại chất lượng tỏi tự động. Trong phạm vi nghiên cứu, luận văn tập
trung vào phần mô hình nhận diện và phân loại ảnh. Việc tích hợp mô hình
vào hệ thống camera, băng chuyền, bộ điều khiển và cơ cấu gạt sản phẩm
chưa được triển khai đầy đủ trong phạm vi luận văn này.

Tuy nhiên, kết quả nghiên cứu có thể được sử dụng làm cơ sở cho các
nghiên cứu tiếp theo nhằm xây dựng hệ thống phân loại tỏi tự động trong
dây chuyền chế biến thực phẩm, góp phần giảm lao động thủ công, nâng cao
năng suất và cải thiện chất lượng sản phẩm.

\begin{enumerate}
\def\labelenumi{\arabic{enumi}.}
\setcounter{enumi}{3}
\item
  \protect\phantomsection\label{_Toc232111987}{}\textbf{Phương pháp
  nghiên cứu}
\end{enumerate}

Luận văn được thực hiện dựa trên sự kết hợp giữa phương pháp nghiên cứu
lý thuyết và phương pháp nghiên cứu thực nghiệm. Trong đó, nghiên cứu lý
thuyết được sử dụng để khảo sát các công trình liên quan, tìm hiểu cơ sở
khoa học của các mô hình học sâu trong phân loại ảnh. Nghiên cứu thực
nghiệm được sử dụng để xây dựng bộ dữ liệu, huấn luyện, so sánh, cải
tiến và đánh giá mô hình phân loại chất lượng tỏi.

Cụ thể, luận văn áp dụng các phương pháp nghiên cứu sau:

\begin{enumerate}
\def\labelenumi{\alph{enumi}.}
\item
  \textbf{Phương pháp nghiên cứu tài liệu}
\end{enumerate}

Luận văn tiến hành thu thập, phân tích và tổng hợp các tài liệu, công
trình nghiên cứu trong và ngoài nước liên quan đến trí tuệ nhân tạo, thị
giác máy tính, học sâu và ứng dụng học sâu trong nông nghiệp. Nội dung
khảo sát tập trung vào các hướng nghiên cứu chính như phân loại ảnh nông
sản, kiểm tra chất lượng thực phẩm, phát hiện lỗi bề mặt, nhận diện hư
hỏng sau thu hoạch và các mô hình học sâu dùng trong phân loại ảnh.

Bên cạnh đó, luận văn cũng nghiên cứu các kiến trúc mạng nơ-ron tích
chập hiện đại như NASNetMobile, InceptionV3, EfficientNetB4, MobileNetV2
và DenseNet121. Đây là các mô hình đã được sử dụng phổ biến trong nhiều
bài toán phân loại ảnh và có thể áp dụng cho bài toán phân loại chất
lượng tỏi. Ngoài ra, các kỹ thuật như Transfer Learning, Data
Augmentation, Global Average Pooling, Global Max Pooling, học đa nhiệm
và hàm mất mát xử lý mất cân bằng dữ liệu cũng được tìm hiểu để làm cơ
sở cho việc đề xuất mô hình cải tiến.

\begin{enumerate}
\def\labelenumi{\alph{enumi}.}
\setcounter{enumi}{1}
\item
  \textbf{Phương pháp thu thập, xử lý và chuẩn bị dữ liệu}
\end{enumerate}

Luận văn sử dụng bộ dữ liệu ảnh tỏi được phân chia thành ba lớp chính
gồm Fully\_Peeled\_Garlic, Partially\_Peeled\_Garlic và Spoiled\_Garlic.
Đây là ba trạng thái chất lượng thường gặp của tỏi sau quá trình bóc vỏ,
tương ứng với tỏi đã bóc sạch, tỏi còn sót vỏ và tỏi hư hỏng.

Dữ liệu sau khi thu thập được tiến hành kiểm tra và tiền xử lý trước khi
đưa vào huấn luyện mô hình. Các bước xử lý bao gồm kiểm tra nhãn dữ
liệu, loại bỏ các ảnh lỗi hoặc không phù hợp, chuẩn hóa kích thước ảnh,
chuẩn hóa giá trị điểm ảnh và tổ chức dữ liệu theo đúng cấu trúc từng
lớp. Sau đó, bộ dữ liệu được chia thành các tập train, validation và
test nhằm phục vụ quá trình huấn luyện, tinh chỉnh và đánh giá mô hình.

Để tăng khả năng khái quát hóa của mô hình và hạn chế hiện tượng quá
khớp, luận văn áp dụng các kỹ thuật tăng cường dữ liệu như xoay ảnh, lật
ảnh, thay đổi độ sáng, phóng to, thu nhỏ hoặc dịch chuyển ảnh. Các kỹ
thuật này giúp mô hình học được nhiều biến thể khác nhau của ảnh tỏi, từ
đó hoạt động ổn định hơn trong điều kiện dữ liệu thực tế.

Ngoài ra, luận văn cũng xem xét vấn đề mất cân bằng dữ liệu giữa các
lớp. Trong bài toán phân loại chất lượng tỏi, số lượng ảnh tỏi sạch, tỏi
còn vỏ và tỏi hư có thể không đồng đều. Vì vậy, việc xử lý mất cân bằng
dữ liệu là cần thiết để tránh mô hình thiên lệch về lớp có số lượng mẫu
lớn.

\begin{enumerate}
\def\labelenumi{\alph{enumi}.}
\setcounter{enumi}{2}
\item
  \textbf{Phương pháp thực nghiệm các mô hình CNN nền}
\end{enumerate}

Sau khi dữ liệu được chuẩn bị, luận văn tiến hành huấn luyện và đánh giá
nhiều mô hình CNN nền khác nhau nhằm xác định mô hình phù hợp nhất với
bài toán phân loại chất lượng tỏi. Các mô hình được sử dụng trong thực
nghiệm bao gồm:

\begin{itemize}
\item
  NASNetMobile
\item
  InceptionV3
\item
  EfficientNetB4
\item
  MobileNetV2
\item
  DenseNet121
\end{itemize}

Các mô hình này được huấn luyện trên cùng bộ dữ liệu và đánh giá theo
cùng hệ thống chỉ số để bảo đảm tính công bằng trong quá trình so sánh.
Mục tiêu của bước thực nghiệm này là xác định mô hình có khả năng trích
xuất đặc trưng ảnh tỏi tốt nhất, đặc biệt trong việc phân biệt các đặc
điểm như vùng hư hỏng, phần vỏ còn sót và trạng thái bề mặt của tép tỏi.

Kết quả thực nghiệm giữa các mô hình CNN nền là cơ sở để lựa chọn
backbone cho mô hình đề xuất. Trong hướng nghiên cứu hiện tại,
EfficientNetB4 được lựa chọn làm mô hình nền chính vì cho thấy khả năng
phù hợp tốt với dữ liệu ảnh tỏi và có khả năng trích xuất đặc trưng
mạnh.

\begin{enumerate}
\def\labelenumi{\alph{enumi}.}
\setcounter{enumi}{3}
\item
  \textbf{Phương pháp đề xuất và cải tiến mô hình}
\end{enumerate}

Trên cơ sở mô hình EfficientNetB4, luận văn đề xuất mô hình học sâu cải
tiến cho bài toán phân loại chất lượng tỏi. Mô hình sử dụng
EfficientNetB4 làm backbone chính để trích xuất đặc trưng ảnh, sau đó
kết hợp hai cơ chế tổng hợp đặc trưng là Global Average Pooling (GAP) và
Global Max Pooling (GMP).

Trong đó, GAP giúp mô hình nắm bắt thông tin tổng quát trên toàn ảnh,
còn GMP giúp giữ lại các đặc trưng nổi bật nhất, chẳng hạn như vùng hư
hỏng hoặc phần vỏ còn sót. Việc kết hợp GAP và GMP giúp mô hình khai
thác đồng thời đặc trưng tổng thể và đặc trưng cục bộ, phù hợp với đặc
điểm của ảnh tỏi sau bóc vỏ.

Khác với mô hình phân loại thông thường chỉ có một đầu ra Softmax ba
lớp, mô hình đề xuất được thiết kế theo hướng học đa nhiệm. Cụ thể, mô
hình gồm một nhiệm vụ chính và hai nhiệm vụ phụ:

Nhiệm vụ chính là phân loại ảnh tỏi thành ba lớp: Fully\_Peeled\_Garlic,
Partially\_Peeled\_Garlic và Spoiled\_Garlic.

Nhiệm vụ phụ thứ nhất, gọi là Spoilage Aux, có nhiệm vụ phân biệt tỏi hư
và tỏi không hư. Nhánh này giúp mô hình học rõ hơn các đặc trưng liên
quan đến hư hỏng như vùng đổi màu, dập nát hoặc bề mặt bất thường.

Nhiệm vụ phụ thứ hai, gọi là Peeling Aux, có nhiệm vụ phân biệt tỏi đã
bóc sạch và tỏi còn sót vỏ trên nhóm ảnh không hư hỏng. Nhánh này giúp
mô hình tập trung hơn vào đặc trưng của phần vỏ còn sót, vốn có thể xuất
hiện ở vùng nhỏ và dễ bị nhầm với bề mặt tỏi bình thường.

Với thiết kế này, mô hình không chỉ học trực tiếp nhãn phân loại ba lớp
mà còn học thêm logic phân loại theo quy trình kiểm tra chất lượng thực
tế: trước hết xác định tỏi có bị hư hay không, sau đó nếu không hư thì
tiếp tục xác định tỏi đã bóc sạch hay còn sót vỏ. Cách tiếp cận này được
kỳ vọng giúp mô hình giảm nhầm lẫn giữa các lớp và cải thiện hiệu quả
phân loại tổng thể.

\begin{enumerate}
\def\labelenumi{\alph{enumi}.}
\setcounter{enumi}{4}
\item
  \textbf{Phương pháp tối ưu hóa quá trình huấn luyện}
\end{enumerate}

Trong quá trình huấn luyện, luận văn tiến hành điều chỉnh các tham số
như số epoch, batch size, learning rate, optimizer và các tham số liên
quan đến quá trình fine-tuning. Việc điều chỉnh này nhằm giúp mô hình
hội tụ ổn định, hạn chế quá khớp và đạt hiệu quả tốt trên tập kiểm tra.

Bên cạnh đó, luận văn sử dụng hàm mất mát phù hợp với bài toán có khả
năng mất cân bằng dữ liệu. Cụ thể, Class-Balanced Loss hoặc biến thể kết
hợp với Focal Loss được sử dụng nhằm giảm ảnh hưởng của sự chênh lệch số
lượng mẫu giữa các lớp. Hàm mất mát này giúp mô hình chú ý hơn đến các
lớp ít mẫu hoặc khó phân loại, đặc biệt là nhóm
Partially\_Peeled\_Garlic và Spoiled\_Garlic.

Đối với mô hình học đa nhiệm, tổng hàm mất mát được xây dựng dựa trên sự
kết hợp giữa mất mát của nhiệm vụ chính và mất mát của các nhiệm vụ phụ.
Trong đó, nhánh phân loại chính giữ vai trò trung tâm, còn hai nhánh phụ
đóng vai trò hỗ trợ quá trình học đặc trưng. Cách tối ưu này giúp mô
hình vừa học được nhãn phân loại cuối cùng, vừa học được các đặc trưng
phụ có ý nghĩa về mặt chất lượng sản phẩm.

\begin{enumerate}
\def\labelenumi{\alph{enumi}.}
\setcounter{enumi}{5}
\item
  \textbf{Phương pháp đánh giá và so sánh mô hình}
\end{enumerate}

Các mô hình được đánh giá thông qua các chỉ số phổ biến trong bài toán
phân loại ảnh, bao gồm Accuracy, Precision, Recall, F1-score, AUC và
Confusion Matrix. Trong đó, Accuracy phản ánh độ chính xác tổng thể;
Precision cho biết mức độ chính xác của các dự đoán thuộc từng lớp;
Recall thể hiện khả năng phát hiện đầy đủ các mẫu thật của từng lớp;
F1-score thể hiện sự cân bằng giữa Precision và Recall; AUC đánh giá khả
năng phân biệt giữa các lớp; và Confusion Matrix giúp phân tích chi tiết
các trường hợp nhầm lẫn.

Luận văn tiến hành so sánh kết quả giữa các mô hình CNN nền và mô hình
đề xuất để xác định mức độ cải thiện. Việc so sánh không chỉ dựa trên độ
chính xác tổng thể mà còn xem xét hiệu quả phân loại từng lớp, đặc biệt
là hai lớp quan trọng trong thực tế gồm tỏi còn sót vỏ và tỏi hư hỏng.

\begin{enumerate}
\def\labelenumi{\alph{enumi}.}
\setcounter{enumi}{6}
\item
  \textbf{Phương pháp phân tích và tổng hợp kết quả}
\end{enumerate}

Sau khi hoàn thành thực nghiệm, luận văn tiến hành phân tích kết quả
huấn luyện và đánh giá mô hình. Nội dung phân tích bao gồm so sánh hiệu
quả giữa các mô hình, đánh giá sự cải thiện của mô hình đề xuất, phân
tích ma trận nhầm lẫn và chỉ ra các trường hợp mô hình còn dự đoán sai.

Thông qua quá trình phân tích, luận văn rút ra nhận xét về ưu điểm, hạn
chế của từng mô hình và mức độ phù hợp của mô hình đề xuất đối với bài
toán phân loại chất lượng tỏi. Trên cơ sở đó, luận văn đề xuất các hướng
phát triển tiếp theo như mở rộng bộ dữ liệu, cải thiện khả năng nhận
diện vùng lỗi nhỏ, tối ưu tốc độ suy luận và tích hợp mô hình vào hệ
thống phân loại tỏi tự động trong dây chuyền chế biến thực phẩm.

\begin{enumerate}
\def\labelenumi{\arabic{enumi}.}
\setcounter{enumi}{4}
\item
  \protect\phantomsection\label{_Toc232111988}{}\textbf{Ý nghĩa nghiên
  cứu}
\end{enumerate}

Đề tài ``Nghiên cứu và phát triển mô hình học sâu cho bài toán phân loại
chất lượng tỏi'' có ý nghĩa cả về mặt khoa học lẫn thực tiễn. Trong bối
cảnh ngành nông nghiệp và chế biến thực phẩm đang từng bước chuyển đổi
theo hướng tự động hóa, việc ứng dụng trí tuệ nhân tạo vào phân loại
chất lượng nông sản là một hướng nghiên cứu cần thiết, có khả năng hỗ
trợ nâng cao năng suất, giảm chi phí lao động và cải thiện chất lượng
sản phẩm đầu ra.

\begin{enumerate}
\def\labelenumi{\arabic{enumi}.}
\item
  \protect\phantomsection\label{_Toc232111989}{}\emph{\textbf{Ý nghĩa
  khoa học}}
\end{enumerate}

Về mặt khoa học, luận văn góp phần nghiên cứu và đánh giá khả năng ứng
dụng các mô hình học sâu trong bài toán phân loại chất lượng tỏi sau bóc
vỏ. Đây là một bài toán có tính đặc thù cao, vì các lớp cần phân loại có
sự khác biệt không quá rõ ràng, đặc biệt là giữa tỏi đã bóc sạch và tỏi
còn sót vỏ. Phần vỏ còn sót có thể xuất hiện ở vùng nhỏ trên bề mặt tép
tỏi, trong khi dấu hiệu hư hỏng có thể biểu hiện dưới nhiều dạng như
biến màu, dập nát hoặc xuất hiện vùng bất thường.

Thông qua việc thực nghiệm các mô hình CNN nền như NASNetMobile,
InceptionV3, EfficientNetB4, MobileNetV2 và DenseNet121, luận văn cung
cấp cơ sở so sánh để đánh giá mức độ phù hợp của từng kiến trúc đối với
bài toán phân loại ảnh tỏi. Kết quả nghiên cứu giúp làm rõ khả năng
trích xuất đặc trưng của các mô hình học sâu trong trường hợp dữ liệu
nông sản có đặc điểm hình ảnh phức tạp và dễ gây nhầm lẫn.

Bên cạnh đó, luận văn đề xuất mô hình cải tiến dựa trên EfficientNetB4
kết hợp GAP/GMP, hai nhánh phụ Spoilage Aux và Peeling Aux, cùng với hàm
mất mát Class-Balanced Loss. Cách tiếp cận này có ý nghĩa khoa học vì mô
hình không chỉ học trực tiếp bài toán phân loại ba lớp, mà còn học thêm
các nhiệm vụ phụ có ý nghĩa theo logic kiểm tra chất lượng tỏi. Cụ thể,
mô hình học cách xác định tỏi có bị hư hay không, sau đó tiếp tục học
cách phân biệt tỏi bóc sạch và tỏi còn vỏ đối với nhóm không hư hỏng.
Điều này giúp tăng khả năng học đặc trưng có định hướng, đặc biệt với
các lớp có sự khác biệt nhỏ.

Ngoài ra, nghiên cứu cũng góp phần bổ sung thêm tài liệu thực nghiệm cho
hướng ứng dụng học sâu trong phân loại chất lượng nông sản sau sơ chế,
một lĩnh vực còn nhiều tiềm năng phát triển tại Việt Nam.

\begin{enumerate}
\def\labelenumi{\arabic{enumi}.}
\setcounter{enumi}{1}
\item
  \protect\phantomsection\label{_Toc232111990}{}\emph{\textbf{Ý nghĩa
  thực tiễn}}
\end{enumerate}

Về mặt thực tiễn, đề tài hướng đến hỗ trợ tự động hóa công đoạn phân
loại tỏi sau bóc vỏ trong dây chuyền chế biến thực phẩm. Hiện nay, công
đoạn này tại nhiều cơ sở sản xuất vẫn phụ thuộc nhiều vào lao động thủ
công. Người lao động phải quan sát bằng mắt thường để phân biệt tỏi
sạch, tỏi còn sót vỏ và tỏi hư. Cách làm này dễ bị ảnh hưởng bởi kinh
nghiệm cá nhân, cường độ làm việc, ánh sáng môi trường và sự mệt mỏi
trong quá trình sản xuất.

Mô hình học sâu được xây dựng trong luận văn có thể hỗ trợ phân loại tỏi
dựa trên ảnh một cách nhanh chóng và nhất quán hơn. Khi được tích hợp
với camera và hệ thống xử lý ảnh, mô hình có thể làm nền tảng cho hệ
thống kiểm tra chất lượng tự động, giúp phát hiện các mẫu tỏi không đạt
yêu cầu trước khi đưa vào công đoạn chế biến tiếp theo.

Việc phân loại chính xác các nhóm tỏi có ý nghĩa trực tiếp đối với chất
lượng sản phẩm. Nếu tỏi hư hoặc tỏi còn sót vỏ bị lẫn vào nhóm tỏi sạch,
sản phẩm cuối cùng có thể bị ảnh hưởng về màu sắc, mùi vị, độ an toàn và
giá trị thương mại. Do đó, hệ thống phân loại tự động có thể góp phần
nâng cao tính đồng nhất của nguyên liệu đầu vào, giảm sai sót trong kiểm
tra chất lượng và hỗ trợ doanh nghiệp cải thiện quy trình sản xuất.

\begin{enumerate}
\def\labelenumi{\arabic{enumi}.}
\setcounter{enumi}{2}
\item
  \protect\phantomsection\label{_Toc232111991}{}\emph{\textbf{Ý nghĩa
  kinh tế và xã hội}}
\end{enumerate}

Việc ứng dụng mô hình học sâu vào phân loại chất lượng tỏi có thể mang
lại lợi ích kinh tế cho các cơ sở chế biến nông sản. Khi công đoạn phân
loại được tự động hóa, doanh nghiệp có thể giảm sự phụ thuộc vào lao
động thủ công, đặc biệt trong các giai đoạn sản xuất có sản lượng lớn.
Điều này góp phần giảm chi phí vận hành, nâng cao năng suất và duy trì
chất lượng phân loại ổn định hơn.

Bên cạnh lợi ích kinh tế, đề tài còn có ý nghĩa xã hội trong việc thúc
đẩy ứng dụng công nghệ vào sản xuất nông nghiệp. Việc sử dụng AI trong
chế biến nông sản giúp nâng cao giá trị sản phẩm, tăng khả năng cạnh
tranh và góp phần hiện đại hóa ngành nông nghiệp Việt Nam. Đây cũng là
hướng đi phù hợp với xu thế phát triển nông nghiệp thông minh, trong đó
các công nghệ như thị giác máy tính, học sâu, cảm biến và tự động hóa
được kết hợp để nâng cao hiệu quả sản xuất.

Ngoài ra, hệ thống phân loại chất lượng tỏi tự động còn có thể hỗ trợ
giảm áp lực cho người lao động trong các công đoạn lặp đi lặp lại, đòi
hỏi quan sát liên tục. Người lao động có thể chuyển sang các công việc
giám sát, vận hành và kiểm soát hệ thống, thay vì thực hiện phân loại
thủ công trong thời gian dài.

\begin{enumerate}
\def\labelenumi{\arabic{enumi}.}
\setcounter{enumi}{3}
\item
  \protect\phantomsection\label{_Toc232111992}{}\emph{\textbf{Ý nghĩa mở
  rộng và định hướng phát triển}}
\end{enumerate}

Kết quả nghiên cứu của luận văn không chỉ có giá trị đối với riêng bài
toán phân loại tỏi mà còn có thể làm cơ sở cho các bài toán phân loại
chất lượng nông sản khác. Các loại nông sản như hành, gừng, khoai tây,
cà rốt hoặc các sản phẩm sau sơ chế cũng có nhu cầu kiểm tra bề mặt,
phát hiện lỗi và phân loại chất lượng tương tự.

Phương pháp đề xuất trong luận văn, đặc biệt là hướng kết hợp backbone
CNN với các nhánh phụ có ý nghĩa theo logic bài toán, có thể được điều
chỉnh để áp dụng cho nhiều bài toán khác trong kiểm tra chất lượng thực
phẩm. Ví dụ, với các sản phẩm khác, mô hình có thể được thiết kế thêm
các nhiệm vụ phụ như phát hiện hư hỏng, phân biệt mức độ chín, nhận diện
phần vỏ còn sót hoặc phát hiện khuyết tật bề mặt.

Trong tương lai, mô hình có thể được phát triển theo hướng tích hợp vào
hệ thống camera, băng chuyền và cơ cấu phân loại tự động. Khi đó, kết
quả nghiên cứu không chỉ dừng lại ở mô hình phân loại ảnh mà có thể trở
thành nền tảng cho một hệ thống kiểm tra chất lượng tỏi hoàn chỉnh, phục
vụ trực tiếp trong dây chuyền chế biến thực phẩm.

Tổng kết, đề tài có ý nghĩa trong việc kết hợp giữa nghiên cứu học sâu
và nhu cầu thực tiễn của ngành chế biến nông sản. Luận văn góp phần
chứng minh tiềm năng của trí tuệ nhân tạo trong tự động hóa phân loại
chất lượng tỏi, đồng thời mở ra hướng phát triển cho các hệ thống kiểm
tra chất lượng nông sản thông minh trong tương lai.

\begin{enumerate}
\def\labelenumi{\arabic{enumi}.}
\setcounter{enumi}{5}
\item
  \protect\phantomsection\label{_Toc232111993}{}\textbf{Cấu trúc luận văn}
\end{enumerate}

Luận văn được cấu trúc thành 5 chương, cụ thể như sau:

\textbf{Chương 1: Giới thiệu tổng quan}

\begin{quote}
Chương 1 trình bày lý do chọn đề tài, tính cấp thiết, tổng quan nghiên
cứu trong và ngoài nước, khoảng trống nghiên cứu, mục tiêu, đối tượng,
phạm vi, phương pháp, ý nghĩa, cấu trúc luận văn và tiến độ thực hiện.
\end{quote}

\textbf{Chương 2: Cơ sở lý thuyết}

\begin{quote}
Chương 2 trình bày tổng quan bài toán phân loại chất lượng tỏi, bộ dữ
liệu Garlic Clove Dataset, cơ sở lý thuyết về học sâu, CNN, học chuyển
giao và các mô hình CNN nền được sử dụng trong nghiên cứu.
\end{quote}

\textbf{Chương 3: Mô hình đề xuất}

\begin{quote}
Chương 3 trình bày quá trình re-train năm mô hình học sâu tiền huấn
luyện, kết quả lựa chọn backbone và mô hình đề xuất
\textbf{CPGC-EfficientNetB4} cho bài toán phân loại chất lượng tỏi.
\end{quote}

\textbf{Chương 4: Kết quả thực nghiệm}

\begin{quote}
Chương 4 trình bày tiêu chí đánh giá, công cụ hỗ trợ, kết quả thực
nghiệm của các mô hình CNN nền và kết quả của mô hình đề xuất.
\end{quote}

\textbf{Chương 5: Kết luận}

\begin{quote}
Chương 5 tổng kết kết quả đạt được, nêu các hạn chế còn tồn tại và đề
xuất hướng phát triển tiếp theo.
\end{quote}

\textbf{Tài liệu tham khảo}

\begin{enumerate}
\def\labelenumi{\arabic{enumi}.}
\setcounter{enumi}{6}
\item
  \protect\phantomsection\label{_Toc232111994}{}\textbf{Tiến độ thực hiện
  luận văn}
\end{enumerate}

{\def\LTcaptype{none} % do not increment counter
\begin{longtable}[]{@{}
  >{\raggedright\arraybackslash}p{(\linewidth - 10\tabcolsep) * \real{0.5888}}
  >{\raggedright\arraybackslash}p{(\linewidth - 10\tabcolsep) * \real{0.0704}}
  >{\raggedright\arraybackslash}p{(\linewidth - 10\tabcolsep) * \real{0.0616}}
  >{\raggedright\arraybackslash}p{(\linewidth - 10\tabcolsep) * \real{0.0529}}
  >{\raggedright\arraybackslash}p{(\linewidth - 10\tabcolsep) * \real{0.0527}}
  >{\raggedright\arraybackslash}p{(\linewidth - 10\tabcolsep) * \real{0.0788}}@{}}
\toprule\noalign{}
\begin{minipage}[b]{\linewidth}\raggedright
Tháng (năm 2025)

Dự kiến nội dung thực hiện
\end{minipage} & \begin{minipage}[b]{\linewidth}\raggedright
5
\end{minipage} & \begin{minipage}[b]{\linewidth}\raggedright
6
\end{minipage} & \begin{minipage}[b]{\linewidth}\raggedright
7
\end{minipage} & \begin{minipage}[b]{\linewidth}\raggedright
8
\end{minipage} & \begin{minipage}[b]{\linewidth}\raggedright
9
\end{minipage} \\
\midrule\noalign{}
\endhead
\bottomrule\noalign{}
\endlastfoot
Tiếp nhận đề tài luận văn & X & & & & \\
Nghiên cứu tổng quan, cơ sở lý thuyết và chuẩn bị dữ liệu & X & X & & & \\
Thực nghiệm các mô hình CNN nền & & X & X & & \\
Xây dựng, tinh chỉnh và đánh giá mô hình đề xuất & & & X & X & \\
Hoàn thiện nội dung luận văn và chuẩn bị bảo vệ & & & & X & X \\
\end{longtable}
}

\section{TÓM TẮT CHƯƠNG 1}\label{tuxf3m-tux1eaft-chux1b0ux1a1ng-1}

Chương 1 đã trình bày những nội dung nền tảng của luận văn
\textbf{``Nghiên cứu và phát triển mô hình học sâu cho bài toán phân
loại chất lượng tỏi''}. Xuất phát từ thực tế trong quá trình chế biến
tỏi, mặc dù các công đoạn như tách tép và bóc vỏ đã được hỗ trợ bằng máy
móc, nhưng công đoạn phân loại tỏi sau bóc vỏ vẫn còn phụ thuộc nhiều
vào lao động thủ công. Việc nhận biết tỏi đã bóc sạch, tỏi còn sót vỏ và
tỏi hư hỏng bằng mắt thường dễ bị ảnh hưởng bởi kinh nghiệm, điều kiện
làm việc và sự mệt mỏi của người lao động, từ đó có thể làm giảm tính ổn
định của chất lượng sản phẩm.

Trên cơ sở đó, chương này đã làm rõ tính cấp thiết của việc ứng dụng trí
tuệ nhân tạo, đặc biệt là học sâu và thị giác máy tính, vào bài toán
phân loại chất lượng tỏi. Các nghiên cứu trong và ngoài nước cho thấy
học sâu đã được ứng dụng hiệu quả trong nhiều bài toán kiểm tra chất
lượng nông sản, phát hiện lỗi bề mặt và hỗ trợ tự động hóa dây chuyền
sản xuất. Tuy nhiên, các nghiên cứu trực tiếp về phân loại tỏi sau bóc
vỏ thành ba nhóm \textbf{tỏi đã bóc sạch}, \textbf{tỏi còn sót vỏ} và
\textbf{tỏi hư hỏng} vẫn còn hạn chế. Đây là khoảng trống để luận văn
tiếp tục nghiên cứu và phát triển.

Chương 1 cũng đã xác định mục tiêu, đối tượng, phạm vi và phương pháp
nghiên cứu của luận văn. Mục tiêu chính là xây dựng mô hình học sâu có
khả năng phân loại chất lượng tỏi dựa trên ảnh. Để thực hiện mục tiêu
này, luận văn tập trung vào việc thu thập và xử lý bộ dữ liệu ảnh tỏi,
thực nghiệm các mô hình CNN nền như \textbf{NASNetMobile},
\textbf{InceptionV3}, \textbf{EfficientNetB4}, \textbf{MobileNetV2} và
\textbf{DenseNet121}, từ đó lựa chọn mô hình phù hợp làm nền tảng cho
hướng cải tiến. Trên cơ sở đó, luận văn định hướng đề xuất mô hình dựa
trên \textbf{EfficientNetB4}, kết hợp \textbf{GAP/GMP}, hai nhánh phụ
\textbf{Spoilage Aux}, \textbf{Peeling Aux} và hàm mất mát
\textbf{CBLoss}.

Ngoài ra, chương này đã nêu ý nghĩa khoa học và thực tiễn của đề tài. Về
mặt khoa học, nghiên cứu góp phần làm rõ khả năng ứng dụng học sâu trong
phân loại chất lượng nông sản sau sơ chế. Về mặt thực tiễn, kết quả
nghiên cứu có thể làm cơ sở cho việc phát triển hệ thống phân loại tỏi
tự động trong dây chuyền chế biến thực phẩm, góp phần giảm phụ thuộc vào
lao động thủ công, nâng cao năng suất và cải thiện chất lượng sản phẩm.
Những nội dung trong chương 1 là cơ sở để các chương tiếp theo trình bày
sâu hơn về cơ sở lý thuyết, phương pháp đề xuất, thực nghiệm và đánh giá
kết quả.

\section{\texorpdfstring{CHƯƠNG 2: CƠ SỞ LÝ THUYẾT VÀ MÔ HÌNH NGHIÊN CỨU
}{CHƯƠNG 2: CƠ SỞ LÝ THUYẾT VÀ MÔ HÌNH NGHIÊN CỨU }}\label{chux1b0ux1a1ng-2-cux1a1-sux1edf-luxfd-thuyux1ebft-vuxe0-muxf4-huxecnh-nghiuxean-cux1ee9u}

\begin{enumerate}
\def\labelenumi{\arabic{enumi}.}
\item
  \protect\phantomsection\label{_Toc232111997}{}\textbf{Tổng quan về bài
  toán phân loại chất lượng tỏi}
\end{enumerate}

Trong quy trình chế biến tỏi, sau khi tỏi được tách thành từng tép và
đưa qua công đoạn bóc vỏ, sản phẩm đầu ra thường chưa hoàn toàn đồng
nhất. Một số tép tỏi đã được bóc sạch, một số vẫn còn sót lại lớp vỏ
mỏng, trong khi một số khác có thể bị hư hỏng, dập nát hoặc biến màu do
ảnh hưởng từ quá trình thu hoạch, bảo quản và xử lý. Vì vậy, công đoạn
phân loại chất lượng tỏi sau bóc vỏ là bước cần thiết trước khi nguyên
liệu được đưa vào các công đoạn chế biến tiếp theo.

Bài toán phân loại chất lượng tỏi trong luận văn này được hiểu là quá
trình xác định trạng thái chất lượng của từng mẫu tỏi sau bóc vỏ. Mỗi
mẫu tỏi cần được phân vào một trong ba nhóm chính: \textbf{tỏi đã bóc
sạch}, \textbf{tỏi còn sót vỏ} và \textbf{tỏi hư hỏng}. Đây là ba trạng
thái thường gặp trong thực tế sản xuất và có ảnh hưởng trực tiếp đến
chất lượng nguyên liệu đầu vào.

Trong đó, \textbf{tỏi đã bóc sạch} là nhóm tỏi đạt yêu cầu, bề mặt tương
đối sạch, không còn lớp vỏ rõ rệt và không có dấu hiệu hư hỏng đáng kể.
\textbf{Tỏi còn sót vỏ} là nhóm tỏi chưa được bóc vỏ hoàn chỉnh, trên bề
mặt vẫn còn một phần vỏ bám lại. Nhóm này cần được phát hiện để có thể
xử lý lại hoặc tách riêng khỏi nhóm đạt yêu cầu. \textbf{Tỏi hư hỏng} là
nhóm tỏi có các biểu hiện không đạt chất lượng như dập nát, thâm đen,
biến màu, mềm nhũn hoặc xuất hiện vùng hư trên bề mặt. Đây là nhóm cần
được loại bỏ để tránh ảnh hưởng đến chất lượng sản phẩm cuối cùng.

Việc phân loại chính xác ba nhóm tỏi trên có ý nghĩa quan trọng trong
dây chuyền chế biến thực phẩm. Nếu tỏi còn sót vỏ bị đưa vào nhóm tỏi
sạch, sản phẩm sau chế biến có thể không đạt yêu cầu về cảm quan và chất
lượng. Nếu tỏi hư bị lẫn vào nguyên liệu đạt chuẩn, chất lượng, mùi vị
và độ an toàn của sản phẩm có thể bị ảnh hưởng. Do đó, công đoạn phân
loại không chỉ nhằm tách sản phẩm theo hình thức bên ngoài mà còn góp
phần kiểm soát chất lượng nguyên liệu trước khi chế biến.

Tuy nhiên, bài toán phân loại chất lượng tỏi có một số khó khăn nhất
định. Tép tỏi thường có kích thước nhỏ, hình dạng không đồng đều và bề
mặt có nhiều biến thiên tự nhiên. Sự khác biệt giữa tỏi đã bóc sạch và
tỏi còn sót vỏ đôi khi không quá rõ ràng, đặc biệt khi phần vỏ còn lại
chỉ chiếm một vùng nhỏ. Bên cạnh đó, dấu hiệu hư hỏng của tỏi cũng có
thể xuất hiện dưới nhiều dạng khác nhau, từ các vết thâm nhỏ đến vùng
dập nát lớn. Những yếu tố này làm cho quá trình phân loại bằng mắt
thường dễ bị nhầm lẫn, nhất là khi sản lượng lớn hoặc tốc độ dây chuyền
cao.

Từ những đặc điểm trên, có thể thấy bài toán phân loại chất lượng tỏi là
một bài toán có ý nghĩa thực tiễn trong chế biến nông sản. Việc nghiên
cứu phương pháp phân loại tự động dựa trên hình ảnh là hướng tiếp cận
phù hợp nhằm hỗ trợ giảm phụ thuộc vào lao động thủ công, nâng cao tính
ổn định của quá trình kiểm tra chất lượng và tạo tiền đề cho việc xây
dựng hệ thống phân loại tỏi tự động trong tương lai.

\begin{enumerate}
\def\labelenumi{\arabic{enumi}.}
\setcounter{enumi}{1}
\item
  \protect\phantomsection\label{_Toc232111998}{}\textbf{Tổng quan về bộ
  dữ liệu ảnh tỏi sử dụng trong nghiên cứu}
\end{enumerate}

Bộ dữ liệu là một trong những thành phần quan trọng quyết định hiệu quả
của mô hình học sâu. Đối với bài toán phân loại chất lượng tỏi, dữ liệu
ảnh cần phản ánh được các trạng thái thực tế của tỏi sau quá trình bóc
vỏ, bao gồm tỏi đã bóc sạch, tỏi còn sót vỏ và tỏi bị hư hỏng. Chất
lượng ảnh, cách tổ chức dữ liệu, độ chính xác của nhãn và mức độ đa dạng
giữa các mẫu đều có ảnh hưởng trực tiếp đến khả năng học đặc trưng cũng
như khả năng khái quát hóa của mô hình.

Trong nghiên cứu này, bộ dữ liệu chính được sử dụng là \textbf{Garlic
Clove Dataset}, bao gồm các ảnh tép tỏi sau quá trình bóc vỏ. Dữ liệu
được thu thập vào tháng 02 năm 2025 tại một cơ sở chế biến tỏi thương
mại ở Thành phố Hồ Chí Minh, Việt Nam. Các ảnh được chụp bằng camera
điện thoại thông minh \textbf{iPhone 14} trong môi trường xử lý tỏi thực
tế. Việc thu thập dữ liệu trong điều kiện sản xuất thực tế giúp bộ dữ
liệu phản ánh tốt hơn các tình huống có thể gặp trong quá trình kiểm tra
chất lượng tỏi sau bóc vỏ.

Các ảnh trong bộ dữ liệu được cắt thủ công thành các vùng ảnh vuông,
trong đó mỗi ảnh tương ứng với một mẫu tép tỏi riêng biệt. Quá trình gán
nhãn được thực hiện bằng công cụ \textbf{LabelImg} dựa trên các tiêu chí
chất lượng đã được xác định trước. Sau khi gán nhãn, dữ liệu được kiểm
tra lại bằng quan sát trực quan nhằm bảo đảm tính nhất quán của nhãn
trong quá trình chuẩn bị dữ liệu.

Trong luận văn, hai bộ dữ liệu được ký hiệu lần lượt là
\textbf{dataset-1} và \textbf{dataset-2} để thuận tiện cho quá trình
trình bày và so sánh kết quả thực nghiệm. Trong đó, \textbf{dataset-1}
là bộ dữ liệu ban đầu gồm \textbf{2.134 ảnh}, được thu thập, cắt vùng
ảnh và gán nhãn theo ba nhóm chất lượng tỏi. \textbf{dataset-2} là bộ dữ
liệu mở rộng gồm \textbf{2.944 ảnh}, trong đó số lượng ảnh của lớp
\textbf{Bóc sạch vỏ} được bổ sung nhằm tăng quy mô dữ liệu và phục vụ
đánh giá mô hình trong điều kiện phân bố lớp khác nhau.

Trong phạm vi nghiên cứu, ảnh tỏi được chia thành ba nhóm chất lượng
chính: \textbf{Tỏi sạch vỏ}, \textbf{Tỏi còn vỏ} và \textbf{Tỏi hư}. Ba
nhóm này tương ứng với ba lớp dữ liệu tiếng Anh là \textbf{Fully Peeled
Garlic}, \textbf{Partially Peeled Garlic} và \textbf{Spoiled Garlic}.
Đây là ba trạng thái thường gặp trong dây chuyền chế biến tỏi sau bóc vỏ
và có ý nghĩa trực tiếp đối với việc kiểm tra chất lượng nguyên liệu.

Bảng 2.1. Định nghĩa các lớp trong bộ dữ liệu ảnh tỏi

{\def\LTcaptype{none} % do not increment counter
\begin{longtable}[]{@{}
  >{\centering\arraybackslash}p{(\linewidth - 6\tabcolsep) * \real{0.0548}}
  >{\raggedright\arraybackslash}p{(\linewidth - 6\tabcolsep) * \real{0.2075}}
  >{\raggedright\arraybackslash}p{(\linewidth - 6\tabcolsep) * \real{0.4978}}
  >{\raggedright\arraybackslash}p{(\linewidth - 6\tabcolsep) * \real{0.1935}}@{}}
\toprule\noalign{}
\begin{minipage}[b]{\linewidth}\centering
\textbf{STT}
\end{minipage} & \begin{minipage}[b]{\linewidth}\centering
\textbf{Tên lớp}
\end{minipage} & \begin{minipage}[b]{\linewidth}\centering
\textbf{Mô tả}
\end{minipage} & \begin{minipage}[b]{\linewidth}\centering
\textbf{Trực quan hóa}
\end{minipage} \\
\midrule\noalign{}
\endhead
\bottomrule\noalign{}
\endlastfoot
1 & Tỏi sạch vỏ

(\emph{Fully Peeled Garlic}) & Các tép tỏi đã được bóc sạch, bề mặt
tương đối nhẵn, đồng đều, không còn vỏ rõ rệt và không có khuyết tật dễ
quan sát. Đây là nhóm tỏi đạt yêu cầu để đóng gói, chế biến công nghiệp
hoặc sử dụng trực tiếp. &
\includegraphics[width=1.29638in,height=1.1811in]{./media/image3.jpeg} \\
2 & Tỏi còn vỏ

\emph{(Partially Peeled Garlic)} & Các tép tỏi vẫn còn sót lại một phần
vỏ bên ngoài trên bề mặt. Những mẫu này thường cần được đưa qua bước làm
sạch hoặc bóc vỏ bổ sung trước khi được xem là tỏi đạt yêu cầu. &
\includegraphics[width=1.27529in,height=1.1811in]{./media/image4.jpeg} \\
3 & Tỏi hư

\emph{(Spoiled Garlic)} & Các tép tỏi có dấu hiệu không đạt chất lượng
như đổi màu, mốc, mềm, dập nát hoặc hư hại bề mặt. Những mẫu này cần
được loại bỏ khỏi dây chuyền do có thể ảnh hưởng đến chất lượng và an
toàn thực phẩm. &
\includegraphics[width=1.28302in,height=1.18056in]{./media/image5.jpeg} \\
\end{longtable}
}

Bộ dữ liệu \textbf{dataset-1} gồm tổng cộng \textbf{2.134 ảnh}, được
phân bố vào ba lớp chất lượng. Trong đó, lớp \textbf{Bóc sạch vỏ} có
\textbf{690 ảnh}, lớp \textbf{Bóc chưa sạch} có \textbf{438 ảnh} và lớp
\textbf{Hư hỏng} có \textbf{1.006 ảnh}. Bộ dữ liệu được chia thành ba
tập con gồm tập huấn luyện, tập kiểm định và tập kiểm tra. Việc chia dữ
liệu được thực hiện ở giai đoạn thu nhận ảnh nhằm hạn chế tình trạng ảnh
trùng lặp hoặc gần trùng lặp xuất hiện ở các tập con khác nhau.

Bộ dữ liệu \textbf{dataset-2} gồm tổng cộng \textbf{2.944 ảnh}. So với
dataset-1, dataset-2 được mở rộng ở lớp \textbf{Bóc sạch vỏ}, nâng số
lượng ảnh của lớp này lên \textbf{1.500 ảnh}. Hai lớp \textbf{Bóc chưa
sạch} và \textbf{Hư hỏng} lần lượt giữ số lượng là \textbf{438 ảnh} và
\textbf{1.006 ảnh}. Việc sử dụng thêm dataset-2 giúp nghiên cứu có thêm
cơ sở đánh giá mô hình trên quy mô dữ liệu lớn hơn, đồng thời xem xét
ảnh hưởng của sự thay đổi phân bố dữ liệu giữa các lớp đến hiệu quả phân
loại.

Bảng 2.2 trình bày chi tiết số lượng ảnh theo từng tập con và từng lớp.
Thông tin này giúp làm rõ quy mô dữ liệu được sử dụng trong các thí
nghiệm, đồng thời cho thấy lớp \textbf{Tỏi còn vỏ} có số lượng mẫu ít
hơn hai lớp còn lại, vì vậy cần chú ý đến vấn đề mất cân bằng dữ liệu
trong quá trình huấn luyện và đánh giá mô hình.

Bảng 2.2. Phân bố ảnh theo lớp và tập con trong dataset-1 và dataset-2

{\def\LTcaptype{none} % do not increment counter
\begin{longtable}[]{@{}
  >{\raggedright\arraybackslash}p{(\linewidth - 10\tabcolsep) * \real{0.1400}}
  >{\raggedright\arraybackslash}p{(\linewidth - 10\tabcolsep) * \real{0.1500}}
  >{\raggedright\arraybackslash}p{(\linewidth - 10\tabcolsep) * \real{0.1800}}
  >{\raggedright\arraybackslash}p{(\linewidth - 10\tabcolsep) * \real{0.1800}}
  >{\raggedright\arraybackslash}p{(\linewidth - 10\tabcolsep) * \real{0.1700}}
  >{\raggedright\arraybackslash}p{(\linewidth - 10\tabcolsep) * \real{0.1000}}@{}}
\toprule\noalign{}
\textbf{Bộ dữ liệu} & \textbf{Tập con} & \textbf{Tỏi sạch vỏ} &
\textbf{Tỏi còn vỏ} & \textbf{Tỏi hư} & \textbf{Tổng} \\
\midrule\noalign{}
\endhead
\bottomrule\noalign{}
\endlastfoot
dataset-1 & Train & 482 & 306 & 704 & 1.492 \\
dataset-1 & Validation & 103 & 65 & 150 & 318 \\
dataset-1 & Test & 105 & 67 & 152 & 324 \\
dataset-1 & Tổng & 690 & 438 & 1.006 & 2.134 \\
dataset-2 & Train & 1.050 & 306 & 704 & 2.060 \\
dataset-2 & Validation & 225 & 65 & 150 & 440 \\
dataset-2 & Test & 225 & 67 & 152 & 444 \\
dataset-2 & Tổng & 1.500 & 438 & 1.006 & 2.944 \\
\end{longtable}
}

Khi sử dụng dataset-2, việc bổ sung ảnh cần được thực hiện theo nguyên
tắc không để các ảnh trùng lặp hoặc gần trùng lặp xuất hiện đồng thời ở
các tập train, validation và test. Nếu ảnh bổ sung được tạo ra từ quá
trình tăng cường dữ liệu, các biến thể của cùng một ảnh gốc phải nằm
trong cùng một tập con để tránh rò rỉ dữ liệu. Nguyên tắc này giúp kết
quả đánh giá phản ánh đúng khả năng khái quát hóa của mô hình trên các
mẫu tỏi chưa từng thấy trong quá trình huấn luyện.

\includegraphics[width=6.52974in,height=2.68595in]{./media/image6.png}

\textbf{Hình 2.1. Phân bố số lượng ảnh và tỷ lệ lớp trong dataset-1 và
dataset-2}

Hình 2.1 minh họa trực quan sự khác biệt về phân bố dữ liệu giữa
\textbf{dataset-1} và \textbf{dataset-2}. Ở Hình 2.1(a), có thể thấy số
lượng ảnh của lớp \textbf{Tỏi sạch vỏ} trong dataset-2 tăng đáng kể so
với dataset-1, từ \textbf{690 ảnh} lên \textbf{1.500 ảnh}. Trong khi đó,
số lượng ảnh của hai lớp \textbf{Tỏi còn vỏ} và \textbf{Tỏi hư} được giữ
nguyên ở cả hai bộ dữ liệu, lần lượt là \textbf{438 ảnh} và
\textbf{1.006 ảnh}.

Ở Hình 2.1(b), sự thay đổi về tỷ lệ lớp được thể hiện rõ hơn. Trong
dataset-1, lớp \textbf{Tỏi hư} chiếm tỷ lệ cao nhất với khoảng
\textbf{47,1\%}, trong khi lớp \textbf{Tỏi sạch vỏ} chiếm khoảng
\textbf{32,3\%} và lớp \textbf{Tỏi còn vỏ} chiếm khoảng \textbf{20,5\%}.
Ngược lại, trong dataset-2, lớp \textbf{Tỏi sạch vỏ} trở thành lớp chiếm
tỷ lệ lớn nhất với khoảng \textbf{51,0\%}, lớp \textbf{Tỏi hư} giảm còn
khoảng \textbf{34,2\%}, và lớp \textbf{Tỏi còn vỏ} chiếm khoảng
\textbf{14,9\%}.

Sự khác biệt này cho thấy \textbf{dataset-2} không chỉ mở rộng quy mô dữ
liệu mà còn làm thay đổi phân bố giữa các lớp. Điều này có ý nghĩa quan
trọng trong quá trình đánh giá mô hình, vì mô hình cần được kiểm tra
trong các điều kiện dữ liệu khác nhau để xem xét khả năng học, khả năng
khái quát hóa và mức độ ổn định khi phân bố lớp thay đổi.

\begin{enumerate}
\def\labelenumi{\arabic{enumi}.}
\setcounter{enumi}{2}
\item
  \protect\phantomsection\label{_Toc232111999}{}\textbf{Cơ sở lý thuyết
  về học sâu}
\end{enumerate}

Học sâu, hay \textbf{Deep Learning}, là một nhánh của học máy, trong đó
mô hình được xây dựng từ nhiều tầng xử lý liên tiếp nhằm tự động học
biểu diễn đặc trưng từ dữ liệu. Thay vì phải thiết kế đặc trưng thủ công
như một số phương pháp học máy truyền thống, mô hình học sâu có khả năng
tự rút trích các đặc điểm quan trọng thông qua quá trình huấn luyện.
Theo LeCun, Bengio và Hinton, học sâu cho phép các mô hình tính toán gồm
nhiều tầng xử lý học được biểu diễn dữ liệu ở nhiều mức độ trừu tượng
khác nhau, qua đó đạt hiệu quả cao trong các bài toán như nhận dạng ảnh,
nhận dạng giọng nói và phát hiện đối tượng.

Trong bài toán phân loại chất lượng tỏi, dữ liệu đầu vào là ảnh các tép
tỏi sau bóc vỏ. Mỗi ảnh có thể chứa nhiều thông tin trực quan như màu
sắc, hình dạng, bề mặt, phần vỏ còn sót hoặc dấu hiệu hư hỏng. Những đặc
điểm này không phải lúc nào cũng dễ mô tả bằng các quy tắc cố định. Vì
vậy, học sâu là hướng tiếp cận phù hợp vì mô hình có thể tự học các đặc
trưng cần thiết từ dữ liệu ảnh để phục vụ quá trình phân loại.

\begin{enumerate}
\def\labelenumi{\arabic{enumi}.}
\item
  \protect\phantomsection\label{_Toc232112000}{}\emph{\textbf{Khái niệm
  học sâu}}
\end{enumerate}

Học sâu có thể được hiểu là phương pháp học biểu diễn dữ liệu thông qua
nhiều tầng tính toán. Mỗi tầng trong mô hình nhận dữ liệu từ tầng trước,
thực hiện một phép biến đổi và truyền kết quả sang tầng tiếp theo. Qua
nhiều tầng liên tiếp, mô hình có thể học được các đặc trưng từ đơn giản
đến phức tạp.

Đối với dữ liệu ảnh, các tầng đầu của mô hình thường học những đặc điểm
cơ bản như màu sắc, vùng sáng tối hoặc đường biên. Các tầng sâu hơn có
thể học các đặc trưng có ý nghĩa cao hơn, chẳng hạn kết cấu bề mặt, hình
dạng đối tượng hoặc những vùng bất thường trên ảnh. Nhờ đó, mô hình học
sâu có thể phân biệt các nhóm ảnh khác nhau dựa trên đặc trưng đã học
được từ dữ liệu huấn luyện.

Trong nghiên cứu này, học sâu được sử dụng để hỗ trợ phân loại ảnh tỏi
thành ba nhóm: \textbf{tỏi sạch vỏ}, \textbf{tỏi còn vỏ} và \textbf{tỏi
hư}. Đây là bài toán mà sự khác biệt giữa các lớp đôi khi không rõ ràng,
chẳng hạn phần vỏ còn sót có thể rất mỏng hoặc dấu hiệu hư hỏng chỉ xuất
hiện ở một vùng nhỏ. Do đó, việc để mô hình tự học đặc trưng từ ảnh là
cần thiết.

\begin{enumerate}
\def\labelenumi{\arabic{enumi}.}
\setcounter{enumi}{1}
\item
  \protect\phantomsection\label{_Toc232112001}{}\emph{\textbf{Cấu trúc
  tổng quát của mạng nơ-ron}}
\end{enumerate}

Mô hình học sâu thường được xây dựng dựa trên \textbf{mạng nơ-ron nhân
tạo}. Một mạng nơ-ron cơ bản gồm ba thành phần chính: lớp đầu vào, các
lớp ẩn và lớp đầu ra. Lớp đầu vào tiếp nhận dữ liệu ban đầu. Trong bài
toán ảnh, dữ liệu đầu vào thường là ma trận các giá trị điểm ảnh. Các
lớp ẩn thực hiện quá trình biến đổi và học đặc trưng từ dữ liệu. Lớp đầu
ra tạo ra kết quả dự đoán cuối cùng, ví dụ xác suất ảnh thuộc về từng
lớp chất lượng tỏi.

Mạng nơ-ron là mô hình học máy gồm các ``nơ-ron'' đơn giản được sắp xếp
theo từng lớp, trong đó mô hình học các trọng số và hệ số lệch từ dữ
liệu để ánh xạ đầu vào thành đầu ra. Mỗi liên kết giữa các nơ-ron được
nhân với một trọng số; đầu vào của hàm kích hoạt có thể được hiểu là
tổng có trọng số của các đầu ra từ lớp trước, thường cộng thêm một hệ số
lệch {[}33{]}.

Công thức tổng quát của một nơ-ron nhân tạo có thể biểu diễn như sau:

\[
y = f\left(\sum_{i=1}^{n} w_i x_i + b\right)
\]

\begin{quote}
Trong đó:
\end{quote}

\begin{itemize}
\item
  \(x_{i}\ \)là giá trị đầu vào thứ \(i\);
\item
  \(w_{i}\ \)\hspace{0pt} là trọng số tương ứng với đầu vào
  \(x_{i}\)\hspace{0pt};
\item
  \(b\) là hệ số lệch;
\item
  \(f\) là hàm kích hoạt;
\item
  \(y\) là đầu ra của nơ-ron.
\end{itemize}

Công thức trên cho thấy mỗi nơ-ron không chỉ nhận dữ liệu đầu vào mà còn
học mức độ quan trọng của từng đầu vào thông qua các trọng số. Trong quá
trình huấn luyện, các trọng số và hệ số lệch được điều chỉnh dần để mô
hình đưa ra dự đoán chính xác hơn.

\includegraphics[width=6.30487in,height=3.56976in]{./media/image7.png}

\textbf{Hình 2.2. Cấu trúc tổng quát của mạng nơ-ron nhân tạo}

\begin{enumerate}
\def\labelenumi{\arabic{enumi}.}
\item
  \protect\phantomsection\label{_Toc232112002}{}\emph{\textbf{Quá trình
  học của mô hình}}
\end{enumerate}

Quá trình học của mô hình học sâu thường gồm hai bước chính: \textbf{lan
truyền tiến} và \textbf{lan truyền ngược}. Ở bước lan truyền tiến, dữ
liệu đầu vào được đưa qua các tầng của mạng để tạo ra kết quả dự đoán.
Kết quả này sau đó được so sánh với nhãn thực tế thông qua một hàm mất
mát nhằm đo mức độ sai lệch giữa dự đoán của mô hình và đáp án đúng.

Sau khi tính được sai số, mô hình thực hiện lan truyền ngược để cập nhật
các tham số. Quá trình lan truyền ngược sử dụng thông tin từ sai số để
xác định mức độ ảnh hưởng của từng trọng số đến kết quả dự đoán, từ đó
điều chỉnh các trọng số theo hướng làm giảm giá trị mất mát. Nhìn chung,
quá trình huấn luyện mạng nơ-ron gồm các bước như lan truyền tiến, tính
sai số, lan truyền ngược và cập nhật trọng số bằng phương pháp tối ưu
như gradient descent {[}33{]}.

Có thể hiểu đơn giản rằng, trong quá trình huấn luyện, mô hình liên tục
thực hiện dự đoán, so sánh với nhãn đúng, sau đó điều chỉnh lại các tham
số để cải thiện kết quả ở những lần học tiếp theo. Quá trình này được
lặp lại nhiều lần trên tập dữ liệu huấn luyện cho đến khi mô hình đạt
được hiệu quả mong muốn hoặc không còn cải thiện đáng kể.

\begin{enumerate}
\def\labelenumi{\arabic{enumi}.}
\setcounter{enumi}{1}
\item
  \protect\phantomsection\label{_Toc232112003}{}\emph{\textbf{Vai trò
  của học sâu trong phân loại ảnh}}
\end{enumerate}

Học sâu có vai trò quan trọng trong các bài toán phân loại ảnh vì mô
hình có khả năng tự động học đặc trưng từ dữ liệu đầu vào. Thay vì phải
xây dựng thủ công các đặc trưng như màu sắc, hình dạng hoặc kết cấu, mô
hình có thể tự tìm ra những đặc điểm phù hợp nhất cho nhiệm vụ phân
loại.

Đối với bài toán phân loại chất lượng tỏi, các dấu hiệu cần nhận diện có
thể rất đa dạng. Tỏi còn vỏ có thể chỉ còn một phần vỏ nhỏ bám trên bề
mặt, trong khi tỏi hư có thể biểu hiện qua vết thâm, vùng mềm, đổi màu
hoặc hư hại bề mặt. Những dấu hiệu này khó có thể được mô tả đầy đủ bằng
một tập quy tắc cố định. Do đó, học sâu giúp mô hình linh hoạt hơn trong
việc nhận diện sự khác biệt giữa các nhóm chất lượng tỏi.

Tóm lại, học sâu cung cấp nền tảng quan trọng cho việc xây dựng mô hình
phân loại chất lượng tỏi trong luận văn. Các nội dung về mạng nơ-ron,
quá trình học và khả năng tự động trích xuất đặc trưng là cơ sở để tiếp
tục trình bày các mô hình chuyên biệt hơn cho dữ liệu ảnh ở những phần
tiếp theo.

\begin{enumerate}
\def\labelenumi{\arabic{enumi}.}
\setcounter{enumi}{3}
\item
  \protect\phantomsection\label{_Toc232112004}{}\textbf{Mạng nơ-ron tích
  chập trong phân loại ảnh}
\end{enumerate}

Mạng nơ-ron tích chập, hay \textbf{Convolutional Neural Network -- CNN},
là một trong những kiến trúc học sâu được sử dụng phổ biến trong các bài
toán xử lý ảnh. CNN được thiết kế phù hợp với dữ liệu có cấu trúc dạng
lưới, đặc biệt là ảnh số, nhờ khả năng khai thác mối quan hệ không gian
giữa các điểm ảnh. CNN thường được sử dụng cho các nhiệm vụ như phân
loại ảnh và nhận diện đối tượng, đồng thời giúp giảm sự phụ thuộc vào
quá trình trích xuất đặc trưng thủ công trong xử lý ảnh {[}34{]}.

Trong bài toán phân loại ảnh, CNN nhận ảnh đầu vào và tự động trích xuất
đặc trưng thông qua nhiều lớp xử lý liên tiếp. Các lớp đầu thường học
những đặc trưng đơn giản như cạnh, đường viền, vùng sáng tối hoặc màu
sắc. Các lớp sâu hơn có thể học những đặc trưng phức tạp hơn như kết cấu
bề mặt, hình dạng đối tượng hoặc các vùng bất thường trên ảnh. Đối với
bài toán phân loại chất lượng tỏi, khả năng này rất quan trọng vì sự
khác biệt giữa các lớp có thể nằm ở những chi tiết nhỏ như phần vỏ còn
sót, vết thâm hoặc vùng hư hỏng trên bề mặt tép tỏi.

\includegraphics[width=5.99174in,height=2.97149in]{./media/image8.png}

\textbf{Hình 2.3.} \emph{Cấu trúc tổng quát của mạng nơ-ron tích chập}

\begin{enumerate}
\def\labelenumi{\arabic{enumi}.}
\item
  \protect\phantomsection\label{_Toc232112005}{}\emph{\textbf{Khái niệm
  mạng nơ-ron tích chập}}
\end{enumerate}

CNN là một dạng mạng nơ-ron nhân tạo được xây dựng chuyên biệt cho dữ
liệu ảnh. Nếu mạng nơ-ron truyền thống thường xử lý dữ liệu dưới dạng
vector một chiều, thì CNN khai thác trực tiếp cấu trúc không gian của
ảnh thông qua chiều rộng, chiều cao và chiều sâu kênh màu. Stanford
CS231n mô tả CNN là mạng gồm các nơ-ron có trọng số và hệ số lệch học
được, nhưng khác với mạng nơ-ron thông thường ở chỗ kiến trúc CNN giả
định đầu vào là ảnh, từ đó tổ chức các lớp theo không gian ba chiều và
giúp giảm số lượng tham số cần học {[}35{]}.

Điểm khác biệt quan trọng của CNN là mỗi nơ-ron trong lớp tích chập
không cần kết nối với toàn bộ ảnh đầu vào, mà chỉ kết nối với một vùng
cục bộ. Nhờ vậy, CNN có thể tập trung phát hiện các đặc trưng nhỏ trong
ảnh, đồng thời giảm đáng kể số lượng tham số so với mạng kết nối đầy đủ.
Đặc điểm này phù hợp với bài toán phân loại tỏi, vì các dấu hiệu cần
nhận diện như vùng vỏ còn sót hoặc vết hư hỏng thường chỉ xuất hiện ở
một phần nhỏ trên ảnh.

Một kiến trúc CNN cơ bản thường bao gồm các thành phần chính như lớp
tích chập, hàm kích hoạt, lớp pooling và lớp phân loại. Các thành phần
này phối hợp với nhau để biến đổi ảnh đầu vào thành các đặc trưng có ý
nghĩa, sau đó sử dụng các đặc trưng đó để đưa ra nhãn phân loại cuối
cùng.

\begin{enumerate}
\def\labelenumi{\arabic{enumi}.}
\setcounter{enumi}{1}
\item
  \protect\phantomsection\label{_Toc232112006}{}\emph{\textbf{Lớp tích
  chập}}
\end{enumerate}

Lớp tích chập là thành phần cốt lõi của CNN. Lớp này sử dụng các bộ lọc,
còn gọi là \textbf{kernel} hoặc \textbf{filter}, để trượt trên ảnh đầu
vào và tạo ra các bản đồ đặc trưng. Mỗi bộ lọc có nhiệm vụ phát hiện một
loại đặc trưng nhất định, chẳng hạn cạnh, vùng sáng tối, hoa văn, kết
cấu hoặc một dạng mẫu cục bộ nào đó. Theo Stanford CS231n, trong quá
trình lan truyền tiến, mỗi bộ lọc được trượt trên chiều rộng và chiều
cao của ảnh đầu vào, sau đó tính tích vô hướng giữa các giá trị trong bộ
lọc và vùng ảnh tương ứng để tạo ra bản đồ kích hoạt {[}35{]}.

Về mặt khái quát, phép tích chập hai chiều có thể được biểu diễn như
sau:

\[
S(i,j)=\sum_m \sum_n X(i+m,j+n)K(m,n)
\]

Trong đó:

\begin{itemize}
\item
  \(X\) là ảnh hoặc bản đồ đặc trưng đầu vào.
\item
  \(K\) là bộ lọc tích chập.
\item
  còn \(S(i,j)\) là giá trị đầu ra tại vị trí\(\ (i,j)\).
\end{itemize}

Công thức này thể hiện việc bộ lọc quét qua ảnh và tính tổng có trọng số
tại từng vùng cục bộ. Tài liệu \textbf{``\emph{Dive into Deep
Learning''}} cũng trình bày lớp tích chập là nền tảng của CNN trong việc
phát hiện mẫu cục bộ và tạo bản đồ đặc trưng cho các tầng tiếp theo
{[}36{]}.

Trong bài toán phân loại chất lượng tỏi, lớp tích chập giúp mô hình phát
hiện các đặc trưng trực quan trên bề mặt tép tỏi. Ở các tầng đầu, mô
hình có thể học các đặc trưng đơn giản như đường biên, vùng sáng tối
hoặc sự thay đổi màu sắc. Ở các tầng sâu hơn, mô hình có thể học các dấu
hiệu phức tạp hơn như vùng vỏ còn sót, vết thâm, dập nát hoặc những bất
thường trên bề mặt tỏi.

\begin{enumerate}
\def\labelenumi{\arabic{enumi}.}
\item
  \protect\phantomsection\label{_Toc232112007}{}\emph{\textbf{Hàm kích
  hoạt trong CNN}}
\end{enumerate}

Sau lớp tích chập, kết quả thường được đưa qua một hàm kích hoạt để bổ
sung tính phi tuyến cho mô hình. Nếu không có hàm kích hoạt, mạng dù có
nhiều tầng vẫn chủ yếu là sự kết hợp của các phép biến đổi tuyến tính,
khó học được các quan hệ phức tạp trong dữ liệu ảnh.

Một trong những hàm kích hoạt được sử dụng phổ biến trong CNN là ReLU.
Hàm ReLU giữ nguyên các giá trị dương và đưa các giá trị âm về 0:

\[
\operatorname{ReLU}(x)=\max(0,x)
\]

Trong bài toán phân loại ảnh tỏi, hàm kích hoạt giúp mô hình biểu diễn
tốt hơn các khác biệt phức tạp giữa các lớp. Ví dụ, vùng vỏ còn sót có
thể có màu sắc gần giống với bề mặt tép tỏi, hoặc dấu hiệu hư hỏng có
thể biểu hiện không đồng nhất giữa các mẫu ảnh. Nhờ tính phi tuyến, mô
hình có thể học được các quan hệ phức tạp hơn thay vì chỉ dựa vào sự
khác biệt tuyến tính đơn giản giữa các giá trị điểm ảnh.

\begin{enumerate}
\def\labelenumi{\arabic{enumi}.}
\setcounter{enumi}{1}
\item
  \protect\phantomsection\label{_Toc232112008}{}\emph{\textbf{Lớp
  pooling}}
\end{enumerate}

Lớp \textbf{pooling} thường được sử dụng sau một hoặc một số lớp tích
chập nhằm giảm kích thước không gian của bản đồ đặc trưng. Việc giảm
kích thước này giúp giảm chi phí tính toán, hạn chế số lượng tham số ở
các tầng sau và làm cho mô hình ít nhạy hơn với những thay đổi nhỏ về vị
trí của đặc trưng trong ảnh. Stanford CS231n trình bày pooling là một
trong các lớp cơ bản của CNN, thường được dùng để thực hiện thao tác
giảm mẫu theo chiều rộng và chiều cao của bản đồ đặc trưng {[}35{]}.

Hai dạng pooling phổ biến là \textbf{Max Pooling} và \textbf{Average
Pooling}. Max Pooling lấy giá trị lớn nhất trong một vùng cục bộ, nhờ đó
giữ lại đặc trưng nổi bật nhất trong vùng đó. Average Pooling lấy giá
trị trung bình, giúp tổng hợp thông tin theo cách mềm hơn.

Đối với ảnh tỏi, pooling giúp giữ lại các đặc trưng quan trọng trong khi
giảm kích thước dữ liệu trung gian. Chẳng hạn, nếu một vùng nhỏ trên ảnh
có dấu hiệu hư hỏng hoặc còn sót vỏ, các lớp tích chập có thể phát hiện
đặc trưng này, còn pooling giúp giữ lại tín hiệu nổi bật để các tầng sau
tiếp tục xử lý.

\begin{enumerate}
\def\labelenumi{\arabic{enumi}.}
\setcounter{enumi}{2}
\item
  \protect\phantomsection\label{_Toc232112009}{}\emph{\textbf{Lớp phân
  loại}}
\end{enumerate}

Sau khi ảnh đầu vào đi qua nhiều lớp tích chập, hàm kích hoạt và
pooling, mô hình thu được các bản đồ đặc trưng biểu diễn thông tin quan
trọng của ảnh. Các đặc trưng này sau đó được đưa vào phần phân loại để
dự đoán nhãn đầu ra. Trong các mô hình CNN truyền thống, phần cuối
thường gồm một hoặc nhiều lớp kết nối đầy đủ, sau đó là lớp đầu ra.

Với bài toán phân loại nhiều lớp, lớp đầu ra thường sử dụng hàm
\textbf{Softmax} để chuyển các giá trị đầu ra thành xác suất thuộc về
từng lớp. Tổng xác suất của tất cả các lớp bằng 1, và lớp có xác suất
cao nhất thường được chọn làm kết quả dự đoán cuối cùng.

Trong bài toán phân loại chất lượng tỏi, lớp đầu ra cần dự đoán một
trong ba nhóm: \textbf{tỏi sạch vỏ}, \textbf{tỏi còn vỏ} hoặc
\textbf{tỏi hư}. Nếu xác suất cao nhất thuộc về lớp tỏi hư, ảnh đầu vào
sẽ được phân vào nhóm tỏi hư. Tương tự, nếu xác suất cao nhất thuộc về
lớp tỏi sạch vỏ hoặc tỏi còn vỏ, mô hình sẽ đưa ra nhãn tương ứng.

\begin{enumerate}
\def\labelenumi{\arabic{enumi}.}
\setcounter{enumi}{3}
\item
  \protect\phantomsection\label{_Toc232112010}{}\emph{\textbf{Vai trò
  của CNN trong phân loại ảnh}}
\end{enumerate}

CNN đã trở thành nền tảng quan trọng trong nhiều hệ thống phân loại ảnh
hiện đại. Một dấu mốc quan trọng là mô hình CNN sâu của Krizhevsky,
Sutskever và Hinton trong bài toán phân loại ImageNet, cho thấy mạng
nơ-ron tích chập có thể đạt hiệu quả vượt trội trong phân loại ảnh quy
mô lớn.

Trong nghiên cứu này, CNN được sử dụng làm nền tảng cho việc phân loại
chất lượng tỏi. Lý do chính là ảnh tỏi chứa nhiều đặc trưng thị giác cần
được khai thác ở nhiều mức độ khác nhau. Các đặc điểm như màu sắc, kết
cấu bề mặt, vùng vỏ còn sót và dấu hiệu hư hỏng đều có thể được học
thông qua các tầng tích chập. Nhờ đó, CNN giúp giảm sự phụ thuộc vào
việc thiết kế đặc trưng thủ công và phù hợp với hướng xây dựng mô hình
phân loại tự động.

\begin{enumerate}
\def\labelenumi{\arabic{enumi}.}
\setcounter{enumi}{4}
\item
  \protect\phantomsection\label{_Toc232112011}{}\textbf{Học chuyển giao
  trong phân loại ảnh}
\end{enumerate}

Học chuyển giao, hay \textbf{Transfer Learning}, là một kỹ thuật quan
trọng trong học sâu, đặc biệt khi dữ liệu huấn luyện của bài toán mới
không quá lớn. Thay vì huấn luyện một mô hình từ đầu với các trọng số
khởi tạo ngẫu nhiên, học chuyển giao tận dụng lại một mô hình đã được
huấn luyện trước trên một bộ dữ liệu lớn. Sau đó, mô hình này được điều
chỉnh để phù hợp với bài toán cụ thể cần giải quyết.

Trong lĩnh vực phân loại ảnh, học chuyển giao thường sử dụng các mô hình
đã được huấn luyện trước trên các bộ dữ liệu quy mô lớn như ImageNet.
Các mô hình này đã học được nhiều đặc trưng thị giác cơ bản và tổng
quát, chẳng hạn cạnh, màu sắc, đường viền, hình dạng và kết cấu. Khi áp
dụng vào một bài toán mới, các đặc trưng đã học trước đó có thể được tái
sử dụng, giúp mô hình huấn luyện nhanh hơn và đạt hiệu quả tốt hơn so
với việc học hoàn toàn từ đầu, đặc biệt trong trường hợp bộ dữ liệu mới
có kích thước hạn chế. Stanford CS231n trình bày học chuyển giao như một
chiến lược phổ biến trong thực tế, trong đó người dùng có thể sử dụng mô
hình CNN đã huấn luyện trước làm bộ trích xuất đặc trưng hoặc tiếp tục
tinh chỉnh mô hình trên bộ dữ liệu mới {[}35{]}.

\includegraphics[width=6.2in,height=4.22441in]{./media/image9.png}

\textbf{\emph{Hình} 2.4} \emph{Mô hình Transfer Learning từ ImageNet}
{[}14{]}, {[}22{]}.

\begin{enumerate}
\def\labelenumi{\arabic{enumi}.}
\item
  \protect\phantomsection\label{_Toc232112012}{}\emph{\textbf{Khái niệm
  học chuyển giao}}
\end{enumerate}

Học chuyển giao có thể hiểu là quá trình chuyển tri thức đã học từ một
bài toán nguồn sang một bài toán đích có liên quan. Trong phân loại ảnh,
bài toán nguồn thường là bài toán phân loại ảnh quy mô lớn, còn bài toán
đích là bài toán cụ thể hơn với tập dữ liệu nhỏ hơn. Mục tiêu của học
chuyển giao là tận dụng những đặc trưng tổng quát mà mô hình đã học được
từ bài toán nguồn để cải thiện hiệu quả học trên bài toán đích.

Một mô hình CNN được huấn luyện trên tập dữ liệu lớn thường học được các
đặc trưng ở nhiều mức độ khác nhau. Các tầng đầu của mô hình thường học
những đặc trưng cơ bản như cạnh, vùng sáng tối hoặc màu sắc. Các tầng
sâu hơn học những đặc trưng có mức trừu tượng cao hơn như kết cấu, bộ
phận của vật thể hoặc đặc điểm liên quan đến từng lớp. Khi chuyển sang
bài toán mới, các đặc trưng cơ bản ở những tầng đầu vẫn có giá trị sử
dụng, vì nhiều loại ảnh khác nhau đều chia sẻ các đặc điểm thị giác
chung.

Đối với bài toán phân loại chất lượng tỏi, số lượng ảnh trong bộ dữ liệu
không quá lớn so với các bộ dữ liệu ảnh quy mô lớn. Vì vậy, nếu huấn
luyện một mô hình sâu từ đầu, mô hình có thể cần nhiều dữ liệu và thời
gian hơn để học được các đặc trưng ổn định. Học chuyển giao giúp giảm
bớt khó khăn này bằng cách sử dụng các trọng số đã học từ trước làm điểm
khởi đầu cho quá trình huấn luyện.

\begin{enumerate}
\def\labelenumi{\arabic{enumi}.}
\setcounter{enumi}{1}
\item
  \protect\phantomsection\label{_Toc232112013}{}\emph{\textbf{Pre-trained
  model trong phân loại ảnh}}
\end{enumerate}

\textbf{Pre-trained model} là mô hình đã được huấn luyện trước trên một
bộ dữ liệu lớn. Trong phân loại ảnh, các mô hình này thường được huấn
luyện trên ImageNet hoặc các bộ dữ liệu thị giác quy mô lớn khác. Nhờ
quá trình huấn luyện trước, mô hình đã học được các đặc trưng thị giác
tổng quát có thể tái sử dụng cho nhiều bài toán khác nhau.

Khi sử dụng pre-trained model cho một bài toán mới, phần đầu của mô hình
thường được giữ lại để làm bộ trích xuất đặc trưng. Phần phân loại cuối
cùng của mô hình được thay thế bằng một lớp đầu ra mới phù hợp với số
lớp của bài toán đích. Ví dụ, một mô hình ban đầu có thể được huấn luyện
để phân loại hàng nghìn lớp ảnh, nhưng khi áp dụng vào bài toán phân
loại tỏi, lớp đầu ra cần được thay đổi để dự đoán ba lớp: \textbf{tỏi
sạch vỏ}, \textbf{tỏi còn vỏ} và \textbf{tỏi hư}.

TensorFlow hướng dẫn học chuyển giao theo cách sử dụng mô hình đã huấn
luyện trước làm nền tảng, sau đó thêm các lớp phân loại mới và huấn
luyện trên bộ dữ liệu riêng của bài toán đích. Đây là cách tiếp cận phổ
biến trong các bài toán phân loại ảnh khi dữ liệu mới có quy mô vừa hoặc
nhỏ {[}37{]}.

\begin{enumerate}
\def\labelenumi{\arabic{enumi}.}
\setcounter{enumi}{2}
\item
  \protect\phantomsection\label{_Toc232112014}{}\emph{\textbf{Các chiến
  lược sử dụng học chuyển giao}}
\end{enumerate}

Trong phân loại ảnh, học chuyển giao thường được triển khai theo hai
chiến lược chính: sử dụng mô hình như bộ trích xuất đặc trưng và tinh
chỉnh mô hình.

Ở chiến lược thứ nhất, các tầng của mô hình đã huấn luyện trước được giữ
nguyên, tức là không cập nhật trọng số trong quá trình huấn luyện trên
dữ liệu mới. Mô hình lúc này đóng vai trò như một bộ trích xuất đặc
trưng, còn phần phân loại mới ở phía sau sẽ được huấn luyện để phù hợp
với bài toán đích. Cách này thường được dùng khi bộ dữ liệu mới có kích
thước nhỏ hoặc khi muốn giảm nguy cơ quá khớp.

Ở chiến lược thứ hai, sau khi huấn luyện phần phân loại mới, một số tầng
phía sau của mô hình pre-trained được mở khóa để tiếp tục cập nhật trọng
số trên dữ liệu mới. Quá trình này được gọi là \textbf{fine-tuning}.
Fine-tuning giúp mô hình điều chỉnh các đặc trưng đã học trước đó sao
cho phù hợp hơn với bài toán cụ thể. Stanford CS231n cho biết có thể
fine-tune toàn bộ các tầng của ConvNet hoặc giữ cố định một số tầng đầu
để hạn chế quá khớp, chỉ tinh chỉnh các tầng cao hơn của mạng {[}35{]}.

\begin{enumerate}
\def\labelenumi{\arabic{enumi}.}
\setcounter{enumi}{3}
\item
  \protect\phantomsection\label{_Toc232112015}{}\emph{\textbf{Fine-tuning
  trong phân loại ảnh}}
\end{enumerate}

Fine-tuning là bước tinh chỉnh mô hình đã huấn luyện trước trên bộ dữ
liệu của bài toán mới. Trong quá trình này, các trọng số của một phần
hoặc toàn bộ mô hình được cập nhật để thích nghi với đặc điểm dữ liệu
mới. Đối với bài toán phân loại ảnh, fine-tuning thường được thực hiện
sau khi phần phân loại mới đã được huấn luyện ổn định.

Lý do không nên fine-tune toàn bộ mô hình ngay từ đầu là vì các tầng đầu
của CNN thường học những đặc trưng cơ bản và tổng quát. Nếu cập nhật tất
cả trọng số quá sớm, mô hình có thể làm mất đi những đặc trưng hữu ích
đã học được từ bộ dữ liệu lớn. Vì vậy, trong nhiều trường hợp, người
nghiên cứu thường đóng băng các tầng đầu, huấn luyện phần phân loại mới,
sau đó mới mở khóa một số tầng cuối để tinh chỉnh nhẹ với tốc độ học nhỏ
hơn.

Trong bài toán phân loại tỏi, fine-tuning giúp mô hình thích nghi tốt
hơn với các đặc trưng riêng của ảnh tỏi, chẳng hạn bề mặt tép tỏi, vùng
vỏ còn sót, vết thâm hoặc vùng hư hỏng. Các đặc trưng này khác với các
đối tượng trong bộ dữ liệu gốc, nên việc tinh chỉnh mô hình trên dữ liệu
tỏi là cần thiết để nâng cao hiệu quả phân loại.

\begin{enumerate}
\def\labelenumi{\arabic{enumi}.}
\setcounter{enumi}{4}
\item
  \protect\phantomsection\label{_Toc232112016}{}\emph{\textbf{Vai trò
  của học chuyển giao trong nghiên cứu}}
\end{enumerate}

Trong luận văn này, học chuyển giao được sử dụng nhằm tận dụng sức mạnh
của các mô hình CNN đã được huấn luyện trước. Các mô hình này đóng vai
trò là nền tảng để trích xuất đặc trưng ảnh, sau đó được điều chỉnh cho
bài toán phân loại chất lượng tỏi.

Việc sử dụng học chuyển giao mang lại một số lợi ích cho nghiên cứu. Thứ
nhất, học chuyển giao giúp giảm thời gian huấn luyện vì mô hình không
cần học mọi đặc trưng từ đầu. Thứ hai, mô hình có thể đạt hiệu quả tốt
hơn khi bộ dữ liệu tỏi không quá lớn. Thứ ba, việc sử dụng các mô hình
pre-trained giúp quá trình so sánh giữa các kiến trúc CNN nền trở nên
thuận lợi hơn, vì các mô hình đều bắt đầu từ những trọng số đã học được
trên dữ liệu ảnh quy mô lớn.

Đối với bài toán phân loại chất lượng tỏi, học chuyển giao đặc biệt phù
hợp vì dữ liệu ảnh có những đặc trưng thị giác phức tạp nhưng số lượng
mẫu còn hạn chế. Nhờ học chuyển giao, các mô hình có thể tận dụng các
đặc trưng tổng quát đã học được trước đó, sau đó điều chỉnh để nhận diện
tốt hơn ba nhóm chất lượng: \textbf{tỏi sạch vỏ}, \textbf{tỏi còn vỏ} và
\textbf{tỏi hư}.

\begin{enumerate}
\def\labelenumi{\arabic{enumi}.}
\setcounter{enumi}{5}
\item
  \protect\phantomsection\label{_Toc232112017}{}\textbf{Các mô hình CNN
  được sử dụng trong nghiên cứu}
\end{enumerate}

Để có cơ sở lựa chọn mô hình phù hợp cho bài toán phân loại chất lượng
tỏi, luận văn tiến hành thực nghiệm với một số kiến trúc CNN đã được sử
dụng rộng rãi trong phân loại ảnh. Các mô hình được lựa chọn gồm
\emph{\textbf{NASNetMobile, InceptionV3, EfficientNetB4, MobileNetV2 và
DenseNet121}}. Đây không chỉ là các mô hình phổ biến trong thư viện học
sâu hiện nay, mà còn đại diện cho những hướng thiết kế khác nhau của
mạng nơ-ron tích chập.

Việc sử dụng nhiều mô hình nền giúp quá trình đánh giá không bị phụ
thuộc vào một kiến trúc duy nhất. Mỗi mô hình có ưu điểm riêng: có mô
hình hướng đến độ chính xác, có mô hình tối ưu tài nguyên tính toán, có
mô hình khai thác đặc trưng đa tỷ lệ, và cũng có mô hình chú trọng việc
tái sử dụng đặc trưng giữa các tầng. Từ kết quả thực nghiệm trên cùng bộ
dữ liệu ảnh tỏi, luận văn có thể xác định mô hình nào phù hợp hơn với
đặc điểm của bài toán, từ đó lựa chọn làm nền tảng cho mô hình đề xuất ở
các bước tiếp theo.

\begin{enumerate}
\def\labelenumi{\arabic{enumi}.}
\item
  \protect\phantomsection\label{_Toc232112018}{}\emph{\textbf{Mô hình
  NASNetMobile}}
\end{enumerate}

NASNetMobile là một phiên bản thuộc họ kiến trúc NASNet, viết tắt của
\textbf{Neural Architecture Search Network}. Đây là nhóm mô hình được
xây dựng dựa trên ý tưởng tìm kiếm kiến trúc mạng nơ-ron một cách tự
động, thay vì thiết kế hoàn toàn thủ công như nhiều mô hình CNN truyền
thống. Trong đó, NASNetMobile được xem là phiên bản gọn nhẹ hơn của
NASNet, hướng đến các môi trường có tài nguyên tính toán hạn chế như
thiết bị di động hoặc các hệ thống cần tốc độ xử lý nhanh.

Điểm khác biệt quan trọng của NASNetMobile nằm ở cách hình thành kiến
trúc mạng. Thay vì con người trực tiếp quyết định toàn bộ số lớp, kiểu
kết nối và phép toán bên trong mô hình, NASNet sử dụng một quá trình tìm
kiếm tự động gọi là \textbf{Neural Architecture Search -- NAS}. Quá
trình này cho phép hệ thống thử nghiệm nhiều cấu trúc mạng con khác
nhau, đánh giá hiệu quả của từng cấu trúc, sau đó dần tìm ra những thiết
kế có khả năng hoạt động tốt hơn. Nhờ vậy, NASNetMobile mang đặc trưng
của một kiến trúc được khám phá tự động, nhưng vẫn có thể triển khai
dưới dạng một mô hình CNN hoàn chỉnh cho bài toán phân loại ảnh.

Về mặt lý thuyết, quá trình tìm kiếm kiến trúc của NASNet sử dụng một
mạng nơ-ron tuần hoàn, thường được gọi là \textbf{bộ điều khiển}. Bộ
điều khiển này có nhiệm vụ sinh ra các kiến trúc mạng con. Mỗi kiến trúc
mạng con sau đó được huấn luyện và đánh giá trên tập xác thực. Độ chính
xác của mạng con trên tập xác thực được xem như tín hiệu phản hồi để cập
nhật lại bộ điều khiển. Qua nhiều vòng lặp, bộ điều khiển có xu hướng
tạo ra những kiến trúc tốt hơn, vì các thiết kế cho kết quả cao sẽ được
ưu tiên với xác suất lớn hơn trong các lần sinh kiến trúc tiếp theo.

Một điểm quan trọng trong NASNet là không tìm kiếm toàn bộ kiến trúc
mạng lớn ngay từ đầu. Cách làm này sẽ tiêu tốn rất nhiều tài nguyên tính
toán, đặc biệt khi áp dụng trực tiếp trên các bộ dữ liệu ảnh lớn như
ImageNet. Thay vào đó, NASNet tìm kiếm một khối kiến trúc cơ sở, gọi là
\textbf{cell}, trên một bộ dữ liệu nhỏ hơn như CIFAR-10. Sau khi tìm
được cell phù hợp, các cell này được chuyển sang bài toán lớn hơn bằng
cách xếp chồng nhiều lần để tạo thành mạng hoàn chỉnh. Đây chính là yếu
tố thể hiện tính chuyển giao trong thiết kế của NASNet: kiến trúc được
tìm trên bài toán nhỏ có thể được mở rộng và sử dụng cho bài toán có quy
mô lớn hơn.

Kiến trúc NASNet được xây dựng từ hai loại cell chính là \textbf{Normal
Cell} và \textbf{Reduction Cell}. Normal Cell có nhiệm vụ xử lý và trích
xuất đặc trưng nhưng vẫn giữ nguyên kích thước không gian của bản đồ đặc
trưng đầu vào. Ngược lại, Reduction Cell làm giảm kích thước chiều cao
và chiều rộng của bản đồ đặc trưng, thường thông qua các phép toán có
bước nhảy lớn hơn. Nhờ sự kết hợp giữa hai loại cell này, mô hình vừa có
thể duy trì thông tin không gian cần thiết ở các tầng đầu, vừa dần nén
và biểu diễn đặc trưng ở mức trừu tượng cao hơn ở các tầng sâu.

Đối với NASNetMobile, phiên bản thường được sử dụng là cấu hình
\textbf{NASNet-A 4 @ 1056}. Trong cấu hình này, Normal Cell được lặp lại
bốn lần giữa các Reduction Cell, còn số lượng bộ lọc ở lớp gần cuối là
1056. Khi kích thước không gian của bản đồ đặc trưng giảm xuống, số
lượng bộ lọc thường được tăng lên để giữ lại lượng thông tin biểu diễn
cần thiết. Cách tổ chức này giúp mô hình kiểm soát được chi phí tính
toán, đồng thời vẫn đảm bảo khả năng học các đặc trưng quan trọng từ ảnh
đầu vào.

Bên trong mỗi cell, kiến trúc chi tiết không được thiết kế thủ công theo
một mẫu cố định, mà là kết quả của quá trình tìm kiếm tự động. Bộ điều
khiển sẽ lựa chọn các đầu vào, các phép toán và cách kết hợp giữa chúng
để tạo thành cấu trúc của cell. Các phép toán thường gặp trong NASNet
gồm tích chập tách rời theo chiều sâu, max pooling, average pooling và
các phép kết hợp đặc trưng. Trong đó, tích chập tách rời theo chiều sâu
được sử dụng nhiều vì có khả năng giảm chi phí tính toán so với tích
chập thông thường, phù hợp với định hướng xây dựng mô hình gọn nhẹ.

Workflow tổng quát của NASNetMobile có thể hiểu như sau: ảnh đầu vào
trước hết đi qua các lớp tích chập ban đầu để tạo bản đồ đặc trưng ở mức
thấp. Sau đó, đặc trưng được truyền qua nhiều nhóm Normal Cell và
Reduction Cell. Normal Cell giúp mô hình tiếp tục khai thác đặc trưng mà
không làm thay đổi kích thước không gian, trong khi Reduction Cell giảm
kích thước bản đồ đặc trưng và giúp mô hình học các biểu diễn sâu hơn.
Cuối cùng, đặc trưng được tổng hợp qua lớp Global Average Pooling và đưa
vào lớp phân loại để dự đoán nhãn đầu ra.

\includegraphics[width=6.36333in,height=3.26986in]{./media/image10.png}

\textbf{\emph{Hình} 2.5} \emph{Cấu trúc Normal Cell và Reduction Cell
trong NASNet} {[}1{]}.

Trong quá trình huấn luyện NASNet, một kỹ thuật chính quy hóa đáng chú ý
là \textbf{ScheduledDropPath}. Kỹ thuật này có ý tưởng gần với DropPath,
tức là loại bỏ ngẫu nhiên một số nhánh kết nối trong mạng để giảm hiện
tượng quá khớp. Tuy nhiên, thay vì giữ xác suất loại bỏ cố định trong
suốt quá trình huấn luyện, ScheduledDropPath tăng dần xác suất này theo
thời gian. Cách làm này giúp mô hình học ổn định hơn ở giai đoạn đầu,
sau đó tăng mức độ chính quy hóa ở các giai đoạn sau để cải thiện khả
năng tổng quát hóa.

\begin{enumerate}
\def\labelenumi{\arabic{enumi}.}
\setcounter{enumi}{1}
\item
  \protect\phantomsection\label{_Toc232112019}{}\emph{\textbf{Mô hình
  Inception}}
\end{enumerate}

InceptionV3 là một kiến trúc mạng nơ-ron tích chập thuộc họ Inception,
được phát triển nhằm nâng cao hiệu quả trích xuất đặc trưng trong các
bài toán phân loại ảnh. Kiến trúc này kế thừa ý tưởng từ GoogLeNet,
trong đó mạng không xử lý đặc trưng theo một luồng tuyến tính duy nhất,
mà tổ chức các phép toán theo nhiều nhánh song song. Nhờ cơ chế này, mô
hình có thể khai thác đồng thời nhiều dạng đặc trưng ở các mức tỷ lệ
khác nhau, từ đặc trưng cục bộ nhỏ đến các biểu diễn có phạm vi không
gian rộng hơn. Để xây dựng cấu trúc này một cách tối ưu, nhóm tác giả đã
đề xuất 4 nguyên tắc thiết kế nền tảng:

\begin{itemize}
\item
  Tránh thắt nút cổ chai biểu diễn (Avoid representational bottlenecks):
  Kích thước biểu diễn (thông tin) nên được giảm đi một cách từ từ từ
  đầu vào cho đến đầu ra, đặc biệt là ở các lớp đầu của mạng, nhằm tránh
  việc nén thông tin quá gắt gao làm mất mát đặc trưng.
\item
  Biểu diễn chiều cao dễ xử lý cục bộ hơn: Tăng số lượng kênh kích hoạt
  (activations) trên mỗi ô (tile) trong mạng tích chập giúp các đặc
  trưng được phân tách rõ ràng hơn, từ đó tăng tốc độ huấn luyện.
\item
  Gộp không gian trên các embedding chiều thấp: Việc giảm chiều của biểu
  diễn đầu vào trước khi thực hiện các phép tích chập lớn (như 3x3)
  không làm suy giảm nghiêm trọng sức mạnh biểu diễn, do các tín hiệu
  lân cận thường có độ tương quan cao.
\item
  Cân bằng giữa chiều rộng và chiều sâu: Hiệu suất tối ưu của mạng đạt
  được khi phân bổ đồng đều ngân sách tính toán cho cả việc tăng số
  lượng bộ lọc (chiều rộng) và tăng số lớp (chiều sâu).
\end{itemize}

Kỹ thuật phân rã phép tích chập (Factorizing Convolutions) là cốt lõi
của Inception-v3 đến từ các chiến lược phân rã (factorization) phép toán
tích chập tinh vi nhằm giảm số lượng tham số trong khi vẫn giữ nguyên,
hoặc tăng cường, sức mạnh biểu diễn:

\includegraphics[width=3.97972in,height=2.77122in]{./media/image11.png}

\textbf{\emph{Hình} 2.6} \emph{Mô-đun Inception với phép tích chập 5x5
được phân rã thành hai phép tích chập 3x3} {[}7{]}.

\begin{itemize}
\item
  Các bộ lọc kích thước lớn như 5x5 có chi phí tính toán đắt đỏ (gấp
  2.78 lần so với 3x3). Do đó, Inception-v3 thay thế một lớp tích chập
  5x5 bằng một chuỗi hai lớp tích chập 3x3. Việc này không những giảm
  thiểu đáng kể số lượng tham số mà còn làm tăng không gian biến thiên
  nhờ việc sử dụng hàm kích hoạt phi tuyến tính (ReLU) sau mỗi lớp.
\end{itemize}

\includegraphics[width=4.03181in,height=4.03181in]{./media/image12.png}

\textbf{\emph{Hình} 2.7} \emph{Phân rã phép tích chập bất đối xứng trong
Inception-v3} {[}7{]}.

\begin{itemize}
\item
  Mô hình tiến xa hơn bằng cách thay thế các phép tích chập n×n thành
  một phép chập 1×n theo sau bởi một phép chập n×1. Ví dụ, một phép chập
  3x3 được phân rã thành 1x3 và 3x1, giúp tiết kiệm 33\% chi phí tính
  toán. Kỹ thuật này đặc biệt hiệu quả trên các lưới đặc trưng có kích
  thước trung bình (từ 12×12 đến 20×20), nơi Inception-v3 áp dụng mạnh
  mẽ các phép chập bất đối xứng kích thước 1x7 và 7x1.
\end{itemize}

\includegraphics[width=5.39659in,height=3.03167in]{./media/image13.png}

\textbf{\emph{Hình} 2.8} \emph{Cấu trúc mô-đun Inception giúp giảm kích
thước lưới và mở rộng tập hợp bộ lọc} {[}7{]}.

Trong các CNN truyền thống, để giảm kích thước không gian của bản đồ đặc
trưng (feature map), người ta thường dùng phép gộp (pooling). Tuy nhiên,
nếu áp dụng pooling trực tiếp sẽ tạo ra nút thắt cổ chai biểu diễn.
Ngược lại, nếu tăng kích thước bộ lọc trước khi pooling thì lại tốn kém
quá nhiều chi phí tính toán. Để giải quyết điều này, Inception-v3 đề
xuất một module giảm lưới hiệu quả: Thiết kế hai nhánh song song, một
nhánh sử dụng phép tích chập với bước nhảy (stride) bằng 2, và nhánh còn
lại sử dụng phép gộp (pooling) với bước nhảy bằng 2. Sau đó, kết quả của
hai nhánh được nối (concatenate) lại với nhau. Kỹ thuật này giúp giảm
không gian lưới một cách rẻ tiền nhưng vẫn mở rộng được số lượng bộ lọc
để tránh thắt nút cổ chai.

\textbf{Bộ phân loại phụ (Auxiliary Classifiers):} Các kiến trúc
Inception trước đây sử dụng các bộ phân loại phụ ở giữa mạng để đẩy
gradient về các lớp thấp, chống lại hiện tượng triệt tiêu gradient
(vanishing gradient). Trong Inception-v3, bộ phân loại phụ còn đóng vai
trò như một cơ chế chính quy hóa (regularizer), đặc biệt khi được kết hợp
với Batch Normalization hoặc Dropout.

\textbf{Label Smoothing Regularization (LSR):} Để hạn chế việc mô hình
gán xác suất quá tuyệt đối cho một nhãn duy nhất, Inception-v3 sử dụng
kỹ thuật làm mịn nhãn. LSR thay thế phân phối nhãn gốc bằng một hỗn hợp
giữa nhãn đúng và một phân phối cố định, thường là phân phối đồng đều.
Cách làm này giúp mô hình giảm quá tự tin, hạn chế quá khớp và cải
thiện khả năng thích ứng với dữ liệu mới.

\begin{enumerate}
\def\labelenumi{\arabic{enumi}.}
\item
  \protect\phantomsection\label{_Toc232112020}{}\emph{\textbf{Mô hình
  EfficientNet}}
\end{enumerate}

EfficientNet là một trong những kiến trúc mạng nơ-ron tích chập (CNN)
nổi bật nhất gần đây, được nhóm nghiên cứu Google Brain giới thiệu vào
năm 2019. Mô hình này không chỉ đạt hiệu suất cao mà còn tối ưu hóa tài
nguyên tính toán, giúp hoạt động hiệu quả trên cả các thiết bị có phần
cứng hạn chế. Các kiến trúc CNN cổ điển như AlexNet, VGG và ResNet đã
tạo ra những bước tiến lớn trong hiệu suất của hệ thống nhận diện hình
ảnh. Tuy nhiên, những mô hình này thường gặp vấn đề về độ phức tạp cao
và yêu cầu tài nguyên phần cứng lớn.

Để nâng cao độ chính xác, các nhà nghiên cứu thường tăng độ sâu (depth),
số kênh (width) hoặc độ phân giải của ảnh đầu vào (resolution) của mô
hình. Tuy nhiên, việc mở rộng này không phải lúc nào cũng hiệu quả,
thường dẫn đến nguy cơ quá khớp (overfitting) hoặc đòi hỏi bộ nhớ lớn
hơn. EfficientNet được thiết kế để khắc phục những hạn chế này thông qua
phương pháp mở rộng cân bằng (compound scaling). Thay vì chỉ tập trung
vào một yếu tố như độ sâu, số kênh hoặc độ phân giải, EfficientNet điều
chỉnh đồng thời tất cả các yếu tố này theo một tỷ lệ mở rộng tối ưu, dựa
trên sự cân bằng giữa hiệu suất và tài nguyên phần cứng.

Kiến trúc của EfficientNet được xây dựng dựa trên nền tảng của
MobileNetV2 và MnasNet, hai mô hình được tối ưu hóa cho các thiết bị di
động. EfficientNet áp dụng phương pháp mở rộng cân bằng (compound
scaling) để phát triển từ phiên bản cơ sở (EfficientNet-B0) thành các
phiên bản lớn hơn (EfficientNet-B1 đến EfficientNet-B7), trong đó
EfficientNet-B7 là phiên bản lớn nhất và đạt độ chính xác cao nhất.

\includegraphics[width=6.19853in,height=2.53908in]{./media/image14.png}

\emph{\textbf{Hình 2.9.} Kiến trúc tổng thể của EfficientNet}
{[}11{]}.\\
Hình 2.9: Mô tả kiến trúc mô hình EfficientNet sử dụng chuỗi block
MBConv với kích

thước kernel khác nhau để trích xuất đặc trưng hiệu quả và tối ưu tài
nguyên.

\textbf{Compound Scaling trong EfficientNet}: là một phương pháp mở rộng
mô hình theo ba chiều: độ sâu (depth), số kênh (width), và độ phân giải
ảnh (resolution). Thay vì chỉ tăng một chiều như các phương pháp truyền
thống, EfficientNet đồng thời điều chỉnh cả ba chiều này theo một tỷ lệ
cân bằng đã được tối ưu hóa. Cách tiếp cận này cải thiện hiệu suất mô
hình mà không làm tăng quá mức số lượng tham số hoặc yêu cầu tài nguyên
tính toán, mang lại độ chính xác cao hơn với hiệu quả tính toán tối ưu.

\textbf{Mobile Inverted Bottleneck Convolution (MBConv):} Một trong
những điểm nổi bật của EfficientNet là sử dụng các khối Mobile Inverted
Bottleneck Convolution (MBConv), tương tự như trong MobileNetV2. Khối
MBConv kết hợp các phép tích chập điểm ảnh 1 × 1 (pointwise convolution)
và tích chập phân tách theo chiều sâu (depthwise convolution), giúp giảm
số lượng tham số mà vẫn duy trì độ chính xác cao.

~\textbf{Squeeze-and-Excitation (SE) Blocks trong EfficientNet:}
EfficientNet tích hợp các khối Squeeze-and-Excitation (SE) để tự động
điều chỉnh trọng số của các kênh, giúp mô hình tập trung vào các kênh
quan trọng hơn. Khối SE hoạt động bằng cách học các hệ số điều chỉnh
(attention) cho từng kênh, dựa trên giá trị tổng hợp từ toàn bộ không
gian đầu vào.

\textbf{Các phiên bản EfficientNet:} EfficientNet là một mô hình mạng
nơ-ron tích chập (CNN) nổi bật, được thiết kế nhằm tối ưu hóa hiệu suất
đồng thời giảm thiểu tài nguyên tính toán. Mô hình này bao gồm nhiều
phiên bản, từ EfficientNet-B0 đến EfficientNet-B7, mỗi phiên bản được mở
rộng từ phiên bản trước thông qua phương pháp mở rộng cân bằng (compound
scaling). Đây là đặc điểm nổi bật giúp EfficientNet duy trì sự cân bằng
giữa độ chính xác và tài nguyên phần cứng khi kích thước mô hình tăng
lên.

\textbf{EfficientNet-B0:} EfficientNet-B0 là phiên bản đầu tiên và nhỏ
nhất trong dòng EfficientNet, được xem là nền tảng để mở rộng các phiên
bản tiếp theo. Kiến trúc của EfficientNet-B0 được tối ưu hóa thông qua
việc sử dụng tìm kiếm cấu trúc neural architecture search (NAS), nhằm
xác định cấu trúc mô hình tối ưu nhất về độ chính xác và hiệu quả tính
toán.

\includegraphics[width=6.2in,height=1.91038in]{./media/image15.png}

\emph{\textbf{Hình 2.10.} Cấu trúc chi tiết của EfficientNet với các
giai đoạn} {[}11{]}.

Hình 2.10: Mô tả chi tiết kiến trúc EfficientNet theo từng giai đoạn:
EfficientNet gồm 8 stage, sử dụng các block MBConv với kích thước kernel
(3 × 3, 5 × 5) và số kênh tăng dần, giúp mở rộng độ sâu và độ rộng mạng
một cách hiệu quả.

\textbf{Các phiên bản của mô hình EfficientNet:}

\begin{itemize}
\item
  \textbf{EfficientNet-B1:} EfficientNet-B1, phiên bản mở rộng đầu tiên
  của dòng EfficientNet, tăng độ sâu và số kênh của các lớp mạng, đồng
  thời nâng độ phân giải ảnh đầu vào từ 224 × 224 pixel lên 240 × 240
  pixel. Phiên bản này có khoảng 7,8 triệu tham số, yêu cầu 0,7 tỷ FLOPs
  và đạt độ chính xác khoảng 79.2\% trên ImageNet.
\item
  \textbf{EfficientNet-B2:} EfficientNet-B2 nâng độ phân giải đầu vào
  lên 260 × 260 pixel, đồng thời tăng thêm một số lớp mạng. Số lượng
  tham số và FLOPs tăng lần lượt lên 9,2 triệu và 1,0 tỷ, với độ chính
  xác đạt khoảng 80.3\% trên ImageNet\textbf{.}
\item
  \textbf{EfficientNet-B3:} EfficientNet-B3 sử dụng độ phân giải đầu vào
  300 × 300 pixel, với khoảng 12 triệu tham số và yêu cầu 1,8 tỷ FLOPs.
  Phiên bản này đạt độ chính xác 81.6\% trên ImageNet.
\item
  \textbf{EfficientNet-B4:} EfficientNet-B4 là một trong những phiên bản
  lớn hơn, sử dụng độ phân giải đầu vào 380 × 380 pixel, với khoảng 19
  triệu tham số và yêu cầu 4,2 tỷ FLOPs. Phiên bản này đạt độ chính xác
  khoảng 82.9\% trên ImageNet, vượt trội hơn nhiều mô hình truyền thống
  có quy mô lớn khác.
\item
  \textbf{EfficientNet-B5:} EfficientNet-B5 sử dụng ảnh đầu vào có độ
  phân giải 456 × 456 pixel, với khoảng 30 triệu tham số và yêu cầu 9,9
  tỷ FLOPs. Phiên bản này đạt độ chính xác 83.6\% trên ImageNet, một kết
  quả ấn tượng so với quy mô của mô hình.
\item
  \textbf{EfficientNet-B6:} EfficientNet-B6 sử dụng độ phân giải ảnh đầu
  vào 528 × 528 pixel, với khoảng 43 triệu tham số và yêu cầu 19,0 tỷ
  FLOPs. Phiên bản này đạt độ chính xác khoảng 84.0\% trên ImageNet.
\item
  \textbf{EfficientNet-B7:} EfficientNet-B7, phiên bản lớn nhất và mạnh
  mẽ nhất trong dòng EfficientNet, sử dụng độ phân giải đầu vào 600 ×
  600 pixel. Phiên bản này có khoảng 66 triệu tham số và yêu cầu 37,0 tỷ
  FLOPs cho mỗi dự đoán. Độ chính xác của EfficientNet-B7 đạt 84.4\%
  trên ImageNet, một trong những kết quả cao nhất trên tập dữ liệu này.

  \begin{enumerate}
  \def\labelenumi{\arabic{enumi}.}
  \item
    \protect\phantomsection\label{_Toc232112021}{}\emph{\textbf{Mô hình
    MobileNet}}
  \end{enumerate}
\end{itemize}

MobileNet là một họ kiến trúc mạng nơ-ron tích chập (CNN) được Google
giới thiệu lần đầu vào năm 2017 {[}24{]}, với mục tiêu tối ưu cho các
thiết bị di động và hệ thống nhúng vốn có hạn chế về tài nguyên tính
toán. Điểm nổi bật của MobileNet là khả năng duy trì độ chính xác trong
các tác vụ phân loại ảnh ở mức chấp nhận được, đồng thời giảm đáng kể số
lượng tham số và chi phí tính toán so với các kiến trúc CNN truyền
thống.

Nguyên lý hoạt động chính của MobileNet là sử dụng Depthwise Separable
Convolution, một kỹ thuật tách riêng quá trình tích chập không gian
(depthwise convolution) và kết hợp kênh (pointwise convolution). Cách
tiếp cận này giúp giảm đáng kể số lượng tham số và phép tính so với tích
chập truyền thống, nhưng vẫn giữ được khả năng trích xuất đặc trưng mạnh
mẽ. Ngoài ra, MobileNet cho phép điều chỉnh mức độ phức tạp của mô hình
thông qua hai siêu tham số là width multiplier (giảm số kênh feature
maps) và resolution multiplier (giảm kích thước đầu vào). Nhờ đó, mô
hình có thể được tùy biến theo yêu cầu về tài nguyên và độ chính xác.
MobileNet sử dụng Depthwise Separable Convolution thay vì Convolution
chuẩn. Phương pháp này tách quá trình chập thành 2 bước:

\begin{itemize}
\item
  \textbf{Depthwise Convolution:} Thực hiện phép tích chập độc lập trên
  từng kênh đầu vào.
\item
  \textbf{Pointwise Convolution (1 x 1 Convolution):} Kết hợp các đặc
  trưng từ bước Depthwise để tạo ra đầu ra mong muốn.
\end{itemize}

\includegraphics[width=6.2in,height=3.55231in]{./media/image16.png}

\emph{\textbf{Hình 2.11.} Standard Convolution và Depthwise Separable
Convolution trong MobileNet {[}24{]}}

Hình 2.11: Minh họa sự khác biệt giữa Standard Convolution và Depthwise
Separable Convolution. Ở Standard Convolution (bên trái), phép tích chập
3 x 3 được áp dụng đồng thời trên toàn bộ các kênh đầu vào, sau đó đi
qua Batch Normalization (BN) và ReLU. Cách này cho độ chính xác cao
nhưng tốn nhiều tham số và chi phí tính toán. Ở Depthwise Separable
Convolution (bên phải), quá trình được tách thành hai bước là Depthwise
và Pointwise Convolution, mỗi bước đều có BN và ReLU. Cách tiếp cận này
giúp giảm mạnh số tham số và phép tính, trong khi vẫn duy trì hiệu quả
mô hình.

\textbf{Các phiên bản MobileNet}

\begin{itemize}
\item
  \textbf{MobileNetV1 (2017):} Phiên bản đầu tiên, tập trung vào việc
  giảm số tham số bằng Depthwise Separable Convolution.
\item
  \textbf{MobileNetV2 (2018):} Giới thiệu cơ chế Inverted Residuals và
  Linear Bottlenecks, giúp cải thiện độ chính xác, ổn định huấn luyện và
  giảm overfitting. Đây là phiên bản được ứng dụng rộng rãi nhất trong
  các nghiên cứu thị giác máy tính.
\item
  \textbf{MobileNetV3 (2019):} Tích hợp tìm kiếm kiến trúc mạng thần
  nơ-ron (Neural Architecture Search -- NAS) và modules (SE --
  Squeeze-and-Excitation), giúp cân bằng tốt hơn giữa tốc độ, kích thước
  và hiệu suất.
\item
  \textbf{Ưu điểm của MobileNet:}
\end{itemize}

\begin{itemize}
\item
  Nhẹ, tốc độ nhanh, phù hợp cho thiết bị di động và IoT.
\item
  Giảm đáng kể số tham số và phép tính so với các mô hình CNN lớn như
  VGG hoặc ResNet.
\item
  Dễ dàng tích hợp vào các ứng dụng thời gian thực (real-time).
\end{itemize}

\begin{itemize}
\item
  \textbf{Hạn chế của MobileNet:}
\end{itemize}

\begin{itemize}
\item
  Độ chính xác thường thấp hơn một chút so với các mô hình lớn hơn
  (DenseNet, EfficientNet) trên cùng một tập dữ liệu.
\item
  Cần tinh chỉnh cẩn thận để tránh mất thông tin, do giảm số lượng tham
  số.
\end{itemize}

\begin{itemize}
\item
  \textbf{Ứng dụng thực tế của MobileNet:} Nhờ sự gọn nhẹ, MobileNet đã
  được áp dụng rộng rãi trong các tác vụ nhận diện hình ảnh trên thiết
  bị di động, phân loại ảnh y tế, giám sát thông minh, và đặc biệt là
  trong nông nghiệp chính xác (Precision Agriculture), nơi cần mô hình
  nhanh, hiệu quả, nhưng vẫn đảm bảo độ chính xác chấp nhận được.
\item
  \textbf{Đánh giá tổng quan:} MobileNet là một trong những kiến trúc
  tiêu biểu cho xu hướng ``lightweight deep learning models''. Ưu điểm
  chính là khả năng triển khai trên thiết bị hạn chế tài nguyên, tốc độ
  suy luận nhanh và dễ dàng tùy biến. Tuy nhiên, độ chính xác của
  MobileNet thường thấp hơn so với các mô hình phức tạp hơn như ResNet
  hay EfficientNet khi huấn luyện trên các tập dữ liệu lớn. Vì vậy,
  MobileNet phù hợp cho các ứng dụng thực tiễn yêu cầu cân bằng giữa độ
  chính xác và tốc độ hơn là các nghiên cứu cần hiệu năng cao nhất.

  \begin{enumerate}
  \def\labelenumi{\arabic{enumi}.}
  \item
    \protect\phantomsection\label{_Toc232112022}{}\emph{\textbf{Mô hình
    DenseNet}}
  \end{enumerate}
\end{itemize}

Trong số các kiến trúc mạng nơ-ron nhân tạo, DenseNet (Mạng tích chập
dày đặc) nổi bật như một giải pháp tiên tiến và hiệu quả cho nhiều bài
toán phức tạp, đặc biệt trong lĩnh vực nhận dạng và phân loại hình ảnh.
DenseNet được giới thiệu lần đầu tiên vào năm 2017 bởi Gao Huang và các
cộng sự trong bài báo ``\emph{Densely Connected Convolutional
Networks}'' {[}23{]}. Kiến trúc này đã mở ra một hướng tiếp cận mới
trong việc thiết kế mạng nơ-ron sâu thông qua việc sử dụng các kết nối
dày đặc.

\begin{itemize}
\item
  \textbf{Cấu trúc và nguyên lý hoạt động:}
\end{itemize}

Cấu trúc của DenseNet được xây dựng dựa trên một ý tưởng đơn giản nhưng
hiệu quả, tăng cường kết nối giữa các lớp trong mạng nơ-ron. Khác với
các kiến trúc mạng nơ-ron truyền thống, trong đó mỗi lớp chỉ nhận đầu
vào từ lớp liền trước và truyền đầu ra cho lớp liền sau, DenseNet cho
phép mỗi lớp nhận đầu vào từ tất cả các lớp trước đó và truyền đầu ra
đến tất cả các lớp tiếp theo. Cơ chế này tạo ra một mạng lưới dày đặc,
trong đó các đặc trưng được chia sẻ và tái sử dụng liên tục giữa các
lớp. Nhờ đó, DenseNet giúp cải thiện hiệu suất mô hình và giảm thiểu
hiện tượng mất mát thông tin trong quá trình lan truyền. Mạng DenseNet
được xây dựng từ các khối Dense (Dense Blocks) và các lớp chuyển tiếp
(Transition Layers).

Dense Block là một thành phần cốt lõi của kiến trúc DenseNet, bao gồm
nhiều lớp được kết nối dày đặc với nhau. Đặc điểm nổi bật của mỗi lớp
trong một Dense Block là khả năng nhận đầu vào từ tất cả các lớp trước
đó và truyền đầu ra đến tất cả các lớp tiếp theo. Mỗi lớp trong Dense
Block thường bao gồm ba thành phần chính là lớp chuẩn hóa theo batch
(Batch Normalization), lớp kích hoạt ReLU (Rectified Linear Unit), và
lớp tích chập (Convolutional Layer). Nhờ cơ chế tái sử dụng các đặc
trưng đã học từ các lớp trước, Dense Block giúp giảm độ phức tạp của mô
hình mà vẫn duy trì được độ chính xác cao.

\textbf{Transition Layer:} Transition Layer được đặt giữa các Dense
Block trong kiến trúc DenseNet, có nhiệm vụ giảm kích thước của bản đồ
đặc trưng (feature map) thông qua việc sử dụng lớp tích chập (1 × 1
Convolution) và phép gộp trung bình (Average Pooling). Lớp chuyển tiếp
đóng vai trò quan trọng trong việc kiểm soát kích thước và độ phức tạp
của mô hình, đồng thời ngăn chặn sự gia tăng quá mức của số lượng đặc
trưng khi mạng trở nên sâu hơn.

Mô hình DenseNet có nhiều ưu điểm như:

\begin{itemize}
\item
  \textbf{Tái sử dụng đặc trưng hiệu quả:} Trong các mạng nơ-ron truyền
  thống, mỗi lớp học các đặc trưng mới dựa trên đầu ra của lớp trước đó.
  Tuy nhiên, trong DenseNet, các đặc trưng đã học từ các lớp trước được
  tái sử dụng trực tiếp ở các lớp tiếp theo thông qua cơ chế kết nối dày
  đặc. Điều này không chỉ giúp mô hình học tập hiệu quả hơn mà còn giảm
  thiểu sự dư thừa không cần thiết, cho phép DenseNet đạt được hiệu suất
  cao với số lượng tham số ít hơn so với các kiến trúc khác.
\item
  \textbf{Giảm nguy cơ tiêu biến gradient (Gradient Vanishing):} Nhờ cơ
  chế kết nối dày đặc, gradient trong quá trình huấn luyện được truyền
  trực tiếp từ các lớp sau về các lớp trước một cách dễ dàng, giảm nguy
  cơ tiêu biến gradient, một vấn đề thường gặp trong các mạng nơ-ron
  sâu. Điều này đặc biệt hữu ích khi huấn luyện các mạng có số lượng lớp
  lớn, giúp DenseNet duy trì hiệu suất ổn định.
\item
  \textbf{Hiệu quả sử dụng bộ nhớ:} Mặc dù có nhiều kết nối hơn,
  DenseNet lại sử dụng bộ nhớ hiệu quả hơn nhờ số lượng tham số thực tế
  ít hơn so với các mô hình có cùng độ sâu, như VGGNet hoặc GoogLeNet.
  Điều này làm cho DenseNet trở thành lựa chọn lý tưởng trong các ứng
  dụng cần tối ưu hóa tài nguyên, chẳng hạn như trên các hệ thống nhúng
  hoặc thiết bị di động.
\item
  \textbf{Giảm nguy cơ quá khớp (Overfitting):} Cơ chế tái sử dụng đặc
  trưng và cấu trúc kết nối dày đặc của DenseNet giúp giảm nguy cơ quá
  khớp, ngay cả khi dữ liệu huấn luyện hạn chế. Đây là một lợi thế lớn
  khi làm việc với các tập dữ liệu nhỏ hoặc trong các môi trường có tài
  nguyên tính toán hạn chế.
\end{itemize}

Kể từ khi được giới thiệu, DenseNet đã được phát triển và mở rộng với
nhiều biến thể khác nhau, nhằm mục đích tối ưu hóa hiệu suất và thích
nghi với các bài toán khác nhau trong học sâu. Một số biến thể đáng chú
ý của DenseNet như DenseNet-121, DenseNet-169, DenseNet-201 và
DenseNet-264, được đặt tên dựa trên tổng số lớp (layers) trong mạng,
tương ứng với 121, 169, 201 và 264 lớp. Sự khác biệt chính giữa các biến
thể này nằm ở số lượng Dense Block và số lớp trong mỗi Dense Block, ảnh
hưởng đến độ sâu và khả năng học của mô hình.

\textbf{Các phiên bản của mô hình DenseNet:}

\begin{itemize}
\item
  \textbf{DenseNet-121:} Đây là phiên bản cơ bản nhất trong số các biến
  thể của DenseNet, thường được sử dụng rộng rãi trong nhiều thử nghiệm
  nhờ kích thước nhỏ gọn và hiệu suất cao. DenseNet-121 bao gồm 4 Dense
  Block với số lớp trong mỗi Dense block lần lượt là 6, 12, 24 và 16.
  Với cấu trúc này, DenseNet-121 đạt được sự cân bằng giữa hiệu suất và
  chi phí tính toán, phù hợp cho các bài toán phân loại hình ảnh trên
  các tập dữ liệu như CIFAR-10 hoặc ImageNet.
\item
  \textbf{DenseNet-169:} DenseNet-169 có số lớp lớn hơn, với 4 Dense
  Block chứa lần lượt 6, 12, 32 và 32 lớp. Sự gia tăng số lớp giúp mạng
  trở nên sâu hơn, cho phép học các đặc trưng phức tạp hơn so với
  DenseNet-121. Biến thể này thường được sử dụng trong các bài toán yêu
  cầu độ chính xác cao hơn, nhưng vẫn duy trì hiệu quả tính toán tương
  đối tốt.
\item
  \textbf{DenseNet-201:} Khác với các mạng nơ-ron tích chập (CNN) truyền
  thống, nơi mỗi lớp chỉ kết nối với lớp liền trước, DenseNet-201 tiếp
  tục khai thác cơ chế kết nối dày đặc ở độ sâu lớn hơn. Phiên bản này
  có khả năng học biểu diễn mạnh hơn DenseNet-121 nhưng cũng yêu cầu
  tài nguyên tính toán và bộ nhớ cao hơn, vì vậy không phải lựa chọn
  chính trong các thử nghiệm của luận văn.
\end{itemize}

\begin{quote}
\includegraphics[width=6.2in,height=2.97083in]{./media/image17.png}

\emph{\textbf{Hình 2.12.} Kiến trúc tổng quan và cơ chế kết nối dày đặc
của DenseNet} {[}5{]}
\end{quote}

\begin{itemize}
\item
  \textbf{DenseNet-264:} Đây là phiên bản sâu nhất trong các biến thể
  được đề cập, với 4 Dense Block chứa lần lượt 6, 12, 64 và 48 lớp.
  DenseNet-264 có khả năng xử lý các tập dữ liệu lớn và phức tạp, như
  ImageNet, nhờ vào độ sâu và khả năng học các đặc trưng ở nhiều cấp độ
  khác nhau. Tuy nhiên, do số lượng lớp lớn, DenseNet-264 đòi hỏi tài
  nguyên tính toán đáng kể, khiến nó phù hợp hơn cho các hệ thống có cấu
  hình mạnh.
\end{itemize}

Sự khác biệt giữa các biến thể này không chỉ nằm ở số lượng lớp mà còn ở
số lượng đặc trưng (feature maps) được tạo ra trong mỗi Dense Block. Các
phiên bản sâu hơn (như DenseNet-201 và DenseNet-264) thường có khả năng
học tốt hơn trên các tập dữ liệu phức tạp, nhưng đi kèm với chi phí tính
toán và bộ nhớ cao hơn. Trong khi đó, DenseNet-121 và DenseNet-169
thường được ưa chuộng trong các ứng dụng thực tế nhờ sự cân bằng giữa
hiệu suất và hiệu quả tài nguyên.

Trong luận văn này, biến thể được đưa vào thực nghiệm là
\textbf{DenseNet121}. Lựa chọn này phù hợp với quy mô bộ dữ liệu ảnh
tỏi vì DenseNet121 vẫn khai thác được ưu điểm tái sử dụng đặc trưng của
DenseNet, đồng thời có chi phí huấn luyện thấp hơn các biến thể sâu như
DenseNet-201 hoặc DenseNet-264.

\section{TÓM TẮT CHƯƠNG 2}\label{tuxf3m-tux1eaft-chux1b0ux1a1ng-2}

Chương 2 đã trình bày các nền tảng lý thuyết và cơ sở khoa học cần thiết
cho việc nghiên cứu, xây dựng và đánh giá mô hình phân loại chất lượng
tỏi. Trước hết, chương đã giới thiệu tổng quan về bài toán phân loại tỏi
sau bóc vỏ, trong đó nhiệm vụ chính là phân loại ảnh tép tỏi vào ba
nhóm: tỏi sạch vỏ, tỏi còn vỏ và tỏi hư. Đây là cơ sở quan trọng để xác
định rõ phạm vi bài toán, đặc điểm dữ liệu đầu vào và nhãn đầu ra của mô
hình.

Bên cạnh đó, bộ dữ liệu ảnh tỏi sử dụng trong nghiên cứu cũng được trình
bày chi tiết. Luận văn sử dụng Garlic Clove Dataset với hai phiên bản dữ
liệu là dataset-1 gồm 2.134 ảnh và dataset-2 gồm 2.944 ảnh. Việc sử dụng
hai phiên bản dữ liệu này giúp quá trình thực nghiệm có thêm cơ sở để
đánh giá hiệu quả mô hình trong các điều kiện phân bố lớp khác nhau, đặc
biệt trong bối cảnh dữ liệu thực tế thường có sự chênh lệch giữa các
nhóm mẫu.

Tiếp theo, chương đã trình bày các kiến thức nền tảng về học sâu và mạng
nơ-ron tích chập. Các thành phần chính trong một mô hình CNN như lớp
tích chập, hàm kích hoạt, lớp pooling và lớp phân loại được giới thiệu
nhằm làm rõ cách mô hình trích xuất và học đặc trưng từ ảnh. Đối với bài
toán phân loại chất lượng tỏi, CNN có vai trò quan trọng vì mô hình có
thể tự động học các đặc trưng thị giác trên bề mặt tép tỏi, từ những đặc
điểm tổng thể đến các dấu hiệu cục bộ như phần vỏ còn sót, vùng thâm,
dập nát hoặc hư hỏng.

Ngoài ra, chương cũng trình bày phương pháp học chuyển giao và kỹ thuật
fine-tuning. Đây là hướng tiếp cận phù hợp trong nghiên cứu này vì bộ dữ
liệu ảnh tỏi có quy mô không quá lớn so với các bộ dữ liệu ảnh phổ biến
trong học sâu. Việc tận dụng các mô hình đã được huấn luyện trước giúp
mô hình kế thừa những đặc trưng thị giác tổng quát, giảm thời gian huấn
luyện và cải thiện khả năng học trên dữ liệu ảnh tỏi.

Cuối cùng, chương đã khảo sát các mô hình CNN nền được sử dụng trong
thực nghiệm, bao gồm NASNetMobile, InceptionV3, EfficientNetB4,
MobileNetV2 và DenseNet121. Mỗi mô hình đại diện cho một hướng thiết kế
kiến trúc khác nhau, chẳng hạn NASNetMobile với kiến trúc được tìm kiếm
tự động, InceptionV3 với khả năng khai thác đặc trưng đa tỷ lệ,
EfficientNetB4 với phương pháp mở rộng mô hình cân bằng, MobileNetV2 với
thiết kế gọn nhẹ và DenseNet121 với cơ chế kết nối dày đặc giúp tái sử
dụng đặc trưng. Việc phân tích các mô hình này tạo cơ sở lý thuyết và
thực nghiệm cho quá trình so sánh, lựa chọn mô hình nền phù hợp, từ đó
phát triển mô hình đề xuất ở các chương tiếp theo.

\section{CHƯƠNG 3: MÔ HÌNH ĐỀ
XUẤT}\label{chux1b0ux1a1ng-3-muxf4-huxecnh-ux111ux1ec1-xuux1ea5t}

Trong chương này, luận văn tập trung trình bày và phân tích hai hướng tiếp
cận chính: Mô hình 1: Các mô hình học sâu tiền huấn luyện cơ sở
(Pre-trained CNN baseline), trong đó tiến hành thực nghiệm và so sánh 5
kiến trúc mạng khác nhau (NASNetMobile, InceptionV3, EfficientNetB4,
MobileNetV2, DenseNet121) để đánh giá hiệu năng trích xuất đặc trưng. Mô
hình 2: Mô hình đề xuất của luận văn (dựa trên kiến trúc đạt kết quả tốt
nhất là EfficientNetB4), kết hợp chiến lược tinh chỉnh (Fine-tuning)
chuyên sâu các khối tích chập (Blocks 2, 3, 4) và áp dụng các kỹ thuật
cải tiến về hàm mất mát, cơ chế học đa nhiệm nhằm tối ưu hóa độ chính
xác cho bài toán phân loại tỏi.

\textbf{Mô hình 1:} Nhóm mô hình học sâu tiền huấn luyện được re-train
trên bộ dữ liệu ảnh tỏi nhằm xác định kiến trúc nền phù hợp nhất cho bài
toán phân loại chất lượng tỏi.

\textbf{Mô hình 2:} Mô hình đề xuất được phát triển dựa trên kiến trúc
đạt kết quả tốt nhất từ Mô hình 1, sau đó bổ sung các cải tiến về chiến
lược fine-tuning, hàm mất mát và thiết kế đầu ra để nâng cao hiệu quả
phân loại.

\begin{quote}
\protect\phantomsection\label{_Toc232112025}{}\textbf{3.1 Các mô hình
học sâu tiền huấn luyện cho phân loại tỏi}

\protect\phantomsection\label{_Toc232112026}{}\emph{\textbf{3.1.1 Mục
tiêu}}
\end{quote}

Mục tiêu của giai đoạn này là re-train và đánh giá các mô hình CNN tiền
huấn luyện trên bộ dữ liệu ảnh tỏi sau bóc vỏ, từ đó lựa chọn một
backbone phù hợp để phát triển mô hình đề xuất ở Mục 3.2. Các mô hình
được đánh giá trên cùng bài toán phân loại ba lớp:
\textbf{Fully\_Peeled\_Garlic}, \textbf{Partially\_Peeled\_Garlic} và
\textbf{Spoiled\_Garlic}. Đây lần lượt là ba trạng thái chất lượng
tương ứng với tỏi đã bóc sạch, tỏi còn sót vỏ và tỏi hư hỏng.

Việc so sánh các mô hình tiền huấn luyện không chỉ nhằm tìm mô hình đạt
độ chính xác cao nhất, mà còn giúp đánh giá khả năng trích xuất đặc
trưng thị giác của từng kiến trúc đối với các dấu hiệu đặc thù của ảnh
tỏi, chẳng hạn vùng vỏ còn sót, vết thâm, vùng đổi màu và kết cấu bề mặt
tép tỏi. Kết quả của giai đoạn này là cơ sở để chọn kiến trúc nền cho
các cải tiến tiếp theo.

Quy trình thực nghiệm được thực hiện theo bốn bước chính: chuẩn bị dữ
liệu, xây dựng mô hình, huấn luyện nhiều lần với các seed khác nhau và
so sánh kết quả trên tập kiểm tra độc lập. Toàn bộ quá trình được tổ
chức thống nhất cho cả năm mô hình nhằm bảo đảm tính công bằng khi so
sánh.

\begin{quote}
\protect\phantomsection\label{_Toc232112027}{}\emph{\textbf{3.1.2 Các mô
hình được lựa chọn}}
\end{quote}

Nghiên cứu lựa chọn năm kiến trúc CNN tiền huấn luyện trên ImageNet gồm
\textbf{NASNetMobile}, \textbf{InceptionV3}, \textbf{EfficientNetB4},
\textbf{MobileNetV2} và \textbf{DenseNet121}. Các kiến trúc này đại diện
cho nhiều hướng thiết kế khác nhau trong học sâu cho ảnh:

\begin{itemize}
\item
  \textbf{NASNetMobile} là kiến trúc được thiết kế thông qua tìm kiếm
  kiến trúc mạng nơ-ron, có ưu điểm gọn nhẹ và phù hợp với các môi
  trường tính toán hạn chế.
\item
  \textbf{InceptionV3} khai thác đặc trưng đa tỷ lệ thông qua các mô-đun
  Inception, phù hợp với ảnh có nhiều vùng đặc trưng ở kích thước khác
  nhau.
\item
  \textbf{EfficientNetB4} sử dụng chiến lược compound scaling để cân
  bằng đồng thời độ sâu, độ rộng và độ phân giải đầu vào của mạng.
\item
  \textbf{MobileNetV2} sử dụng depthwise separable convolution và inverted
  residual block, phù hợp với các bài toán cần mô hình nhẹ.
\item
  \textbf{DenseNet121} khai thác cơ chế kết nối dày đặc giữa các lớp,
  giúp tái sử dụng đặc trưng và cải thiện lan truyền gradient.
\end{itemize}

\protect\phantomsection\label{_bookmark21}{}\emph{\textbf{3.1.3 Bộ dữ
liệu và tiền xử lý}}

Bộ dữ liệu sử dụng trong giai đoạn re-train là
\textbf{Garlic Clove Dataset -- dataset-1}, gồm \textbf{2.134 ảnh} tép
tỏi sau bóc vỏ. Dữ liệu được chia sẵn thành ba tập con: training,
validation và test. Cụ thể, tập training có \textbf{1.492 ảnh}, tập
validation có \textbf{318 ảnh} và tập test có \textbf{324 ảnh}. Ba lớp dữ
liệu được phân bố như sau: Fully\_Peeled\_Garlic gồm 690 ảnh,
Partially\_Peeled\_Garlic gồm 438 ảnh và Spoiled\_Garlic gồm 1.006 ảnh.

Việc chia dữ liệu được thực hiện trước khi huấn luyện để hạn chế rò rỉ dữ
liệu giữa các tập con. Mỗi mô hình chỉ học trên tập training, sử dụng tập
validation để theo dõi quá trình hội tụ và lựa chọn checkpoint tốt nhất,
sau đó đánh giá cuối cùng trên tập test độc lập.

Trước khi đưa vào mô hình, ảnh được resize theo kích thước đầu vào phù
hợp với từng kiến trúc. NASNetMobile và MobileNetV2 sử dụng ảnh
\(224 \times 224\), InceptionV3 sử dụng ảnh \(299 \times 299\), còn
EfficientNetB4 và DenseNet121 sử dụng ảnh \(380 \times 380\). Mỗi mô hình
sử dụng hàm tiền xử lý tương ứng trong thư viện Keras Applications để bảo
đảm ảnh đầu vào phù hợp với backbone đã được huấn luyện trước.

\protect\phantomsection\label{_Toc232112029}{}\emph{\textbf{3.1.4 Xây
dựng và huấn luyện mô hình}}

Các mô hình được xây dựng theo cùng một nguyên tắc transfer learning.
Phần backbone được khởi tạo bằng trọng số ImageNet và loại bỏ tầng phân
loại gốc (\texttt{include\_top=False}). Trên đầu backbone, nghiên cứu bổ
sung một classification head thống nhất gồm:

\begin{itemize}
\item
  Global Average Pooling để chuyển bản đồ đặc trưng thành vector đặc
  trưng;
\item
  Batch Normalization để ổn định phân bố đặc trưng trong quá trình học;
\item
  Dense layer 128 nút với hàm kích hoạt ReLU và L2 regularization
  \(1\times10^{-5}\);
\item
  Dropout 0.5 để giảm quá khớp;
\item
  Dense Softmax ba nút để dự đoán xác suất thuộc ba lớp chất lượng tỏi.
\end{itemize}

Trong giai đoạn so sánh baseline, các mô hình được huấn luyện theo chiến
lược \textbf{head-only fine-tuning}: backbone được đóng băng và chỉ các
tầng phân loại mới được cập nhật trọng số. Cách huấn luyện này giúp đánh
giá trực tiếp chất lượng đặc trưng đã học sẵn của từng kiến trúc khi áp
dụng vào bộ dữ liệu ảnh tỏi.

Để tăng khả năng khái quát hóa, tập training được áp dụng data
augmentation gồm xoay ảnh, dịch chuyển ngang/dọc, phóng to/thu nhỏ, lật
ngang, lật dọc và thay đổi độ sáng. Do dữ liệu có sự mất cân bằng giữa
các lớp, class weights được tính từ tần suất xuất hiện của từng lớp trong
tập training và đưa vào quá trình huấn luyện. Hàm mất mát sử dụng là
Categorical Cross-Entropy có label smoothing 0.15. Mỗi mô hình được chạy
ba lần với các seed \textbf{42}, \textbf{123} và \textbf{456}; kết quả
báo cáo là giá trị trung bình và độ lệch chuẩn trên ba lần chạy.

\emph{\textbf{Bảng 3.1.} Cấu hình re-train của năm mô hình tiền huấn
luyện}

{\def\LTcaptype{none} % do not increment counter
\begin{longtable}[]{@{}
  >{\raggedright\arraybackslash}p{(\linewidth - 8\tabcolsep) * \real{0.2200}}
  >{\centering\arraybackslash}p{(\linewidth - 8\tabcolsep) * \real{0.1600}}
  >{\centering\arraybackslash}p{(\linewidth - 8\tabcolsep) * \real{0.1300}}
  >{\centering\arraybackslash}p{(\linewidth - 8\tabcolsep) * \real{0.1700}}
  >{\raggedright\arraybackslash}p{(\linewidth - 8\tabcolsep) * \real{0.2600}}@{}}
\toprule\noalign{}
\begin{minipage}[b]{\linewidth}\raggedright
\textbf{Mô hình}
\end{minipage} & \begin{minipage}[b]{\linewidth}\centering
\textbf{Input}
\end{minipage} & \begin{minipage}[b]{\linewidth}\centering
\textbf{Batch}
\end{minipage} & \begin{minipage}[b]{\linewidth}\centering
\textbf{Learning rate}
\end{minipage} & \begin{minipage}[b]{\linewidth}\centering
\textbf{Thiết lập chung}
\end{minipage} \\
\midrule\noalign{}
\endhead
\bottomrule\noalign{}
\endlastfoot
NASNetMobile & \(224 \times 224\) & 32 & \(1\times10^{-5}\) &
Head-only, ImageNet, Adam + ExponentialDecay, tối đa 50 epoch,
EarlyStopping patience = 15 \\
InceptionV3 & \(299 \times 299\) & 32 & \(1\times10^{-5}\) &
Head-only, ImageNet, Adam + ExponentialDecay, tối đa 50 epoch,
EarlyStopping patience = 15 \\
EfficientNetB4 & \(380 \times 380\) & 64 & \(1\times10^{-4}\) &
Head-only, ImageNet, Adam + ExponentialDecay, tối đa 50 epoch,
EarlyStopping patience = 15 \\
MobileNetV2 & \(224 \times 224\) & 32 & \(1\times10^{-5}\) &
Head-only, ImageNet, Adam + ExponentialDecay, tối đa 50 epoch,
EarlyStopping patience = 15 \\
DenseNet121 & \(380 \times 380\) & 32 & \(1\times10^{-5}\) &
Head-only, ImageNet, Adam + ExponentialDecay, tối đa 50 epoch,
EarlyStopping patience = 15 \\
\end{longtable}
}

\emph{\textbf{Bảng 3.2.} Kết quả re-train năm mô hình tiền huấn luyện
trên dataset-1}

{\def\LTcaptype{none} % do not increment counter
\begin{longtable}[]{@{}
  >{\raggedright\arraybackslash}p{(\linewidth - 8\tabcolsep) * \real{0.2600}}
  >{\centering\arraybackslash}p{(\linewidth - 8\tabcolsep) * \real{0.1700}}
  >{\centering\arraybackslash}p{(\linewidth - 8\tabcolsep) * \real{0.1900}}
  >{\centering\arraybackslash}p{(\linewidth - 8\tabcolsep) * \real{0.1700}}
  >{\centering\arraybackslash}p{(\linewidth - 8\tabcolsep) * \real{0.1800}}@{}}
\toprule\noalign{}
\begin{minipage}[b]{\linewidth}\raggedright
\textbf{Mô hình}
\end{minipage} & \begin{minipage}[b]{\linewidth}\centering
\textbf{Accuracy}
\end{minipage} & \begin{minipage}[b]{\linewidth}\centering
\textbf{Precision}
\end{minipage} & \begin{minipage}[b]{\linewidth}\centering
\textbf{Recall}
\end{minipage} & \begin{minipage}[b]{\linewidth}\centering
\textbf{F1-Score}
\end{minipage} \\
\midrule\noalign{}
\endhead
\bottomrule\noalign{}
\endlastfoot
NASNetMobile & \(0.7706 \pm 0.0154\) & \(0.7897 \pm 0.0021\) &
\(0.7706 \pm 0.0154\) & \(0.7724 \pm 0.0134\) \\
InceptionV3 & \(0.7994 \pm 0.0157\) & \(0.8001 \pm 0.0157\) &
\(0.7994 \pm 0.0157\) & \(0.7988 \pm 0.0154\) \\
\textbf{EfficientNetB4} & \(\mathbf{0.8714} \pm \mathbf{0.0052}\) &
\(\mathbf{0.8885} \pm \mathbf{0.0049}\) &
\(\mathbf{0.8714} \pm \mathbf{0.0052}\) &
\(\mathbf{0.8722} \pm \mathbf{0.0052}\) \\
MobileNetV2 & \(0.7891 \pm 0.0038\) & \(0.8086 \pm 0.0080\) &
\(0.7891 \pm 0.0038\) & \(0.7910 \pm 0.0015\) \\
DenseNet121 & \(0.7901 \pm 0.0076\) & \(0.8015 \pm 0.0092\) &
\(0.7901 \pm 0.0076\) & \(0.7893 \pm 0.0073\) \\
\end{longtable}
}

Kết quả ở Bảng 3.2 cho thấy \textbf{EfficientNetB4} đạt hiệu năng cao
nhất trong cả bốn chỉ số đánh giá. Accuracy trung bình của EfficientNetB4
đạt \(0.8714\), cao hơn InceptionV3 khoảng \(0.0720\), MobileNetV2 khoảng
\(0.0823\), DenseNet121 khoảng \(0.0813\) và NASNetMobile khoảng
\(0.1008\). Đồng thời, độ lệch chuẩn của EfficientNetB4 tương đối nhỏ
(\(\pm 0.0052\) đối với Accuracy và F1-Score), cho thấy mô hình có độ ổn
định tốt qua ba lần chạy độc lập.

Từ kết quả này, EfficientNetB4 được lựa chọn làm backbone chính cho mô
hình đề xuất ở phần tiếp theo. Lựa chọn này phù hợp với đặc thù của bài
toán phân loại tỏi vì mô hình vừa khai thác tốt đặc trưng thị giác ở độ
phân giải cao, vừa duy trì hiệu năng ổn định khi dữ liệu có sự mất cân
bằng giữa các lớp.

\begin{quote}
\protect\phantomsection\label{_Toc232112030}{}\textbf{3.2 Mô hình đề
xuất}

\protect\phantomsection\label{_Toc232112031}{}\emph{\textbf{3.2.1 Hạn
chế của mô hình}}
\end{quote}

\protect\phantomsection\label{_Toc232112032}{}\emph{\textbf{3.2.2 Ý
tưởng cải tiến}}

\begin{quote}
\emph{\textbf{Bảng 3.3.} Bảng Hyper-parameters~của mô hình đề xuất}
\end{quote}

{\def\LTcaptype{none} % do not increment counter
\begin{longtable}[]{@{}
  >{\raggedright\arraybackslash}p{(\linewidth - 4\tabcolsep) * \real{0.2456}}
  >{\centering\arraybackslash}p{(\linewidth - 4\tabcolsep) * \real{0.3421}}
  >{\centering\arraybackslash}p{(\linewidth - 4\tabcolsep) * \real{0.3333}}@{}}
\toprule\noalign{}
\begin{minipage}[b]{\linewidth}\raggedright
\end{minipage} & \begin{minipage}[b]{\linewidth}\centering
\end{minipage} & \begin{minipage}[b]{\linewidth}\centering
\end{minipage} \\
\midrule\noalign{}
\endhead
\bottomrule\noalign{}
\endlastfoot
& & \\
& & \\
& & \\
& & \\
& & \\
& & \\
\end{longtable}
}

Bảng 3.3

\protect\phantomsection\label{_Toc232112033}{}\emph{\textbf{3.2.3 Kiến
trúc}}

\section{CHƯƠNG 4: KẾT QUẢ THỰC
NGHIỆM}\label{chux1b0ux1a1ng-4-kux1ebft-quux1ea3-thux1ef1c-nghiux1ec7m}

Trong Chương 3, luận văn đã trình bày các mô hình được sử dụng, bao gồm
nhóm mô hình học sâu tiền huấn luyện (Pre-trained CNN) và mô hình cải
tiến

\textbf{Trọng tâm của chương là:}

Kết quả trong chương này sẽ làm cơ sở cho phần Kết luận ở Chương 5.

\begin{quote}
\protect\phantomsection\label{_Toc232112035}{}\textbf{4.1 Tiêu chí đánh
giá và công cụ hỗ trợ}
\end{quote}

\protect\phantomsection\label{_Toc232112036}{}\emph{\textbf{4.1.1 Tiêu
chí đánh giá}}

Để đánh giá khách quan hiệu năng của các mô hình, nghiên cứu sử dụng bốn
tiêu chí phân loại chuẩn, được tính toán trên cả tập validation và test
độc lập, bao gồm:

\begin{itemize}
\item
  \textbf{Độ chính xác (Accuracy):}
\end{itemize}

{\def\LTcaptype{none} % do not increment counter
\begin{longtable}[]{@{}
  >{\raggedright\arraybackslash}p{(\linewidth - 6\tabcolsep) * \real{0.0548}}
  >{\raggedright\arraybackslash}p{(\linewidth - 6\tabcolsep) * \real{0.8159}}
  >{\raggedright\arraybackslash}p{(\linewidth - 6\tabcolsep) * \real{0.0507}}
  >{\raggedright\arraybackslash}p{(\linewidth - 6\tabcolsep) * \real{0.0484}}@{}}
\toprule\noalign{}
\begin{minipage}[b]{\linewidth}\raggedright
\end{minipage} & \begin{minipage}[b]{\linewidth}\raggedright
\begin{quote}
\[Accuracy = \frac{TP + TN}{TP + TN + FP + FN}\]
\end{quote}
\end{minipage} & \begin{minipage}[b]{\linewidth}\raggedright
(7)
\end{minipage} & \begin{minipage}[b]{\linewidth}\raggedright
\end{minipage} \\
\midrule\noalign{}
\endhead
\bottomrule\noalign{}
\endlastfoot
\end{longtable}
}

Chỉ số phản ánh tỷ lệ dự đoán đúng trên tổng số mẫu, thể hiện mức độ
chính xác tổng thể của mô hình.

\begin{itemize}
\item
  \textbf{Độ chính xác dự đoán (Precision):}
\end{itemize}

{\def\LTcaptype{none} % do not increment counter
\begin{longtable}[]{@{}
  >{\raggedright\arraybackslash}p{(\linewidth - 6\tabcolsep) * \real{0.0548}}
  >{\raggedright\arraybackslash}p{(\linewidth - 6\tabcolsep) * \real{0.8159}}
  >{\raggedright\arraybackslash}p{(\linewidth - 6\tabcolsep) * \real{0.0507}}
  >{\raggedright\arraybackslash}p{(\linewidth - 6\tabcolsep) * \real{0.0484}}@{}}
\toprule\noalign{}
\begin{minipage}[b]{\linewidth}\raggedright
\end{minipage} & \begin{minipage}[b]{\linewidth}\centering
\[Precision = \frac{TP}{TP + FP}\]
\end{minipage} & \begin{minipage}[b]{\linewidth}\raggedright
(8)
\end{minipage} & \begin{minipage}[b]{\linewidth}\raggedright
\end{minipage} \\
\midrule\noalign{}
\endhead
\bottomrule\noalign{}
\endlastfoot
\end{longtable}
}

Phản ánh mức độ chính xác trong các dự đoán dương tính của mô hình.
Precision cao đặc biệt quan trọng trong kiểm tra chất lượng tỏi, vì giúp
hạn chế các cảnh báo sai khi xác định mẫu tỏi không đạt yêu cầu.

\begin{itemize}
\item
  \textbf{Độ bao phủ (Recall):}
\end{itemize}

{\def\LTcaptype{none} % do not increment counter
\begin{longtable}[]{@{}
  >{\raggedright\arraybackslash}p{(\linewidth - 6\tabcolsep) * \real{0.0548}}
  >{\raggedright\arraybackslash}p{(\linewidth - 6\tabcolsep) * \real{0.8159}}
  >{\raggedright\arraybackslash}p{(\linewidth - 6\tabcolsep) * \real{0.0507}}
  >{\raggedright\arraybackslash}p{(\linewidth - 6\tabcolsep) * \real{0.0484}}@{}}
\toprule\noalign{}
\begin{minipage}[b]{\linewidth}\raggedright
\end{minipage} & \begin{minipage}[b]{\linewidth}\raggedright
\begin{quote}
\[Recall = \frac{TP}{TP + FN}\]
\end{quote}
\end{minipage} & \begin{minipage}[b]{\linewidth}\raggedright
(9)
\end{minipage} & \begin{minipage}[b]{\linewidth}\raggedright
\end{minipage} \\
\midrule\noalign{}
\endhead
\bottomrule\noalign{}
\endlastfoot
\end{longtable}
}

Chỉ số đo lường khả năng phát hiện đầy đủ các mẫu thuộc từng lớp thực tế,
đặc biệt là các mẫu tỏi còn vỏ hoặc tỏi hư cần được tách khỏi nhóm đạt
yêu cầu. Đây là tiêu chí quan trọng để hạn chế bỏ sót mẫu lỗi trong dây
chuyền chế biến.

\begin{itemize}
\item
  \textbf{F1-Score:}
\end{itemize}

{\def\LTcaptype{none} % do not increment counter
\begin{longtable}[]{@{}
  >{\raggedright\arraybackslash}p{(\linewidth - 6\tabcolsep) * \real{0.0548}}
  >{\raggedright\arraybackslash}p{(\linewidth - 6\tabcolsep) * \real{0.8159}}
  >{\raggedright\arraybackslash}p{(\linewidth - 6\tabcolsep) * \real{0.0634}}
  >{\raggedright\arraybackslash}p{(\linewidth - 6\tabcolsep) * \real{0.0484}}@{}}
\toprule\noalign{}
\begin{minipage}[b]{\linewidth}\raggedright
\end{minipage} & \begin{minipage}[b]{\linewidth}\centering
\[F1 - Score = 2.\frac{Precision.Recall}{Precision + Recall}\]
\end{minipage} & \begin{minipage}[b]{\linewidth}\raggedright
(10)
\end{minipage} & \begin{minipage}[b]{\linewidth}\raggedright
\end{minipage} \\
\midrule\noalign{}
\endhead
\bottomrule\noalign{}
\endlastfoot
\end{longtable}
}

Là chỉ số trung bình giữa Precision và Recall, giúp cân bằng khi dữ liệu
bị mất cân bằng giữa các lớp.

\textbf{Trong đó:}

\begin{itemize}
\item
  \textbf{TP (True Positive):} số trường hợp dương tính được dự đoán
  đúng.
\item
  \textbf{TN (True Negative):} số trường hợp âm tính được dự đoán đúng.
\item
  \textbf{FP (False Positive):} số trường hợp âm tính nhưng bị dự đoán
  sai thành dương tính.
\item
  \textbf{FN (False Negative):} số trường hợp dương tính nhưng bị dự
  đoán sai thành âm tính.
\end{itemize}

Ngoài ra, nghiên cứu còn sử dụng:

\begin{quote}
\textbf{Ma trận nhầm lẫn (Confusion Matrix):} dùng để quan sát chi tiết
các lớp chất lượng tỏi dễ bị nhầm lẫn.
\end{quote}

Những tiêu chí này không chỉ phản ánh khả năng phân loại tổng thể mà còn
làm rõ mức độ tin cậy và khả năng khái quát hóa của mô hình khi áp dụng
trong thực tế sản xuất nông nghiệp.

\protect\phantomsection\label{_Toc232112037}{}\emph{\textbf{4.1.2 Công
cụ hỗ trợ}}

Trong quá trình xây dựng và huấn luyện mô hình, nghiên cứu sử dụng các
công cụ chính sau:

\begin{itemize}
\item
  \textbf{TensorFlow và Keras:} Framework học sâu chủ đạo, cung cấp sẵn
  các mô hình Pre-trained (NASNetMobile, DenseNet, EfficientNet,
  Inception, MobileNet).
\item
  \textbf{Scikit-learn:} Hỗ trợ tính toán các chỉ số đánh giá
  (Precision, Recall, F1, Confusion Matrix, Top-K Accuracy).
\item
  \textbf{Matplotlib và Seaborn:} Dùng để trực quan hóa quá trình huấn
  luyện (Accuracy, Loss) và ma trận nhầm lẫn.
\item
  \textbf{Google Colab Pro}: Môi trường huấn luyện chính, tận dụng
  GPU/TPU để tăng tốc quá trình tính toán.
\item
  \textbf{Google Drive:} Lưu trữ dữ liệu, checkpoint mô hình, lịch sử
  huấn luyện và kết quả thực nghiệm.
\end{itemize}

Việc kết hợp các tiêu chí đánh giá khoa học cùng công cụ hỗ trợ mạnh mẽ
đảm bảo rằng quá trình thực nghiệm khách quan, minh bạch, có thể tái lập
và ứng dụng thực tiễn.

\begin{quote}
\protect\phantomsection\label{_Toc232112038}{}\textbf{4.2 Kết quả của 5
mô hình Pre-trained}
\end{quote}

Các mô hình được tiến hành huấn luyện và đánh giá trên tập kiểm định,
với việc áp dụng các kỹ thuật như EarlyStopping, ReduceLROnPlateau và
ModelCheckpoint nhằm tối ưu quá trình huấn luyện và ngăn ngừa hiện tượng
quá khớp. Sau khi quá trình huấn luyện hoàn tất, mô hình được kiểm tra
trên tập dữ liệu độc lập để đánh giá khả năng tổng quát hóa. Các chỉ số
được báo cáo bao gồm: Accuracy, Precision, Recall, F1-score và ma trận
nhầm lẫn, giúp phân tích chi tiết hiệu quả phân loại ở từng lớp chất
lượng tỏi.

\textbf{Kết quả theo từng mô hình:}

\emph{\textbf{Bảng 4.1.}} \emph{Bảng kết quả của năm mô hình
Pre-trained}

{\def\LTcaptype{none} % do not increment counter
\begin{longtable}[]{@{}
  >{\centering\arraybackslash}p{(\linewidth - 10\tabcolsep) * \real{0.1222}}
  >{\centering\arraybackslash}p{(\linewidth - 10\tabcolsep) * \real{0.3232}}
  >{\centering\arraybackslash}p{(\linewidth - 10\tabcolsep) * \real{0.1185}}
  >{\centering\arraybackslash}p{(\linewidth - 10\tabcolsep) * \real{0.1271}}
  >{\centering\arraybackslash}p{(\linewidth - 10\tabcolsep) * \real{0.1008}}
  >{\centering\arraybackslash}p{(\linewidth - 10\tabcolsep) * \real{0.1232}}@{}}
\toprule\noalign{}
\begin{minipage}[b]{\linewidth}\centering
No
\end{minipage} & \begin{minipage}[b]{\linewidth}\centering
Models
\end{minipage} & \begin{minipage}[b]{\linewidth}\centering
Accuracy
\end{minipage} & \begin{minipage}[b]{\linewidth}\centering
F1-Score
\end{minipage} & \begin{minipage}[b]{\linewidth}\centering
Recall
\end{minipage} & \begin{minipage}[b]{\linewidth}\centering
Precision
\end{minipage} \\
\midrule\noalign{}
\endhead
\bottomrule\noalign{}
\endlastfoot
\textbf{1} & & & & & \\
2 & & & & & \\
3 & & & & & \\
4 & & & & & \\
5 & & & & & \\
\end{longtable}
}

Bảng 4.1 So sánh hiệu năng của năm mô hình học sâu

\begin{quote}
\protect\phantomsection\label{_Toc232112039}{}\textbf{4.3 Kết quả của mô
hình đề xuất}
\end{quote}

\emph{\textbf{Bảng 4.2.} So sánh các mô hình CNN trong phân loại chất
lượng tỏi}

{\def\LTcaptype{none} % do not increment counter
\begin{longtable}[]{@{}
  >{\centering\arraybackslash}p{(\linewidth - 10\tabcolsep) * \real{0.0699}}
  >{\centering\arraybackslash}p{(\linewidth - 10\tabcolsep) * \real{0.2815}}
  >{\centering\arraybackslash}p{(\linewidth - 10\tabcolsep) * \real{0.1671}}
  >{\centering\arraybackslash}p{(\linewidth - 10\tabcolsep) * \real{0.1847}}
  >{\centering\arraybackslash}p{(\linewidth - 10\tabcolsep) * \real{0.1056}}
  >{\centering\arraybackslash}p{(\linewidth - 10\tabcolsep) * \real{0.1144}}@{}}
\toprule\noalign{}
\begin{minipage}[b]{\linewidth}\centering
\textbf{No}
\end{minipage} & \begin{minipage}[b]{\linewidth}\centering
\textbf{Models}
\end{minipage} & \begin{minipage}[b]{\linewidth}\centering
\end{minipage} & \begin{minipage}[b]{\linewidth}\centering
\end{minipage} & \begin{minipage}[b]{\linewidth}\centering
\end{minipage} & \begin{minipage}[b]{\linewidth}\centering
\end{minipage} \\
\midrule\noalign{}
\endhead
\bottomrule\noalign{}
\endlastfoot
1 & & & & & \\
2 & & & & & \\
3 & & & & & \\
4 & & & & & \\
5 & & & & & \\
6 & & & & & \\
7 & & & & & \\
8 & & & & & \\
\end{longtable}
}

\emph{Hình 4.13. Ma trận nhầm lẫn (Confusion Matrix)}

\section{CHƯƠNG 5: KẾT LUẬN}\label{chux1b0ux1a1ng-5-kux1ebft-luux1eadn}

\begin{quote}
\protect\phantomsection\label{_Toc232112041}{}\textbf{5.1 Tổng kết quá
trình nghiên cứu}

\protect\phantomsection\label{_Toc232112042}{}\textbf{5.2 Hạn chế và
hướng phát triển}
\end{quote}

\protect\phantomsection\label{_Toc232112043}{}\emph{\textbf{5.2.1 Hạn
chế:}}

\protect\phantomsection\label{_Toc232112044}{}\emph{\textbf{5.2.2 Hướng
phát triển trong tương lai:}}

\begin{quote}
\protect\phantomsection\label{_Toc232112045}{}\textbf{5.3 Khả năng ứng
dụng thực tiễn}
\end{quote}

\begin{itemize}
\tightlist
\item
\end{itemize}

\section{TÀI LIỆU THAM KHẢO}\label{tuxe0i-liux1ec7u-tham-khux1ea3o}

1. B. Zoph, V. Vasudevan, J. Shlens, and Q. V. Le, ``Learning
transferable architectures for scalable image recognition,'' in Proc.
IEEE Conf. Computer Vision and Pattern Recognition (CVPR), 2018, pp.
8697--8710.

2. Thủ tướng Chính phủ, ``\emph{Quyết định số 127/QĐ-TTg về việc phê
duyệt Chiến lược quốc gia về nghiên cứu, phát triển và ứng dụng trí tuệ
nhân tạo đến năm 2030},'' Hà Nội, Việt Nam, 2021. {[}Online{]}.
Available: \url{https://vanban.chinhphu.vn/default.aspx?docid=202565&pageid=27160}.

3. Rashid, Mohammed, et al. ``\emph{Plant Disease Detection Using Deep
Neural Networks and Computer Vision.}'', 2023.

4. Trịnh, C. Đ. (2023). Ứng dụng học sâu để nhận dạng và phân loại bệnh
trên cây lúa: Nghiên cứu sử dụng kiến trúc mạng dư (ResNet). Luận văn
thạc sĩ, Đại học Nông nghiệp Việt Nam, Hà Nội.

5. G. Huang, Z. Liu, L. van der Maaten, and K. Q. Weinberger, ``Densely
Connected Convolutional Networks,'' in Proc. IEEE Conf. Computer Vision
and Pattern Recognition (CVPR), Honolulu, HI, 2017, pp. 2261--2269.

6. M. Tan and Q. V. Le, ``\emph{EfficientNetV2: Smaller Models and
Faster Training,}'' in Proc. 38th Int. Conf. Machine Learning, vol. 139,
pp. 10096--10106, 2021.

7. C. Szegedy, V. Vanhoucke, S. Ioffe, J. Shlens, and Z. Wojna,
``Rethinking the Inception Architecture for Computer Vision,'' in Proc.
IEEE Conf. Computer Vision and Pattern Recognition (CVPR), Las Vegas,
NV, 2016, pp. 2818--2826.

8. M. Sandler, A. Howard, M. Zhu, A. Zhmoginov, and L. C. Chen,
``MobileNetV2: Inverted Residuals and Linear Bottlenecks,'' in Proc.
IEEE/CVF Conf. Computer Vision and Pattern Recognition (CVPR), Salt Lake
City, UT, 2018, pp. 4510--4520.

9. F. Chollet, ``Xception: Deep learning with depthwise separable
convolutions,'' in Proc. IEEE Conf. Computer Vision and Pattern
Recognition (CVPR), Honolulu, HI, 2017, pp. 1251--1258.

10. D. Chicco, and G. Jurman, ``The advantages of the Matthews
correlation coefficient (MCC) over F1 score and Accuracy in binary
classification evaluation,'' IEEE Access, vol. 7, pp. 749--762, 2019.

11. M. Tan and Q. V. Le, ``EfficientNet: Rethinking model scaling for
convolutional neural networks,'' in Proc. 36th Int. Conf. Machine
Learning (ICML), 2019, pp. 6105--6114.

12. M. Meenu, C. Kurade, B. C. Neelapu, S. Kalra, H. S. Ramaswamy, and
Y. Yu, ``A concise review on food quality assessment using digital image
processing,'' Trends in Food Science \& Technology, vol. 118, pp.
106--124, 2021.

13. K. Wonggasem, P. Chakranon, and P. Wongchaisuwat, ``Automated quality
inspection of baby corn using image processing and deep learning,''
Artificial Intelligence in Agriculture, vol. 11, pp. 61--69, 2024.

14. S. J. Pan and Q. Yang, ``A Survey on Transfer Learning,'' IEEE
Transactions on Knowledge and Data Engineering, vol. 22, no. 10, pp.
1345--1359, Oct. 2010.

15. J. Yosinski, J. Clune, Y. Bengio, and H. Lipson, ``How transferable
are features in deep neural networks?'' in Proc. Advances in Neural
Information Processing Systems (NeurIPS), vol. 27, 2014.

16. S. Thrun and L. Pratt, Eds., Learning to Learn. New York, NY, USA:
Springer, 1998, pp. 1--392.

17. Dosovitskiy, A., Beyer, L., Kolesnikov, A., Weissenborn, D., Zhai,
X., Unterthiner, T., Dehghani, M., Minderer, M., Heigold, G., Gelly, S.,
Uszkoreit, J., \& Houlsby, N. (2021). An Image is Worth 16x16 Words:
Transformers for Image Recognition at Scale.

18. Szegedy, C., Vanhoucke, V., Ioffe, S., \& Shlens, J. (2016).
Rethinking the Inception Architecture for Computer Vision. In Proc. IEEE
Conf. Computer Vision and Pattern Recognition (CVPR), Las Vegas, NV, pp.
2818--2826.

19. Chollet, F. (2017). Xception: Deep Learning with Depthwise Separable
Convolutions. In Proceedings of the IEEE Conference on Computer Vision
and Pattern Recognition (CVPR).

20. Sumalatha, K., Hisaria, S., Sharma, P., \& Gupta, R. (2023). An
Analysis of Multi-Criteria Performance in Deep Learning-Based Medical
Image Classification: A Comprehensive Review. Journal Name, vol. X, pp.
Y--Z.

21. C. Szegedy, W. Liu, Y. Jia, P. Sermanet, S. Reed, D. Anguelov, D.
Erhan, V. Vanhoucke, and A. Rabinovich, ``Going deeper with
convolutions,'' in Proc. IEEE Conf. Comput. Vis. Pattern Recognit.
(CVPR), 2015, pp. 1--9.

22. O. Russakovsky et al., ``ImageNet large scale visual recognition
challenge,'' Int. J. Comput. Vis., vol. 115, no. 3, pp. 211--252, Dec.
2015.

23. Huang, G., Liu, Z., Van Der Maaten, L., \& Weinberger, K. Q. (2017).
Densely connected convolutional networks. Proceedings of the IEEE
Conference on Computer Vision and Pattern Recognition (CVPR), 4700-4708.

24. A. G. Howard et al., ``MobileNets: Efficient convolutional neural
networks for mobile vision applications,'' arXiv preprint
arXiv:1704.04861, 2017.

25. P. Hoang and T. Le, ``Deep Learning-Based Classification of Peeled
Garlic Cloves for Smart Food Processing,'' manuscript, HUTECH
University, 2026.

26. Cục Thống kê, ``Báo cáo tình hình kinh tế -- xã hội quý IV và năm
2025,'' 2026. {[}Online{]}. Available:
\url{https://www.nso.gov.vn/du-lieu-va-so-lieu-thong-ke/2026/01/bao-cao-tinh-hinh-kinh-te-xa-hoi-quy-iv-va-nam-2025/}.

27. K. G. Liakos, V. Athanasiadis, E. Bozinou, and S. I. Lalas,
``Machine Learning for Quality Control in the Food Industry: A Review,''
Foods, vol. 14, no. 19, article 3424, 2025. {[}Online{]}. Available:
\url{https://www.mdpi.com/2304-8158/14/19/3424}.

28. P. T. Q. Anh, D. Q. Thuyet, and Y. Kobayashi, ``Image classification
of root-trimmed garlic using multi-label and multi-class classification
with deep convolutional neural network,'' Postharvest Biology and
Technology, vol. 190, article 111956, 2022. {[}Online{]}. Available:
\url{https://doi.org/10.1016/j.postharvbio.2022.111956}.

29. PlantwisePlus Blog, ``AI Plant Doctor app launched to help diagnose
pests and diseases on dragon fruit crops in Vietnam,'' 2021.
{[}Online{]}. Available:
\url{https://blog.plantwise.org/2021/02/18/ai-plant-doctor-app-launched-to-help-diagnose-pests-and-diseases-on-dragon-fruit-crops-in-vietnam/}.

30. CESTI, ``Ứng dụng AI trong quản lý bệnh hại cây trồng,'' 2025.
{[}Online{]}. Available:
\url{https://thongke.cesti.gov.vn/an-pham-thong-ke/thong-tin-chuyen-de-khoa-hoc-cong-nghe-va-doi-moi-sang-tao/1431-ung-dung-ai-trong-quan-ly-benh-hai-cay-trong}.

31. Học viện Nông nghiệp Việt Nam, ``Nghiên cứu thuật toán học sâu mạng
nơ-ron tích chập và ứng dụng trong nhận dạng hình ảnh ong mật mang
phấn,'' 2023. {[}Online{]}. Available:
\url{https://vnua.edu.vn/tin-tuc-su-kien/nghien-cuu-khoa-hoc/nghien-cuu-thuat-toan-hoc-sau-mang-noron-tich-chap-va-ung-dung-trong-nhan-dang-hinh-anh-ong-mat-mang-phan-55308}.

32. Khoa Công nghệ thực phẩm, Học viện Nông nghiệp Việt Nam, ``Ứng dụng
trí tuệ nhân tạo (AI) vào phát hiện độc tố vi nấm trong thực phẩm,''
2025. {[}Online{]}. Available:
\url{https://cntp.vnua.edu.vn/tin-tuc-va-su-kien/ung-dung-tri-tue-nhan-tao-ai-vao-phat-hien-doc-to-vi-nam-trong-thuc-pham-64051}.

33. IBM, ``What Is a Neural Network?'' {[}Online{]}. Available:
\url{https://www.ibm.com/think/topics/neural-networks}.

34. IBM, ``What are Convolutional Neural Networks?'' {[}Online{]}.
Available: \url{https://www.ibm.com/think/topics/convolutional-neural-networks}.

35. Stanford University, ``CS231n: Convolutional Neural Networks for
Visual Recognition.'' {[}Online{]}. Available:
\url{https://cs231n.github.io/convolutional-networks/}.

36. A. Zhang, Z. C. Lipton, M. Li, and A. J. Smola, ``Convolutional
Neural Networks,'' in \emph{Dive into Deep Learning}. {[}Online{]}.
Available: \url{https://d2l.ai/chapter_convolutional-neural-networks/}.

37. TensorFlow, ``Transfer learning and fine-tuning.'' {[}Online{]}.
Available: \url{https://www.tensorflow.org/tutorials/images/transfer_learning}.

38. I. Rojas Santelices, S. Cano, F. Moreira, and A. Peña Fritz,
``Artificial Vision Systems for Fruit Inspection and Classification:
Systematic Literature Review,'' Sensors, vol. 25, no. 5, article 1524,
2025. {[}Online{]}. Available: \url{https://doi.org/10.3390/s25051524}.

39. F. Liu, Y. Zhang, C. Du, X. Ren, B. Huang, and X. Chai, ``Design and
Experimentation of a Machine Vision-Based Cucumber Quality Grader,''
Foods, vol. 13, no. 4, article 606, 2024. {[}Online{]}. Available:
\url{https://doi.org/10.3390/foods13040606}.

40. Y. Gao, X. Ma, Z. Xia, T. Qi, X. Wang, Z. He, and G. Chen,
``Garlic-YOLO-DD: a lightweight object detection algorithm for garlic
damage detection,'' Frontiers in Plant Science, vol. 16, article 1702045,
2026. {[}Online{]}. Available:
\url{https://doi.org/10.3389/fpls.2025.1702045}.

41. Z. He, X. Ma, Y. Jiang, G. Chen, Y. Gao, X. Wang, T. Qi, and Z. Xia,
``DH-GarlicNet: a precise identification method for garlic damage based
on the improved residual network,'' Frontiers in Plant Science, vol. 17,
article 1801346, 2026. {[}Online{]}. Available:
\url{https://doi.org/10.3389/fpls.2026.1801346}.

\end{document}
